%%%%%%%%%%%%%%%%%%%%%%%%%%%%%%%%%%%%%%%%%%%%%%%%%%%%%%%%%%%%%%%%%%%%%%
% Lớp 11 | Hàm số lượng giác và phương trình lượng giác | Giá trị lượng giác của một góc lượng giác
% Dạng 1D1?2-2: Xét dấu giá trị lượng giác. Tính giá trị lượng giác của một góc
% Số câu: 149  (sinh tự động bởi tach_cau.py; từ nay thầy sửa trực tiếp file này)
%%%%%%%%%%%%%%%%%%%%%%%%%%%%%%%%%%%%%%%%%%%%%%%%%%%%%%%%%%%%%%%%%%%%%%

%%%%%%%%%%%%%%%% TRẮC NGHIỆM 4 PHƯƠNG ÁN %%%%%%%%%%%%%%%%

% nguồn: Ok_Toan11_TracNghiem_Lan2
\begin{ex}%[1D1N2-2]%[ma:e1ea04f9]
	Cho $\dfrac{\pi}{2} <\alpha <\pi $. Kết quả đúng là 
	 \choice
		{$\sin \alpha >0,$  $\cos \alpha >0$}
		{$\sin \alpha <0,$ $\cos \alpha  <0$}
		{\True $\sin  \alpha >0,$ $\cos \alpha <0$}
		{$\sin \alpha  <0,$ $\cos \alpha >0$}
	\loigiai{ 
	Vì $\dfrac{\pi}{2} <\alpha <\pi $ nên $\alpha $ nằm trong  góc phần tư thứ $II$ nên $\sin \alpha >0$, $\cos \alpha <0$.}
\end{ex}

% nguồn: Ok_Toan11_TracNghiem_Lan2
\begin{ex}%[1D1N2-2]%[ma:0bae8baf]
	Cho $0<\alpha <\dfrac{\pi}{2}$. Khẳng định nào  sau đây đúng?
	\choice 
		{\True $\tan \alpha  >0,\cot \alpha >0$}
		{$\tan \alpha >0,\cot \alpha <0$} 
		{$\tan \alpha <0,\cot \alpha <0$} 
		{$\tan \alpha <0,\cot  \alpha >0$}
	\loigiai{ 
		 Điểm biểu diễn cung lượng giác $\alpha $ nằm góc phần tư thứ nhất nên $\tan \alpha >0$, $\cot \alpha >0$.
	} 
\end{ex}

% nguồn: Ok_Toan11_TracNghiem_Lan2
\begin{ex}%[1D1N2-2]%[ma:a7fa6522]
	 Cho $2\pi <\alpha <\dfrac{5\pi}{2}$. Khẳng định nào sau đây đúng?
	\choice 
		 {\True $\tan \alpha >0$; $\cot \alpha >0$}
		{$\tan \alpha <0$; $\cot  \alpha <0$}
		 {$\tan \alpha >0$; $\cot \alpha <0$}
		{$\tan  \alpha <0$; $\cot \alpha >0$}
	\loigiai{ 
		Ta có $2\pi <\alpha <\dfrac{5\pi}{2}$, suy ra điểm cuối cung $\alpha -\pi $ thuộc góc phần  tư thứ $I$. \\
	 Do đó $\tan \alpha >0$ và $\cot \alpha >0$.}
\end{ex}

% nguồn: Ok_Toan11_TracNghiem_Lan2
\begin{ex}%[1D1N2-2]%[ma:b5c8fedf]
	 Cho góc $\alpha $ thỏa mãn $2\pi <\alpha <\dfrac{5\pi}{2} $. Khẳng định nào sau đây \textbf{sai}?
	\choice 
		{\True $\tan  \alpha <0$}
		{$\cot  \alpha >0$}
		{$\sin  \alpha >0$}
		{$\cos \alpha >0$} 
	 \loigiai{
	Với $2\pi <\alpha  <\dfrac{5\pi}{2} $ ta có $\sin \alpha >0$, $\cos \alpha >0$, $\tan \alpha >0$, $\cot \alpha >0$.}
\end{ex}

% nguồn: Ok_Toan11_TracNghiem_Lan2
\begin{ex}%[1D1N2-2]%[ma:2c19a23c]
	Cho $0< \alpha < \dfrac{\pi}{2}$.  Tính $\sqrt{\dfrac{1+\sin \alpha}{1-\sin \alpha}}+\sqrt{\dfrac{1-\sin \alpha}{1+\sin \alpha}}$.
	 \choice
		{$-\dfrac{2}{\cos  \alpha}$}
		{\True $\dfrac{2}{\cos \alpha}$} 
		 {$\dfrac{2}{\sin \alpha}$}
		{$-\dfrac{2}{\sin \alpha}$} 
	 \loigiai{
		Đặt $A=\sqrt{\dfrac{1+\sin \alpha}{1-\sin \alpha}}+\sqrt{\dfrac{1-\sin  \alpha}{1+\sin \alpha}}$.\\
		Khi đó 
		\[A^2=\left(\sqrt{\dfrac{1+\sin \alpha}{1-\sin  \alpha}}+\sqrt{\dfrac{1-\sin \alpha}{1+\sin \alpha}} \right)^2=\dfrac{1+\sin \alpha}{1-\sin \alpha}+\dfrac{1-\sin \alpha}{1+\sin \alpha}+2=\dfrac{4}{1-\sin ^2\alpha} =\dfrac{4}{\cos ^2\alpha}.\]
		Vì $0< \alpha <  \dfrac{\pi}{2}$ nên $\cos \alpha > 0$ do đó $A=\dfrac{2}{\cos \alpha}$.
	 }
\end{ex}

% nguồn: Ok_Toan11_TracNghiem_Lan2
\begin{ex}%[1D1N2-2]%[ma:cf06db9e]
	Cho biết  $\tan \alpha =\dfrac{1}{2} $. Tính $\cot \alpha $
	\choice 
		{\True $\cot \alpha  =2$}
		{$\cot \alpha =\dfrac{1}{4} $} 
		{$\cot \alpha =\dfrac{1}{2}  $}
		{$\cot \alpha  =\sqrt{2} $}
	 \loigiai{
	Ta có $\tan \alpha \cdot\cot \alpha  =1\Rightarrow \cot \alpha =\dfrac{1}{\tan \alpha} =\dfrac{1}{\dfrac{1}{2}} =2$.}
\end{ex}

% nguồn: Ok_Toan11_TracNghiem_Lan2
\begin{ex}%[1D1N2-2]%[ma:c529d290]
	Cho $0<\alpha <\dfrac{\pi}{2} $. Mệnh đề nào sau đây \textbf{sai}? 
	\choice 
		 {$\sin (\alpha +\pi)<0$}
		 {$\cos (\alpha +\pi)>0$}
		{\True $\tan (\pi -\alpha)>0$} 
		 {$\cot (\pi -\alpha)<0$}
	\loigiai{ 
		Ta có $\tan (\pi -\alpha)=-\tan \alpha <0$, $\forall  \alpha \in \left(0;\dfrac{\pi}{2} \right)$.
	 }
\end{ex}

% nguồn: Ok_Toan11_TracNghiem_Lan2
\begin{ex}%[1D1N2-2]%[ma:2a611eb0]
	Đơn giản biểu  thức $A=\cos \left(\alpha-\dfrac{\pi}{2} \right)+\sin (\alpha-\pi)$, ta được
	\choice 
		{$A=\cos \alpha+\sin  \alpha$}
		 {$A=2\sin \alpha$}
		{$A=\sin \alpha  \cos \alpha$}
		{\True $A=0$} 
	 \loigiai{
		Ta có $A=\cos \left(\alpha-\dfrac{\pi}{2} \right)+\sin \left(\alpha-\pi  \right)=\cos \left(\dfrac{\pi}{2}-\alpha \right)-\sin \left(\pi-\alpha \right)=\sin \alpha-\sin \alpha=0$.
	} 
\end{ex}

% nguồn: Ok_Toan11_TracNghiem_Lan2
\begin{ex}%[1D1N2-2]%[ma:94444fd9]
	Cho $\dfrac{7\pi}{4} <\alpha <2\pi $. Khẳng định nào sau đây  đúng?
	\choice 
		{\True $\cos \alpha  >0$}
		 {$\sin \alpha >0$}
		{$\tan \alpha  >0$}
		 {$\cot \alpha >0$}
	 \loigiai{
		Ta có $\dfrac{7\pi}{4} <\alpha <2\pi$ nên $\alpha$ thuộc cung phần tư  thứ $IV$.\\
	Do đó  $\cos \alpha >0$.}
\end{ex}

% nguồn: Ok_TracNghiemLop11_Lan1
\begin{ex}%[1D1N2-2]%[ma:3da6181e]
	 \immini{Trên đường tròn lượng giác, cho điểm $M(a;b)$. Góc lượng giác sđ$(OA,OM)=\alpha $. Chọn khẳng định đúng?
	\choice 
		{$\sin \alpha  =a$}
		{\True $\sin \alpha  =b$}
		{$\sin \alpha  =\dfrac{a}{b}$}
		{$\sin \alpha =\dfrac{b}{a}$}}{ 
		\begin{tikzpicture}[scale=0.6,font=\footnotesize,line join=round,line cap=round,>=stealth]
				\path (0,0) coordinate (O) 
					(3,0) coordinate (A) 
					(0,3) coordinate (B) 
					(0,-3) coordinate (B') 
					(-3,0)  coordinate (A')
					 (0:0)++(150:3) coordinate (M)
					($(O)!(M)!(A')$) coordinate  (a)
					 ($(O)!(M)!(B)$) coordinate (b);
				\draw[->]  (-4,0) -- (4,0) node[above,blue]{$x$};
				\draw[->] (0,-4) --  (0.,4) node[left,blue]{$y$};
				\draw[orange]  (O) circle (3cm);
				\draw[rotate=0,->,green!50!black] (0.5,0)  arc (0:150:0.5cm);
				 \draw (0.35,0.25) node[above,blue] {$\alpha$};
				\draw[dashed] (a)--(M)--(b); 
				\draw[green!50!black] (M)--(O); 
				\foreach \p/\r  in {A/-45,M/150,a/-90,O/-150,A'/-135,B'/-45,B/45,b/0}
				\fill (\p) circle (1pt)  node[shift={(\r:3mm)},blue]{$\p$};
			\end{tikzpicture}
		 }
	 \loigiai{
	Theo định nghĩa giá trị lượng giác của góc lượng giác nên đáp án là $\sin \alpha =b$.} 
\end{ex}

% nguồn: Ok_TracNghiemLop11_Lan1
\begin{ex}%[1D1N2-2]%[ma:e0e7b926]
	Cho $0<\alpha <\dfrac{\pi}{2}.$ Khẳng định  nào sau đây đúng?
	\choice 
		 {$\sin \alpha<0$}
		{$\cos  \alpha<0$}
		{$\tan \alpha<0$} 
		{\True $\cot \alpha>0$} 
	 \loigiai{
		Do $0<\alpha <\dfrac{\pi}{2}$  thì $\sin \alpha >0$ và $\cos\alpha >0$ nên $\cot \left(\alpha\right)=\dfrac{\cos \alpha}{\sin \alpha}>0$
	} 
\end{ex}

% nguồn: Ok_TracNghiemLop11_Lan1
\begin{ex}%[1D1N2-2]%[ma:86e92a6c]
	\immini[thm]{Cho góc hình học $uOv=30^\circ$. Xác  định số đo của góc lượng giác $(Ou,Ov)$.
	 \choice
		{  $\mathrm{\text{sđ}}(Ou,Ov)=-30^\circ+k360^\circ$, $k \in \mathbb{Z}$}
		 {\True $\mathrm{\text{sđ}}(Ou,Ov)=30^\circ+k360^\circ$, $k \in \mathbb{Z}$}
		{$\mathrm{\text{sđ}}(Ou,Ov)=150^\circ+k360^\circ$, $k \in \mathbb{Z}$} 
		{$\mathrm{\text{sđ}}(Ou,Ov)=-150^\circ+k360^\circ$, $k \in \mathbb{Z}$}}{\begin{tikzpicture}
		\draw (0,0)--(5,0); 
		\draw (0,0)--(30:4); 
		\draw  (5,0) node[below]{$u$}
			 (0,0) node[left]{$O$};
		 \draw (30:4) node[left]{$v$};
		 \draw (5,0) coordinate (A)--(0:0) coordinate (B)--(30:4) coordinate (C)
			pic [draw,->]{angle=A--B--C}; 
	\end{tikzpicture}
		} 
	\loigiai{ 
		Các góc lượng  giác tia đầu $Ou$, tia cuối $Ov$ có số đo là $\mathrm{\text{sđ}}(Ou,Ov)=30^\circ+k360^\circ$, $k \in \mathbb{Z}$.
	} 
\end{ex}

% nguồn: Ok_TracNghiemLop11_Lan1
\begin{ex}%[1D1N2-2]%[ma:98834b2e]
	Cho $\cos \alpha = \dfrac{5}{13}$ với $-\dfrac{\pi}{2} < \alpha < 0$. Giá trị  của $\sin \alpha$ là
	\choice 
		{$\dfrac{12}{13}$} 
		{\True  $-\dfrac{12}{13}$}
		 {$-\dfrac{5}{13}$}
		{$\dfrac{13}{5}$} 
	\loigiai{ 
		Ta có $\sin^2\alpha + \cos^2\alpha  = 1 \Leftrightarrow \sin^2\alpha = 1 - \cos^2\alpha = 1 - \left(\dfrac{5}{13}\right)^2 = \left(\dfrac{12}{13}\right)^2$.\\
		$\Leftrightarrow \sin\alpha = \dfrac{12}{13}$ hoặc  $\sin\alpha = -\dfrac{12}{13}$.\\
		Vì $-\dfrac{\pi}{2} < \alpha < 0$ nên $\sin\alpha  < 0$.\\
		Vậy $\sin\alpha = -\dfrac{12}{13}$. 
	} 
\end{ex}

% nguồn: Ok_TracNghiemLop11_Lan1
\begin{ex}%[1D1N2-2]%[ma:110d2eb9]
	 Cho $ 0<\alpha <\dfrac{\pi}{2}$. Mệnh đề nào sau đây sai?
	\choice 
		{$\sin (\alpha\,+\,\pi)<0$} 
		{$\cos (\alpha\,+\,\pi)<0$} 
		 {\True $\tan (\pi\,-\,\alpha)>0$}
		{$\cot  (\pi\,-\,\alpha)<0$}
	 \loigiai{$\tan (\pi-\alpha)=-\tan\alpha <0\,\forall\,\alpha\in\left(0\,;\,\dfrac{\pi}{2}\right)$.}
\end{ex}

% nguồn: Ok_TracNghiemLop11_Lan1
\begin{ex}%[1D1N2-2]%[ma:ce9ac395]
	Cho góc $\alpha $ thỏa mãn $ 2\pi <\alpha  <\dfrac{5\pi}{2}$. Khẳng định nào sau đây {\bf sai}?
	\choice 
		{\True  $\tan\alpha <0$}
		{$\cot\alpha >0$} 
		 {$\sin\alpha >0$}
		 {$\cos\alpha >0$}
	\loigiai{ Với $ 2\pi <\alpha <\dfrac{5\pi}{2}$ ta có $\sin\alpha >0$, $\cos\alpha >0$, $\tan\alpha >0$, $\cot\alpha >0$. 
	 }
\end{ex}

% nguồn: Ok_TracNghiemLop11_Lan1
\begin{ex}%[1D1N2-2]%[ma:8776ac07]
	Cho góc $\alpha$  thỏa mãn $0<\alpha<\dfrac{\pi}{2}$. Khẳng định nào sau đây đúng?
	\choice 
		{\True $\sin \alpha>0$} 
		 { $\cot \alpha<0$}
		 { $\sin \alpha<0$}
		{ $\cos \alpha<0$} 
	 \loigiai{
		Góc $\alpha$ thỏa mãn $0<\alpha<\dfrac{\pi}{2}$ nằm ở góc phần tư  số $I$.\\
	Do đó $\sin \alpha>0$; $\cos \alpha>0$; $\tan \alpha>0$;  $\cot \alpha>0$.}
\end{ex}

% nguồn: Ok_TracNghiemLop11_Lan1
\begin{ex}%[1D1N2-2]%[ma:4ae6bf6d]
	Cho  cấp số cộng $(u_n)$ có $u_1=-3$, $u_7=21$. Khi đó công sai $d$ là
	\choice 
		 {$3$}
		{\True  $4$}
		 {$24$}
		{$\dfrac{24}{7}$} 
	\loigiai{Áp dụng công  thức $u_n=u_1+(n-1)d$. Khi đó
		 $$
		u_7=u_1+6d \Rightarrow d=\dfrac{u_7-u_1}{6}=\dfrac{21-(-3)}{6}=4. 
	 $$}
\end{ex}

% nguồn: Ok_TracNghiemLop11_Lan1
\begin{ex}%[1D1N2-2]%[ma:c20a7417]
	Cho $-\pi<\alpha<-\dfrac{\pi}{2}$. Khẳng  định nào sau đây \textbf{đúng}?
	\choice 
		{$\sin \alpha>0$, $\cos  \alpha< 0$}
		{\True $\cos \alpha<0$,  $\tan \alpha>0$}
		{$\sin \alpha<0$, $\tan \alpha<0$} 
		 {$\tan \alpha>0$, $\cot \alpha<0$}
	\loigiai{ 
		Vì $-\pi<\alpha<-\dfrac{\pi}{2}$ nên điểm biểu diễn của góc lượng giác $\alpha$ thuộc  góc phần tư thứ III. Do đó $\sin \alpha<0$ và $\cos \alpha<0$. Khi đó $\tan \alpha >0$.
	} 
\end{ex}

% nguồn: Ok_TracNghiemLop11_Lan1
\begin{ex}%[1D1N2-2]%[ma:4ae63314]
	 Cho $\pi <\alpha<\dfrac{3\pi}{2}$. Khẳng định nào sau đây là đúng?
	\choice 
		{$\sin \alpha >0$} 
		{$\cos \alpha >0$} 
		 {\True $\sin \left(\dfrac{\pi}{2}-\alpha\right) <0$}
		 {$\cos \left(\dfrac{\pi}{2}-\alpha\right) >0$}
	\loigiai{Ta có $\pi <\alpha<\dfrac{3\pi}{2}$ nên $\sin \alpha  <0$ và $\cos \alpha <0$.\\
		Mà $\sin  \left(\dfrac{\pi}{2}-\alpha\right) =\cos \alpha<0$ nên
		$\cos \left(\dfrac{\pi}{2}-\alpha\right) =\sin  \alpha<0$.
	} 
\end{ex}

% nguồn: Ok_Toan11_TracNghiem_Lan2
\begin{ex}%[1D1H2-2]%[ma:d19bb816]
	 Nếu $\sin x=3\cos x$ thì $\sin x\cos x$ bằng
	 \choice
		 {\True $\dfrac{3}{10}$}
		 {$\dfrac{2}{9}$}
		{$\dfrac{1}{4}$} 
		 {$\dfrac{1}{6}$}
	\loigiai{ 
		Ta có $\sin^2x+\cos^2x=1\Rightarrow  \left(3\cos x\right)^2+\cos^2x=1\Rightarrow \cos^2x=\dfrac{1}{10}\Rightarrow \sin^2x=\dfrac{9}{10}$.\\
		Khi đó \[\sin x=3\cos x\Rightarrow \left(\sin x-3\cos x\right)^2=0\Leftrightarrow \sin^2x+9\cos^2x-6\sin x\cos x=0\Rightarrow \sin x\cdot \cos x=\dfrac{3}{10}.\] 
	} 
\end{ex}

% nguồn: Ok_Toan11_TracNghiem_Lan2
\begin{ex}%[1D1H2-2]%[ma:a7a2921a]
	Cho $\dfrac{\pi}{2} <\alpha <\pi$. Giá trị  lượng giác nào sau đây luôn dương?
	\choice 
		 {$\sin (\pi +\alpha)$}
		{\True $\cot \left(\dfrac{\pi}{2}  -\alpha \right)$}
		{$\cos  (-\alpha)$}
		{$\tan (\pi +\alpha)$} 
	 \loigiai{
		Ta  có
		\begin{itemize}
			\item $\sin (\pi +\alpha)=-\sin \alpha $; 
			\item $\cot \left(\dfrac{\pi}{2} -\alpha  \right)=\sin \alpha $;
			 \item $\cos (-\alpha)=\cos \alpha $;
			\item $\tan (\pi  +\alpha)=\tan \alpha$.
		\end{itemize}
		Do $\dfrac{\pi}{2} <\alpha <\pi$ suy  ra
		\begin{itemize}
			\item $\sin \alpha >0$. 
			\item $\cos \alpha  <0$.
			\item $\tan \alpha  <0$.
		\end{itemize}
	 }
\end{ex}

% nguồn: Ok_Toan11_TracNghiem_Lan2
\begin{ex}%[1D1H2-2]%[ma:b5b6676c]
	Trên đường  tròn lượng giác, lấy điểm $M\left(-\dfrac{\sqrt{3}}{4};-\dfrac{\sqrt{13}}{4} \right)$ có góc lượng giác $(OA,OM)$ có số đo $\alpha+k2\pi$. Khi đó, giá trị của $\sin \alpha$ là giá trị nào trong các giá trị sau?
	 \choice
		{\True  $-\dfrac{\sqrt{13}}{4}$}
		 {$-\dfrac{\sqrt{3}}{4}$}
		{$\dfrac{\sqrt{3}}{4}$} 
		{$\dfrac{\sqrt{13}}{4}$} 
	\loigiai{ 
		Điểm $M\left(-\dfrac{\sqrt{3}}{4};-\dfrac{\sqrt{13}}{4} \right)$ có góc lượng giác $\left(OA;OM\right)$ có số đo $\alpha+k2\pi$.  \\
		Khi  đó, giá trị của $\sin \alpha=y_M=-\dfrac{\sqrt{13}}{4}$.
	} 
\end{ex}

% nguồn: Ok_Toan11_TracNghiem_Lan2
\begin{ex}%[1D1H2-2]%[ma:243cc58e]
	Cho góc $\alpha $ thỏa  $\cot\alpha =\dfrac{3}{4} $ và $0^\circ <\alpha <90^\circ$. Khẳng định nào sau đây đúng?
	\choice 
		{$\cos\alpha =\dfrac{4}{5} $} 
		{\True $\sin \alpha =\dfrac{4}{5}  $}
		{$\sin\alpha =-\dfrac{4}{5} $} 
		{$\cos\alpha =-\dfrac{4}{5} $} 
	\loigiai{ 
		Áp dụng công thức $1+\cot ^{2} \alpha =\dfrac{1}{\sin ^{2} \alpha} $ $\Leftrightarrow  1+\dfrac{9}{16} =\dfrac{1}{\sin ^{2} \alpha} $ $\Leftrightarrow \sin ^{2} \alpha =\dfrac{16}{25} $.\\
	Do $0^\circ <\alpha <90^\circ \Rightarrow \sin \alpha >0 \Rightarrow \sin \alpha =\dfrac{4}{5}  $.}
\end{ex}

% nguồn: Ok_Toan11_TracNghiem_Lan2
\begin{ex}%[1D1H2-2]%[ma:91ff6adb]
	Cho $0<\alpha <\dfrac{\pi}{2}$. Khẳng định nào  sau đây đúng?
	 \choice
		{$\cot \left(\alpha +\dfrac{\pi}{2}  \right)>0$}
		{$\cot  \left(\alpha +\dfrac{\pi}{2} \right)\ge 0$}
		{$\tan (\alpha  +\pi)<0$}
		{\True  $\tan (\alpha +\pi)>0$}
	 \loigiai{
		Ta  có
		\begin{itemize}
			\item $0<\alpha <\dfrac{\pi}{2} \Rightarrow \dfrac{\pi}{2} <\alpha  +\dfrac{\pi}{2} <\pi \Rightarrow \cot \left(\alpha +\dfrac{\pi}{2} \right)<0$.
			\item $0<\alpha <\dfrac{\pi}{2} \Rightarrow \pi <\alpha +\pi <\dfrac{3\pi}{2}  \Rightarrow \tan (\alpha +\pi)>0$.
		\end{itemize}
	 }
\end{ex}

% nguồn: Ok_Toan11_TracNghiem_Lan2
\begin{ex}%[1D1H2-2]%[ma:0d923685]
	Cho $\sin \alpha =\dfrac{3}{5} \left(\dfrac{\pi}{2}  <\alpha <\pi \right)$. Giá trị của $\cos \alpha $ bằng
	 \choice
		 {\True $\cos \alpha =-\dfrac{4}{5} $}
		{$\cos \alpha =\dfrac{4}{5} $} 
		{$\cos \alpha  =\dfrac{2}{5} $}
		{$\cos \alpha  =-\dfrac{2}{5} $}
	\loigiai{ 
		Ta có $\sin ^{2} \alpha +\cos ^{2} \alpha =1\Rightarrow \cos ^{2} \alpha =1-\sin ^{2} \alpha =1-\dfrac{9}{25}  =\dfrac{16}{25} $.\\
	Mà $\dfrac{\pi}{2} <\alpha <\pi \Rightarrow  \cos \alpha =-\dfrac{4}{5} $.}
\end{ex}

% nguồn: Ok_Toan11_TracNghiem_Lan2
\begin{ex}%[1D1H2-2]%[ma:e0e488d0]
	Cho  $\sin \alpha =\dfrac{4}{5} $, $(90^\circ <\alpha <180^\circ)$. Tính $\cos \alpha $.
	 \choice
		{$\cos  \alpha =-\dfrac{4}{5} $}
		{\True $\cos \alpha =-\dfrac{3}{5} $} 
		{$\cos \alpha =\dfrac{5}{3} $} 
		{$\cos \alpha =\dfrac{3}{5} $} 
	 \loigiai{
		Ta  có
		\begin{eqnarray*}
			&&\sin ^{2} \alpha +\cos ^{2}  \alpha =1\\
			&\Leftrightarrow& \cos ^{2} \alpha  =1-\sin ^{2} \alpha =\dfrac{9}{25} \\
			&\Rightarrow& \cos \alpha =\pm  \dfrac{3}{5}.
		\end{eqnarray*}
	Vì $90^\circ <\alpha  <180^\circ $ nên $\cos \alpha =-\dfrac{3}{5} $.}
\end{ex}

% nguồn: Ok_Toan11_TracNghiem_Lan2
\begin{ex}%[1D1H2-2]%[ma:0e44e38a]
	 Cho $\pi <\alpha <\dfrac{3\pi}{2}$. Khẳng định nào sau đây đúng?
	\choice 
		{$\tan \left(\dfrac{3\pi}{2} -\alpha  \right)<0$}
		{\True  $\tan \left(\dfrac{3\pi}{2} -\alpha \right)>0$}
		{$\tan \left(\dfrac{3\pi}{2} -\alpha  \right)\le 0$}
		{$\tan \left(\dfrac{3\pi}{2} -\alpha \right)\ge  0$}
	\loigiai{ 
		 Ta có $\pi <\alpha <\dfrac{3\pi}{2} \Rightarrow 0<\dfrac{3\pi}{2} -\alpha <\dfrac{\pi}{2} \Rightarrow \heva{&\sin \left(\dfrac{3\pi}{2} -\alpha \right)>0\\ &\cos \left(\dfrac{3\pi}{2} -\alpha \right)>0}\Rightarrow \tan \left(\dfrac{3\pi}{2} -\alpha \right)>0$.
	} 
\end{ex}

% nguồn: Ok_Toan11_TracNghiem_Lan2
\begin{ex}%[1D1H2-2]%[ma:39f7b3d7]
	Cho $\dfrac{7\pi}{4} < \alpha < 2\pi$. Khẳng định  nào sau đây \textbf{đúng}?
	\choice 
		{\True $\cos \alpha >  0$}
		{$\sin \alpha > 0$} 
		{$\tan  \alpha > 0$}
		 {$\cot \alpha > 0$}
	\loigiai{ 
		 Ta có $\dfrac{7\pi}{4} < \alpha < 2\pi \Leftrightarrow \dfrac{3\pi}{2}+\dfrac{\pi}{4} < \alpha < 2\pi$.\\
		Khi đó $  \alpha $ thuộc cung phần tư thứ IV vì vậy $\cos \alpha > 0$.
	} 
\end{ex}

% nguồn: Ok_Toan11_TracNghiem_Lan2
\begin{ex}%[1D1H2-2]%[ma:cb6d22d9]
	Trên đường  tròn lượng giác, lấy điểm $M\left(\dfrac{\sqrt{5}}{4};-\dfrac{\sqrt{11}}{4} \right)$ có góc lượng giác $(OA, OM)$ có số đo $\alpha+k2\pi$. Khi đó, giá trị của $\sin \alpha$ là giá trị nào trong các giá trị sau?
	 \choice
		 {\True $-\dfrac{\sqrt{11}}{4}$}
		{$-\dfrac{\sqrt{5}}{4}$} 
		 {$\dfrac{\sqrt{5}}{4}$}
		 {$\dfrac{\sqrt{11}}{4}$}
	 \loigiai{
		Điểm $M\left(\dfrac{\sqrt{5}}{4};-\dfrac{\sqrt{11}}{4} \right)$ có góc lượng giác $(OA,OM)$ có số đo $\alpha+k2\pi$.\\ 
		Khi đó, giá trị của $\sin \alpha=y_M=-\dfrac{\sqrt{11}}{4}$. 
	} 
\end{ex}

% nguồn: Ok_Toan11_TracNghiem_Lan2
\begin{ex}%[1D1H2-2]%[ma:b6e567db]
	 Rút gọn biểu thức $M = \cos^4 15^\circ - \sin^4 15^\circ$.
	\choice 
		{$M=1$} 
		{\True  $M=\dfrac{\sqrt{3}}{2}$}
		{$M=\dfrac{1}{4}$} 
		{$M=0$} 
	\loigiai{ 
		Ta có $M = \cos^4 15^\circ  - \sin^4 15^\circ = (\cos^2 15^\circ)^2 - (\sin^2 15^\circ)^2$.\\
		$= (\cos^2 15^\circ  - \sin^2 15^\circ)(\cos^2 15^\circ + \sin^2 15^\circ) = \cos^2 15^\circ - \sin^2 15^\circ = \cos(2 \cdot 15^\circ) = \cos 30^\circ = \dfrac{\sqrt{3}}{2}$.
	 }
\end{ex}

% nguồn: Ok_Toan11_TracNghiem_Lan2
\begin{ex}%[1D1H2-2]%[ma:6920fb08]
	Cho $\sin \alpha=-\dfrac{3}{5}$ và $\pi < \alpha < \dfrac{3\pi}{2}$. Khi đó giá trị của $\cos\alpha$ và $\tan \alpha$ lần  lượt là
	 \choice
		 {\True $-\dfrac{4}{5}; \dfrac{3}{4}$}
		{$-\dfrac{4}{5};-\dfrac{3}{4}$} 
		{$\dfrac{4}{5};-\dfrac{3}{4}$} 
		 {$\dfrac{3}{4};-\dfrac{4}{5}$}
	\loigiai{ 
		Ta có $\cos^2\alpha=1-\sin^2\alpha = 1-\dfrac{9}{25}=\dfrac{16}{25} \Rightarrow \cos\alpha=\pm\dfrac{4}{5}$.\\ 
		Do $\pi < \alpha < \dfrac{3\pi}{2}$ nên $\cos \alpha  < 0\Rightarrow \cos \alpha=-\dfrac{4}{5}$.\\
		Khi đó $\tan\alpha=\dfrac{\sin\alpha}{\cos\alpha}=\dfrac{-\dfrac{3}{5}}{-\dfrac{4}{5}}=\dfrac{3}{4}$.\\  Vậy $\cos\alpha=-\dfrac{4}{5}; \tan\alpha=\dfrac{3}{4}$.
	 }
\end{ex}

% nguồn: Ok_Toan11_TracNghiem_Lan2
\begin{ex}%[1D1H2-2]%[ma:f4fca742]
	Cho $0<\alpha  <\dfrac{\pi}{2}$. Khẳng định nào sau đây đúng?
	\choice 
		{$\sin (\alpha  -\pi)\ge 0$}
		{$\sin (\alpha -\pi)\le 0$} 
		{$\sin (\alpha  -\pi)>0$}
		{\True $\sin (\alpha  -\pi)<0$}
	 \loigiai{
		Ta có $0<\alpha <\dfrac{\pi}{2} \Rightarrow -\pi <\alpha  -\pi <-\dfrac{\pi}{2}$.\\
	Điểm cuối cung $\alpha -\pi $ thuộc góc phần tư thứ $III$ suy  ra $\sin (\alpha -\pi)<0$.}
\end{ex}

% nguồn: Ok_Toan11_TracNghiem_Lan2
\begin{ex}%[1D1H2-2]%[ma:2b76f196]
	Trên đường tròn lượng  giác, lấy điểm $M\left(\dfrac{\sqrt{5}}{4};-\dfrac{\sqrt{11}}{4} \right)$ có góc lượng giác $(OA,OM)$ có số đo $\alpha+k2\pi$. Khi đó, giá trị của $\cos \alpha$ là giá trị nào trong các giá trị sau?
	\choice 
		 {$-\dfrac{\sqrt{11}}{4}$}
		 {$-\dfrac{\sqrt{5}}{4}$}
		{\True $\dfrac{\sqrt{5}}{4}$} 
		 {$\dfrac{\sqrt{11}}{4}$}
	 \loigiai{
		Điểm $M\left(\dfrac{\sqrt{5}}{4};-\dfrac{\sqrt{11}}{4} \right)$ có góc lượng giác $(OA,OM)$ có số đo  $\alpha+k2\pi$.\\
		Khi đó, giá trị của  $\cos \alpha=x_M=\dfrac{\sqrt{5}}{4}$.
	} 
\end{ex}

% nguồn: Ok_Toan11_TracNghiem_Lan2
\begin{ex}%[1D1H2-2]%[ma:4fe21021]
	Tính $\sin  \alpha$, biết $\cos \alpha=\dfrac{\sqrt{5}}{3}$ và $\dfrac{3\pi}{2} < \alpha < 2\pi$.
	 \choice
		 {$\dfrac{1}{3}$}
		{$-\dfrac{1}{3}$} 
		 {$\dfrac{2}{3}$}
		 {\True $-\dfrac{2}{3}$}
	 \loigiai{
		Ta có $\sin ^2\alpha=1-\cos ^2\alpha=1-\dfrac{5}{9}=\dfrac{4}{9}$ $\Leftrightarrow \sin \alpha=\pm \dfrac{2}{3}$.\\ 
		Do $\dfrac{3\pi}{2} < \alpha < 2\pi$ nên  $\sin \alpha < 0$.\\
		 Vậy $\sin \alpha=-\dfrac{2}{3}$.
	} 
\end{ex}

% nguồn: Ok_Toan11_TracNghiem_Lan2
\begin{ex}%[1D1H2-2]%[ma:8562ee07]
	Trên đường tròn lượng giác, lấy điểm $M\left(\dfrac{\sqrt{5}}{4};-\dfrac{\sqrt{11}}{4} \right)$ có góc lượng giác $(OA,OM)$ có số đo $\alpha+k2\pi$. Khi đó, giá trị của $\tan \alpha$ là giá trị  nào trong các giá trị sau?
	\choice 
		 {$\dfrac{\sqrt{55}}{11}$}
		{$-\dfrac{\sqrt{55}}{11}$} 
		 {\True $-\dfrac{\sqrt{55}}{5}$}
		 {$\dfrac{\sqrt{55}}{5}$}
	\loigiai{ 
		Điểm $M\left(\dfrac{\sqrt{5}}{4};-\dfrac{\sqrt{11}}{4} \right)$ có góc  lượng giác $(OA,OM)$ có số đo $\alpha+k2\pi$. \\
		Khi đó, giá trị của $\tan  \alpha=\dfrac{y_M}{x_M}=-\dfrac{\sqrt{11}}{4}:\dfrac{\sqrt{5}}{4}=-\dfrac{\sqrt{55}}{5}$.
	} 
\end{ex}

% nguồn: Ok_Toan11_TracNghiem_Lan2
\begin{ex}%[1D1H2-2]%[ma:a426e34f]
	Cho  $\pi <\alpha <\dfrac{3\pi}{2}$. Mệnh đề nào sau đây đúng?
	\choice 
		{$\sin (\pi  +\alpha)<0$}
		 {\True $\sin (\pi -\alpha)<0$}
		{$\cos (\pi  -\alpha)<0$}
		{$\cos (\pi +\alpha)<0$} 
	 \loigiai{
		Ta có $\pi <\alpha <\dfrac{3\pi}{2} \Rightarrow 2\pi <\alpha +\pi <\dfrac{5\pi}{2} \Rightarrow \sin  (\pi +\alpha)>0$; $\cos (\pi +\alpha)>0$.\\
	Mặt khác $\pi <\alpha <\dfrac{3\pi}{2}  \Rightarrow -\dfrac{\pi}{2} <\pi -\alpha <0\Rightarrow \sin (\pi -\alpha)<0$; $\cos (\pi -\alpha)>0$.}
\end{ex}

% nguồn: Ok_Toan11_TracNghiem_Lan2
\begin{ex}%[1D1H2-2]%[ma:eb407428]
	Cho $3\pi < \alpha < \dfrac{10\pi}{3}$. Khẳng  định nào sau đây \textbf{đúng}?
	 \choice
		{$\cos  \alpha > 0$}
		{\True $\sin  \alpha < 0$}
		 {$\tan \alpha < 0$}
		{$\cot \alpha < 0$} 
	 \loigiai{
		Ta có $3\pi < \alpha < \dfrac{10\pi}{3} \Leftrightarrow 2\pi+\pi  < \alpha < 2\pi+\pi+\dfrac{\pi}{3}$.\\
		Khi đó $\alpha$ thuộc cung phần tư thứ  III vì vậy $\sin \alpha < 0$.
	 }
\end{ex}

% nguồn: Ok_Toan11_TracNghiem_Lan2
\begin{ex}%[1D1H2-2]%[ma:dcad0b4b]
	Cho $\sin \alpha=\dfrac{1}{4}$, với $0^\circ< \alpha < 90^\circ$. Tính $\cos \alpha$; $\tan \alpha$. 
	 \choice
		{$\cos  \alpha=-\dfrac{\sqrt{15}}{4};\tan \alpha=\dfrac{\sqrt{15}}{15}$}
		{$\cos \alpha=-\dfrac{\sqrt{15}}{4};\tan  \alpha=-\dfrac{\sqrt{15}}{15}$}
		{$\cos \alpha=\dfrac{\sqrt{15}}{4};\tan \alpha=-\dfrac{\sqrt{15}}{15}$} 
		 {\True $\cos \alpha=\dfrac{\sqrt{15}}{4};\tan \alpha=\dfrac{\sqrt{15}}{15}$}
	 \loigiai{
		Ta có $\cos ^2\alpha=1-\sin ^2\alpha=\dfrac{15}{16} \Rightarrow \cos  \alpha=\pm \dfrac{\sqrt{15}}{4}$.\\
		Do $0^\circ< \alpha < 90^\circ$ nên $\cos \alpha=\dfrac{\sqrt{15}}{4}$.\\ 
		 Ta có $\tan \alpha=\dfrac{\sin \alpha}{\cos \alpha}=\dfrac{1}{\sqrt{15}}=\dfrac{\sqrt{15}}{15}$.
	 }
\end{ex}

% nguồn: Ok_Toan11_TracNghiem_Lan2
\begin{ex}%[1D1H2-2]%[ma:b9ca9a5f]
	Trên đường tròn lượng giác, lấy điểm $M\left(-\dfrac{3}{5};\dfrac{4}{5} \right)$ có góc lượng giác $(OA,OM)$ có số đo $\alpha+k2\pi$.  Khi đó, giá trị của $\cot \alpha$ là giá trị nào trong các giá trị sau?
	\choice 
		{$\dfrac{3}{4}$} 
		 {\True $-\dfrac{3}{4}$}
		{$-\dfrac{4}{3}$} 
		{$\dfrac{4}{3}$} 
	 \loigiai{
		Điểm $M\left(-\dfrac{3}{5};\dfrac{4}{5} \right)$ có góc lượng giác $(OA,OM)$ có số  đo $\alpha+k2\pi$. \\
		Khi đó, giá trị của  $\cot \alpha=\dfrac{x_M}{y_M}=-\dfrac{3}{5}:\dfrac{4}{5}=-\dfrac{3}{4}$.
	 }
\end{ex}

% nguồn: Ok_Toan11_TracNghiem_Lan2
\begin{ex}%[1D1H2-2]%[ma:f0d17009]
	Trên đường tròn lượng giác, lấy điểm $M\left(-\dfrac{3}{5};\dfrac{4}{5} \right)$ có góc lượng giác $(OA,OM)$ có số  đo $\alpha+k2\pi$. Khi đó, giá trị của $\tan \alpha$ là giá trị nào trong các giá trị sau?
	\choice 
		 {$\dfrac{3}{4}$}
		{$-\dfrac{3}{4}$} 
		 {\True $-\dfrac{4}{3}$}
		{$\dfrac{4}{3}$} 
	\loigiai{ 
		Điểm $M\left(-\dfrac{3}{5};\dfrac{4}{5} \right)$  có góc lượng giác $(OA,OM)$ có số đo $\alpha+k2\pi$. \\
		Khi đó, giá trị của  $\tan \alpha=\dfrac{y_M}{x_M}=\dfrac{4}{5}:\left(-\dfrac{3}{5} \right)=-\dfrac{4}{3}$.
	 }
\end{ex}

% nguồn: Ok_Toan11_TracNghiem_Lan2
\begin{ex}%[1D1H2-2]%[ma:14d13bb8]
	Trên đường tròn lượng giác, lấy điểm $M\left(-\dfrac{3}{5};\dfrac{4}{5} \right)$ có góc lượng giác $(OA,OM)$ có số đo $\alpha+k2\pi$. Khi đó, giá trị của $\sin \alpha$ là giá trị nào trong các giá  trị sau?
	\choice 
		{\True  $\dfrac{4}{5}$}
		 {$-\dfrac{4}{5}$}
		 {$\dfrac{3}{5}$}
		 {$-\dfrac{3}{5}$}
	 \loigiai{
		Điểm $M\left(-\dfrac{3}{5};\dfrac{4}{5}  \right)$ có góc lượng giác $(OA,OM)$ có số đo $\alpha+k2\pi$.\\
		Khi  đó, giá trị của $\sin \alpha=y_M=\dfrac{4}{5}$.
	} 
\end{ex}

% nguồn: Ok_Toan11_TracNghiem_Lan2
\begin{ex}%[1D1H2-2]%[ma:7b8ac34f]
	Cho $\sin\alpha =\dfrac{4}{5} $ với $\dfrac{\pi}{2} <\alpha <\pi  $. Tính $\cos\alpha $?
	\choice 
		 {$\dfrac{1}{5} $}
		{$\dfrac{3}{5} $} 
		{\True $-\dfrac{3}{5}  $}
		{$\dfrac{4}{5} $} 
	\loigiai{ 
		Ta có 
		\begin{eqnarray*}
			 &&\sin ^{2} \alpha +\cos ^{2} \alpha =1\\
			&\Leftrightarrow&  \cos ^{2} \alpha =1-\sin ^{2} \alpha =\dfrac{9}{25} \\
			&\Rightarrow& \cos \alpha =\pm \dfrac{3}{5}. 
		\end{eqnarray*}
	Vì $\dfrac{\pi}{2} <\alpha <\pi$ nên  $\cos \alpha =-\dfrac{3}{5} $.}
\end{ex}

% nguồn: Ok_Toan11_TracNghiem_Lan2
\begin{ex}%[1D1H2-2]%[ma:9821727a]
	Cho $3\cos x+2\sin  x=2$ và $\sin x<0$, khẳng định nào sau đây đúng?
	 \choice
		 {\True $\sin x=-\dfrac{5}{13}$}
		{$\sin  x=-\dfrac{7}{13}$}
		 {$\sin x=-\dfrac{9}{13}$}
		{$\sin x=-\dfrac{12}{13}$} 
	\loigiai{ 
		 Ta có
		\allowdisplaybreaks\begin{eqnarray*}
		3\cos  x+2\sin x=2&\Leftrightarrow&\left(3\cos x+2\sin x\right)^2=4\\
		&\Leftrightarrow &9\cos^2x+12\cos  x\cdot\sin x+4\sin^2x=4\\
		&\Leftrightarrow& 5\cos^2x+12\cos  x\cdot \sin x=0\\
		&\Leftrightarrow &\cos x\left(5\cos  x+12\sin x\right)=0\\
		&\Leftrightarrow& \hoac{&\cos x=0 \\&5\cos  x+12\sin x=0.}
	\end{eqnarray*}
	\begin{itemize}
		\item Với $\cos x=0\Rightarrow\sin x=1$  (loại vì $\sin x<0$).
		\item Với $5\cos x+12\sin x=0$, ta có  hệ phương trình
		\[\heva{&5\cos x+12\sin x=0 \\&3\cos x+2\sin  x=2 }\Leftrightarrow \heva{&\sin x=-\dfrac{5}{13} \\&\cos x=\dfrac{12}{13}.}\]
	\end{itemize}
	Vậy  $\sin x=-\dfrac{5}{13}$.
 }
\end{ex}

% nguồn: Ok_Toan11_TracNghiem_Lan2
\begin{ex}%[1D1H2-2]%[ma:5a04e4a1]
	 Nếu $\tan \alpha=\dfrac{3}{4}$ thì $\sin ^2\alpha$ bằng
	\choice 
		 {$\dfrac{16}{25}$}
		{\True  $\dfrac{9}{25}$}
		 {$\dfrac{25}{16}$}
		 {$\dfrac{25}{9}$}
	 \loigiai{
		 Ta có $\dfrac{1}{\cos ^2\alpha}=1+\tan ^2\alpha=1+\left(\dfrac{3}{4} \right)^2=\dfrac{25}{16} \Rightarrow \cos ^2\alpha=\dfrac{16}{25}$.\\
		$\Rightarrow \sin  ^2\alpha=1-\cos ^2\alpha=1-\dfrac{16}{25}=\dfrac{9}{25}$.
	 }
\end{ex}

% nguồn: Ok_Toan11_TracNghiem_Lan2
\begin{ex}%[1D1H2-2]%[ma:689ddb7e]
	Cho $\cos\alpha=\dfrac{-1}{6}$ với $ -\pi < \alpha  < \dfrac{-\pi}{2} $. Tính $\sin\alpha$.
	 \choice
		{\True $\sin \alpha=\dfrac{-\sqrt{35}}{6}$} 
		{$\sin  \alpha=\dfrac{35}{36}$}
		{$\sin  \alpha=\dfrac{5}{6}$}
		 {$\sin \alpha=\dfrac{\sqrt{35}}{6}$}
	\loigiai{ 
		Ta có $-\pi < \alpha <  \dfrac{-\pi}{2} \Rightarrow \sin \alpha < 0$.\\
		Ta lại có $ \sin^2 \alpha = 1 - \cos^2 \alpha  = 1-\left(\dfrac{-1}{6} \right)^2 = \dfrac{35}{36}$.\\
		Mà $ \sin \alpha < 0 $ nên $ \sin \alpha  = -\dfrac{\sqrt{35}}{6}$.
	 }
\end{ex}

% nguồn: Ok_Toan11_TracNghiem_Lan2
\begin{ex}%[1D1H2-2]%[ma:42b799b6]
	Trên đường tròn lượng giác, lấy điểm $M\left(\dfrac{\sqrt{5}}{4};-\dfrac{\sqrt{11}}{4} \right)$ có góc lượng giác $(OA,OM)$ có số đo $\alpha+k2\pi$. Khi đó, giá trị của $\cot \alpha$ là giá  trị nào trong các giá trị sau?
	\choice 
		 {$\dfrac{\sqrt{55}}{11}$}
		 {\True $-\dfrac{\sqrt{55}}{11}$}
		{$-\dfrac{\sqrt{55}}{5}$} 
		 {$\dfrac{\sqrt{55}}{5}$}
	\loigiai{ 
		Điểm $M\left(\dfrac{\sqrt{5}}{4};-\dfrac{\sqrt{11}}{4} \right)$  có góc lượng giác $(OA,OM)$ có số đo $\alpha+k2\pi$.\\
		Khi đó,  giá trị của $\cot \alpha=\dfrac{x_M}{y_M}=\dfrac{\sqrt{5}}{4}:\left(-\dfrac{\sqrt{11}}{4} \right)=-\dfrac{\sqrt{55}}{11}$.
	 }
\end{ex}

% nguồn: Ok_Toan11_TracNghiem_Lan2
\begin{ex}%[1D1H2-2]%[ma:7af0ac66]
	Cho $\sin \alpha=\dfrac{3}{5}$, với  $\dfrac{\pi}{2} < \alpha < \pi$. Giá trị của $\cos\alpha$ là
	 \choice
		{$\dfrac{4}{5}$} 
		{\True  $-\dfrac{4}{5}$}
		{$\pm \dfrac{4}{5}$} 
		{$\dfrac{16}{25}$} 
	\loigiai{ 
		Do $\dfrac{\pi}{2} < \alpha  < \pi$ nên $\cos\alpha<0$.\\
		Ta có $\sin ^2\alpha+\cos ^2\alpha=1$ $\Rightarrow  \cos ^2\alpha=1-\sin^2\alpha=1-\dfrac{9}{25}=\dfrac{16}{25}$.\\
		Suy ra $\cos\alpha=-\dfrac{4}{5}$. 
	 }
\end{ex}

% nguồn: Ok_Toan11_TracNghiem_Lan2
\begin{ex}%[1D1H2-2]%[ma:28423d81]
	Bất đẳng thức nào dưới đây  là \textbf{đúng}?
	 \choice
		{$\sin 90^\circ <  \sin 150^\circ$}
		{$\sin 90^\circ  15' < \sin 90^\circ 30'$}
		 {\True $\cos 90^\circ 30' > \cos 100^\circ$}
		{$\cos 150^\circ  > \cos 120^\circ$}
	\loigiai{ 
		Các góc trong đề bài thuộc $ \left[90^\circ; 180^\circ \right] $ tức thuộc góc phần tư thứ II  nên khi góc $ a<b $ thì $ \sin a > \sin b $ và $ \cos a > \cos b $.\\
		Từ đó dễ  nhận thấy $\cos 90^\circ 30' > \cos 100^\circ$ là phương án đúng.
	} 
\end{ex}

% nguồn: Ok_Toan11_TracNghiem_Lan2
\begin{ex}%[1D1H2-2]%[ma:7b8eb7c1]
	Cho $0< \alpha < \dfrac{\pi}{2}$. Khẳng định nào sau  đây \textbf{đúng}?
	 \choice
		{$\cot \left(\alpha+\dfrac{\pi}{2} \right) >  0$}
		{$\cot \left(\alpha+\dfrac{\pi}{2} \right)\ge  0$}
		 {$\tan \left(\alpha+\pi \right) < 0$}
		{\True $\tan \left(\alpha+\pi \right) >  0$}
	 \loigiai{
		Ta có $\heva{&{0< \alpha  < \dfrac{\pi}{2}} \\
		 &{0< \alpha < \dfrac{\pi}{2}}
		 } \Rightarrow
		\heva{&\dfrac{\pi}{2} <  \alpha+\dfrac{\pi}{2} < \pi \\ &\pi < \alpha+\pi < \dfrac{3\pi}{2} } \Rightarrow
		\heva{&\cot \left(\alpha+\dfrac{\pi}{2} \right) < 0 \\ &\tan  \left(\alpha+\pi \right) > 0.} $
	} 
\end{ex}

% nguồn: Ok_TracNghiemLop11_Lan1
\begin{ex}%[1D1H2-2]%[ma:0875acc3]
	Cho $\sin\alpha=-\dfrac{4}{5}$ và $\pi<\alpha<\dfrac{3\pi}{2}$. Giá trị của $\cos\alpha$ là 
	\choice 
		{$\dfrac{3}{5}$} 
		{\True $-\dfrac{3}{5}$} 
		 {$-\dfrac{9}{25}$}
		 {$\dfrac{9}{25}$}
	\loigiai{ 
		 Ta có $\sin^2\alpha+\cos^2\alpha=1\Rightarrow\cos^2\alpha=1-\sin^2\alpha=1-\left(-\dfrac{4}{5}\right)^2=\dfrac{9}{25}\Leftrightarrow\hoac{&\cos\alpha=\dfrac{3}{5}\\&\cos\alpha=-\dfrac{3}{5}.}$\\
	Vì  $\pi<\alpha<\dfrac{3\pi}{2}$ nên $\cos\alpha<0\Rightarrow\cos\alpha=-\dfrac{3}{5}$.}
\end{ex}

% nguồn: Ok_TracNghiemLop11_Lan1
\begin{ex}%[1D1H2-2]%[ma:261dd639]
	Biết rằng $\tan \alpha = 2$. Giá trị biểu  thức $\dfrac{\sin \alpha +2 \cos \alpha}{3\sin \alpha - \cos \alpha}\, (\cos \alpha \ne 0)$ là
	\choice 
		 {$\dfrac{4}{5}$}
		{$1$} 
		{$\dfrac{3}{5}$} 
		{\True $\dfrac{5}{3}$} 
	\loigiai{ 
		Ta có: $\dfrac{\sin \alpha +2 \cos \alpha}{3\sin \alpha - \cos \alpha} = \dfrac{\dfrac{\sin \alpha}{\cos \alpha} + \dfrac{2 \cos \alpha}{\cos \alpha}}{\dfrac{3\sin \alpha}{\cos \alpha} - \dfrac{\cos \alpha}{\cos \alpha}}=\dfrac{\tan \alpha +2}{3\tan \alpha -1}=\dfrac{2+2}{3 \cdot 2 -1}=\dfrac{4}{5}$. 
	} 
\end{ex}

% nguồn: Ok_TracNghiemLop11_Lan1
\begin{ex}%[1D1H2-2]%[ma:1a34bd03]
	Cho $\tan x=\dfrac{3}{4}$ và  góc $x$ thỏa mãn $\pi<x<\dfrac{3\pi}{2}$. Giá trị $\cos x$ bằng
	\choice 
		{$\dfrac{3}{5}$} 
		{$\dfrac{4}{5}$} 
		 {$-\dfrac{3}{5}$}
		{\True  $-\dfrac{4}{5}$}
	 \loigiai{
		Ta có $\pi<x<\dfrac{3\pi}{2}\Rightarrow  \cos x<0$.\\
		Lại  có
		\begin{eqnarray*}
			&&1+\tan^2x=\dfrac{1}{\cos^{2}x}\\ 
			&\Rightarrow&  \cos^{2}x=\dfrac{1}{1+\tan^2x}\\
			&\Rightarrow& \cos  x=-\sqrt{\dfrac{1}{1+\tan^2x}}=-\sqrt{\dfrac{1}{1+\left(\dfrac{3}{4}\right)^{2}}}=-\dfrac{4}{5}.
		\end{eqnarray*}
	} 
\end{ex}

% nguồn: Ok_TracNghiemLop11_Lan1
\begin{ex}%[1D1H2-2]%[ma:47445273]
	Cho biết $\tan\alpha=\dfrac{1}{2}$. Tính  $\cot\alpha $
	 \choice
		{\True $\cot\alpha=2$} 
		{$\cot\alpha=\dfrac{1}{4}$} 
		{$\cot\alpha=\dfrac{1}{2}$} 
		 {$\cot\alpha=\sqrt{2}$}
	\loigiai{Ta có  $\tan\alpha \cdot \cot\alpha=1 \Rightarrow\cot\alpha=\dfrac{1}{\tan\alpha}=\dfrac{1}{\dfrac{1}{2}}=2$.}
\end{ex}

% nguồn: Ok_TracNghiemLop11_Lan1
\begin{ex}%[1D1H2-2]%[ma:520648a2]
	Cho góc $\alpha $ thỏa mãn $\sin \alpha =\dfrac{{12}}{{13}}$ và $\dfrac{\pi}{2}<\alpha <\pi $. Tính $\cos \alpha.$ 
	 \choice
		 {$\cos \alpha =\dfrac{1}{{13}}$}
		{$\cos  \alpha =\dfrac{5}{{13}}$}
		{$\cos \alpha  =-\dfrac{5}{{13}}$}
		{\True  $\cos \alpha =-\dfrac{1}{{13}}$}
	\loigiai{Ta có  $\heva{& \cos \alpha =\pm \sqrt{{1-{\sin}^2\alpha}}=\pm \dfrac{5}{{13}} \\
		& \dfrac{\pi}{2}<\alpha <\pi}\Rightarrow\cos  \alpha =-\dfrac{5}{{13}}.$
	 }
\end{ex}

% nguồn: Ok_TracNghiemLop11_Lan1
\begin{ex}%[1D1H2-2]%[ma:0fa3624c]
	Giá trị của biểu thức $ S=3-\sin^290^\circ+2\cos^260^\circ-3\tan^245^\circ $ bằng 
	 \choice
		{$\dfrac{1}{2}$} 
		{\True  $-\dfrac{1}{2}$}
		{$  1$}
		{$  3$}
	\loigiai{Ta có $  S=3-\sin^290^\circ+2\cos^260^\circ-3\tan^245^\circ =3-1^2+2\cdot \left(\dfrac{1}{2}\right)^2-3\cdot 1^2=-\dfrac{1}{2}$.}
\end{ex}

% nguồn: Ok_TracNghiemLop11_Lan1
\begin{ex}%[1D1H2-2]%[ma:45154d98]
	Tìm khẳng định {\bf sai} trong  các khẳng định sau đây?
	\choice 
		{$\tan 45^\circ <\tan 60^\circ  $}
		{$\cos 45^\circ\le\sin 45^\circ  $}
		{$\sin  60^\circ <\sin 80^\circ $}
		{\True $\cos 35^\circ  >\cos 10^\circ $}
	\loigiai{Khi $\alpha\in\left(0^\circ\,;\,90^\circ\right)$ hàm $\cos\alpha $ là hàm giảm nên $\cos  35^\circ <\cos 10^\circ $ suy ra $\cos 35^\circ >\cos 10^\circ $ sai.}
\end{ex}

% nguồn: Ok_TracNghiemLop11_Lan1
\begin{ex}%[1D1H2-2]%[ma:a92b80a2]
	Trong các đẳng thức sau,  đẳng thức nào đúng?
	 \choice
		{$\cos 150^\circ=\dfrac{\sqrt{3}}{2}$} 
		{$\cot 150^\circ=\sqrt{3}$} 
		{\True  $\tan 150^\circ=-\dfrac{1}{\sqrt{3}}$}
		 {$\sin 150^\circ=-\dfrac{\sqrt{3}}{2}$}
	\loigiai{ $\cos 150^\circ=-\dfrac{\sqrt{3}}{2}$ suy ra  $\cos 150^\circ=\dfrac{\sqrt{3}}{2}$ sai.\\
	$\cot 150^\circ=-\sqrt{3}$ suy ra $\cot 150^\circ=\sqrt{3}$  sai.\\
	$\tan 150^\circ=-\dfrac{1}{\sqrt{3}}$  suy ra $\tan 150^\circ=-\dfrac{1}{\sqrt{3}}$ đúng.\\
	$\sin 150^\circ=\dfrac{1}{2}$ suy ra $\sin 150^\circ=-\dfrac{\sqrt{3}}{2}$ sai.} 
\end{ex}

% nguồn: Ok_Toan11_TracNghiem_Lan2
\begin{ex}%[1D1V2-2]%[ma:aa172855]
	Cho $\cot \alpha=4\tan \alpha$ và $\alpha \in \left(\dfrac{\pi}{2};\pi  \right)$. Khi đó $\sin \alpha$ bằng
	\choice 
		{$-\dfrac{\sqrt{5}}{5}$} 
		 {$\dfrac{1}{2}$}
		{$\dfrac{2\sqrt{5}}{5}$} 
		 {\True $\dfrac{\sqrt{5}}{5}$}
	\loigiai{ 
		 Ta có
		 \allowdisplaybreaks
		\begin{eqnarray*}
			&\cot \alpha=4\tan \alpha &\Leftrightarrow  \dfrac{\cot \alpha}{\tan \alpha}=4\Leftrightarrow \cot ^2\alpha=4\Leftrightarrow 1+\cot ^2\alpha=5 \\
			&&\Leftrightarrow \dfrac{1}{\sin ^2\alpha}=5\Leftrightarrow  \sin ^2\alpha=\dfrac{1}{5} \Leftrightarrow \sin \alpha=\pm \dfrac{\sqrt{5}}{5}.
		\end{eqnarray*}
		Vì $\alpha \in \left(\dfrac{\pi}{2};\pi \right)$  nên $\sin \alpha=\dfrac{\sqrt{5}}{5}$.
	 }
\end{ex}

% nguồn: Ok_Toan11_TracNghiem_Lan2
\begin{ex}%[1D1V2-2]%[ma:f39e805c]
	Nếu $\sin x+\cos x=\dfrac{1}{2}$ thì $3 \sin x+2  \cos x$ bằng
	\choice 
		{\True  $\dfrac{5-\sqrt{7}}{4}$ hoặc $\dfrac{5+\sqrt{7}}{4}$}
		{$\dfrac{5-\sqrt{5}}{7}$ hoặc $\dfrac{5+\sqrt{5}}{4}$} 
		{  $\dfrac{2-\sqrt{3}}{5}$ hoặc $\dfrac{2+\sqrt{3}}{5}$}
		 {$\dfrac{3-\sqrt{2}}{5}$ hoặc $\dfrac{3+\sqrt{2}}{5}$}
	\loigiai{ 
		 Ta có
		\[(3 \sin x+2 \cos  x)=2(\sin x+\cos x)+\sin x=1+\sin x.\]
		 Từ $\sin x+\cos x=\dfrac{1}{2}\Rightarrow \sin x \cdot \cos x=-\dfrac{3}{8}$.\\
		 Khi đó $\sin x$, $\cos x$ là nghiệm của phương trình $t^2-\dfrac{1}{2}t-\dfrac{3}{8}=0$. \\
		 Mặt khác
		\[t^2-\dfrac{1}{2}t-\dfrac{3}{8}=0\Leftrightarrow \hoac{&t=\dfrac{1+\sqrt{7}}{4} \\&t=\dfrac{1-\sqrt{7}}{4}.}\] 
		\begin{itemize}
			\item Với $\sin x=\dfrac{1+\sqrt{7}}{4}$ suy ra $(3  \sin x+2 \cos x)=1+\sin x=1+\dfrac{1+\sqrt{7}}{4}=\dfrac{5+\sqrt{7}}{4}$.
			 \item Với $\sin x=\dfrac{1-\sqrt{7}}{4}$ suy ra $(3 \sin x+2 \cos x)=1+\sin x=1+\dfrac{1-\sqrt{7}}{4}=\dfrac{5-\sqrt{7}}{4}$.
		\end{itemize}
	 }
\end{ex}

% nguồn: Ok_TracNghiemLop11_Lan1
\begin{ex}%[1D1V2-2]%[ma:4609ab3c]
	Biểu thức $A=\sin^210^\circ+\sin^220^\circ+\cdots +\sin^2180^\circ  $ có giá trị bằng
	 \choice
		 {$A=6$}
		 {$A=8$}
		{$A=3$} 
		{\True  $A=9$}
	\loigiai{ Ta có $\sin\left(90^\circ+\alpha\right)=\cos\alpha  $.\\
	Suy ra  $\sin 100^\circ=\cos 10^\circ\Rightarrow{\sin^2}100^\circ=\cos^210^\circ $.\\
	Tương tự  ta có $\sin^2110^\circ=\cos^220^\circ $, $\sin^2120^\circ=\cos^230^\circ$, $\sin^2130^\circ=\cos^240^\circ $,\\
	 $\sin^2150^\circ=\cos^240^\circ $, $\sin^2160^\circ=\cos^270^\circ $, $\sin^2170^\circ=\cos^280^\circ $, $\sin^2180^\circ=\cos^290^\circ $.\\
	 Vậy ta có $A=\left(\sin^210^\circ+\cos^210^\circ\right)+\left(\sin^220^\circ+\cos^220^\circ\right)+ \cdots +\left(\sin^290^\circ+\cos^290^\circ\right)$ \\ $\Rightarrow A=1+1+\cdots +1=9$.}
\end{ex}

% nguồn: Ok_TracNghiemLop11_Lan1
\begin{ex}%[1D1V2-2]%[ma:26843690]
	Giá trị  của $\tan\left(\alpha+\dfrac{\pi}{3}\right)$ bằng bao nhiêu khi $\sin\alpha=\dfrac{3}{5}$ $\left(\dfrac{\pi}{2}<\alpha <\pi\right)$?
	\choice 
		{$\dfrac{-48+25\sqrt{3}}{11}$} 
		{$\dfrac{48+25\sqrt{3}}{11}$} 
		{$-\dfrac{48+25\sqrt{3}}{11}$} 
		{\True $\dfrac{48-25\sqrt{3}}{11}$} 
	\loigiai{ 
		 Ta có $\sin^2\alpha+\cos^2\alpha=1$
		$\Rightarrow\cos\alpha=\pm\sqrt{1-\sin^2\alpha}$ 
		 $\Rightarrow\cos\alpha=\pm\sqrt{1-\left(\dfrac{3}{5}\right)^2}=\pm\dfrac{4}{5}$.\\
		Mà  $\dfrac{\pi}{2}<\alpha <\pi$ nên $\cos\alpha=-\dfrac{4}{5}$. Suy ra $\tan\alpha=\dfrac{-3}{4}$.\\
	Khi đó ta có  $\tan\left(\alpha+\dfrac{\pi}{3}\right)=\dfrac{\tan\alpha+\tan\dfrac{\pi}{3}}{1-\tan\alpha \cdot \tan\dfrac{\pi}{3}}=\dfrac{\dfrac{-3}{4}+\sqrt{3}}{1-\left(\dfrac{-3}{4}\right)\cdot \sqrt{3}}=\dfrac{48-25\sqrt{3}}{11}$.}
\end{ex}

% nguồn: Ok_TracNghiemLop11_Lan1
\begin{ex}%[1D1V2-2]%[ma:fa36680a]
	Biết $\sin\alpha+\cos\alpha=\dfrac{\sqrt{3}}{2}$. Trong các kết  quả sau, kết quả nào {\bf sai}?
	 \choice
		{$\sin\alpha .\cos\alpha=-\dfrac{1}{8}$} 
		{$\sin\alpha-\cos\alpha=\pm\dfrac{\sqrt{5}}{2}$} 
		{\True  $\sin^4\alpha+\cos^4\alpha=\dfrac{5}{4}$}
		 {$\tan^2\alpha+\cot^2\alpha=62$}
	\loigiai{		$\begin{aligned}
	\sin\alpha+\cos\alpha=\dfrac{\sqrt{3}}{2}&\Rightarrow{\left(\sin\alpha+\cos\alpha\right)^2}=\left(\dfrac{\sqrt{3}}{2}\right)^2\\ 
	&\Leftrightarrow{\sin^2}\alpha+\cos^2\alpha+2\sin\alpha \cdot \cos\alpha=\dfrac{3}{4}\\ 
	 &\Leftrightarrow 1+2\sin\alpha \cdot \cos\alpha=\dfrac{3}{4}\\
	&\Leftrightarrow\sin\alpha \cdot  \cos\alpha=-\dfrac{1}{8}.
\end{aligned}$\\
$\left(\sin\alpha-\cos\alpha\right)^2=\sin^2\alpha+\cos^2\alpha-2\sin\alpha  \cdot \cos\alpha=1+\dfrac{1}{4}=\dfrac{5}{4}$
 $\Rightarrow\sin\alpha-\cos\alpha=\pm\dfrac{\sqrt{5}}{2}$.\\
 $\sin^4\alpha+\cos^4\alpha=\left(\sin^2\alpha+\cos^2\alpha\right)^2-2\sin^2\alpha \cdot \cos^2\alpha=1-2\cdot \left(-\dfrac{1}{8}\right)^2=\dfrac{31}{32}$.\\
 $\tan^2\alpha+\cot^2\alpha=\dfrac{\sin^2\alpha}{\cos^2\alpha}+\dfrac{\cos^2\alpha}{\sin^2\alpha}=\dfrac{\sin^4\alpha+\cos^4\alpha}{\sin^2\alpha \cdot \cos^2\alpha}=\dfrac{\dfrac{31}{32}}{\dfrac{1}{64}}=62$.
} 
\end{ex}

% nguồn: Ok_TracNghiemLop11_Lan1
\begin{ex}%[1D1B2-2]%[ma:9dc2d584]
	Cho $\dfrac{\pi}{2}<\alpha <\pi $. Kết quả đúng là 
	 \choice
		{$\sin\alpha  >0,$ $\cos\alpha >0$}
		{$\sin\alpha <0,$  $\cos\alpha <0$}
		{\True $\sin\alpha >0,$  $\cos\alpha <0$}
		{$\sin\alpha <0,$ $\cos\alpha  >0$}
	 \loigiai{ Vì $\dfrac{\pi}{2}<\alpha <\pi $ nên $\alpha $ nằm trong góc phần tư thứ $ 2$ nên $\sin\alpha >0,\cos\alpha <0$.}
\end{ex}


%%%%%%%%%%%%%%%% ĐÚNG/SAI %%%%%%%%%%%%%%%%

% nguồn: Ok_dungsai_lop11_Lan1
\begin{ex}%[1D1N2-2]%[ma:a2f213d4]
 \immini{Cho  đường tròn lượng giác gốc $A$ (tham khảo hình vẽ). Gọi $M$ là điểm biểu diễn góc lượng giác có số đo $\alpha$. Biết rằng $M$ không trùng với một trong các điểm $A, A', B', B$. Mỗi khẳng định sau đúng hay sai?
  \choiceTF
  {\True Nếu $M$  thuộc cung $AB$ thì $\sin\alpha >0$}
  {\True Nếu $M$ thuộc cung $BA'$ thì $\cos\alpha<0$} 
  {Nếu $M$ thuộc cung  $A'B'$ thì $\sin\alpha>0$}
  {\True Nếu $M$ thuộc cung $B'A$ thì $\cos\alpha>0$} 
  } 
  { \begin{tikzpicture}[line join = round,line cap = round,>=stealth,font=\footnotesize,scale=0.77]
  \tikzset{label  style/.style={font=\footnotesize}}
  \path (0,0) coordinate  (O)
   (3,0) coordinate  (A)
   (0,3) coordinate (B) 
   (0,-3)  coordinate (B')
   (-3,0) coordinate  (A')
   (0:0)++(-120:3)  coordinate (M)
   ($(O)!(M)!(A')$) coordinate (H) 
   ($(O)!(M)!(B)$) coordinate  (K);
  \draw[->] (-4,0) -- (4,0)  node[above,blue]{$x$};
  \draw[->] (0,-4) -- (0.,4)  node[left,blue]{$y$};
  \draw[orange] (O) circle  (3cm);
  \foreach \p/\r in  {A/-45,O/-45,A'/-135,B'/-45,B/45}
  \fill (\p) circle (1pt)  node[shift={(\r:3mm)},blue]{$\p$};
 \end{tikzpicture}
   }
  \loigiai{Ta có
  \begin{itemchoice}
   \itemch $M$  thuộc cung $AB$ thì $\sin\alpha>0$.
   \itemch  $M$ thuộc cung $BA'$ thì $\cos\alpha<0$.
   \itemch  $M$ thuộc cung $A'B'$ thì $\sin\alpha<0$.
   \itemch $M$ thuộc  cung $B'A$ thì $\cos\alpha>0$.
  \end{itemchoice}
  }
\end{ex}

% nguồn: Ok_dungsai_lop11_Lan1
\begin{ex}%[1D1H2-2]%[ma:751e7f8b]
 Biết $\sin a=\dfrac{8}{17}$, $\tan b=\dfrac{5}{12}$ và $a$, $b$ là các  góc nhọn. Khi đó.
 \choiceTF[2] 
  {\True $\tan  a=\dfrac{8}{15}$}
   {\True $\sin (a-b)=\dfrac{21}{221}$}
  {$\cos  (a+b)=\dfrac{14}{22}$}
  {$\tan 2(a+b)=\dfrac{171}{140}$} 
  \loigiai{
  \begin{itemchoice}
   \itemch Đúng. Vì $a$, $b$  là các góc nhọn nên $\cos a>0$, $\cos b>0$.\\
   Ta có $\cos a=\sqrt{1-\sin ^2 a}=\dfrac{15}{17} \Rightarrow \tan  a=\dfrac{\sin a}{\cos a}=\dfrac{8}{15}$.
   \itemch Đúng. Ta có $\cos b=\sqrt{\dfrac{1}{1+\tan ^2 b}}=\dfrac{12}{13} \Rightarrow \sin b=\cos  b \tan b=\dfrac{5}{13}$.\\
   Khi đó 
   $\sin (a-b)=\sin a \cos  b-\cos a \sin b=\dfrac{8}{17} \cdot \dfrac{12}{13}-\dfrac{15}{17} \cdot \dfrac{5}{13}=\dfrac{21}{221}$.
   \itemch Sai.  Ta có $\cos (a+b)=\cos a \cos b-\sin a \sin b=\dfrac{15}{17} \cdot \dfrac{12}{13}-\dfrac{8}{17} \cdot \dfrac{5}{13}=\dfrac{140}{221}$.
   \itemch Sai. Ta có $\tan (a+b)=\dfrac{\tan a+\tan  b}{1-\tan a \tan b}=\dfrac{\dfrac{8}{15}+\dfrac{5}{12}}{1-\dfrac{8}{15} \cdot \dfrac{5}{12}}=\dfrac{171}{140}$.\\
   Vậy $\tan 2(a+b)=\dfrac{2 \tan (a+b)}{1-\tan ^2(a+b)}=\dfrac{2  \cdot \dfrac{171}{140}}{1-\left(\dfrac{171}{140}\right)^2} \approx-5$.
  \end{itemchoice}
 } 
\end{ex}

% nguồn: Ok_dungsai_lop11_Lan1
\begin{ex}%[1D1H2-2]%[ma:2c3ee103]
 Cho góc lượng giác $\alpha$ thỏa mãn $\cos  \alpha=\dfrac{3}{5}$, $ 0<\alpha<\dfrac{\pi}{2}$. Xét tính đúng - sai của các phát biểu sau.
 \choiceTF 
  {Giá trị sin của  góc lượng giác $\alpha$ âm}
  { \True Giá trị tan của góc  lượng giác $\alpha$ dương}
  {Giá trị của biểu thức $\mathrm{P}=\sin ^2 \alpha-2 \cos  ^2 \alpha=\dfrac{2}{25}$}
  { \True Giá trị của biểu thức  $Q=\tan \alpha+\cot \alpha>2$}
 \loigiai{ 
  \begin{itemchoice}
   \itemch Sai. Vì $0<\alpha<\dfrac{\pi}{2}$ nên $\sin  \alpha>0$.
   \itemch Đúng. Vì $0<\alpha<\dfrac{\pi}{2}$ nên $\sin \alpha>0$,  $\cos \alpha>0 $ suy ra $ \tan \alpha>0$.
   \itemch Sai.  Vì
   $\heva{&\cos \alpha=\dfrac{3}{5}\\ &0<\alpha<\dfrac{\pi}{2}} \Rightarrow \sin \alpha=\sqrt{1-\cos ^2 \alpha}=\dfrac{4}{5} $.\\ 
   Vậy $P=\sin  ^2 \alpha-2 \cos ^2 \alpha=\dfrac{16}{25}-\dfrac{18}{25}=-\dfrac{2}{25}$.
   \itemch  Đúng. $Q=\tan \alpha+\cot \alpha=\dfrac{4}{3}+\dfrac{3}{4}=\dfrac{25}{12}>2$.
  \end{itemchoice}
  }
\end{ex}

% nguồn: Ok_dungsai_lop11_Lan1
\begin{ex}%[1D1H2-2]%[ma:ecb9d5ed]
 \immini{Cho đường tròn lượng giác gốc $A$ (tham khảo hình vẽ). Gọi $M$ là điểm biểu diễn góc lượng giác có số đo $\alpha$. Biết rằng $M$ không trùng với một trong các điểm  $A, A', B', B$, mỗi khẳng định sau đúng hay sai?
 \choiceTF 
  {\True Nếu  $M$ thuộc cung $BA'$ thì $\cot\alpha<0$}
  {Nếu $M$ thuộc cung  $AB$ thì $\tan\alpha >0$}
  {Nếu $M$ thuộc cung $A'B'$ và  $ \cos \alpha=-\dfrac{3}{5}$ thì $\sin \alpha =\dfrac{4}{5} $}
  {\True Nếu $M$ thuộc cung  $B'A$ và $\tan \alpha=-2$ thì $\cos \alpha=\dfrac{\sqrt{5}}{5}$}
  } 
  { \begin{tikzpicture}[line join = round,line cap = round,>=stealth,font=\footnotesize,scale=0.77]
   \tikzset{label style/.style={font=\footnotesize}}
  \path (0,0) coordinate  (O)
   (3,0) coordinate  (A)
   (0,3)  coordinate (B)
    (0,-3) coordinate (B')
   (-3,0) coordinate (A') 
    (0:0)++(-120:3) coordinate (M)
   ($(O)!(M)!(A')$) coordinate  (H)
    ($(O)!(M)!(B)$) coordinate (K);
  \draw[->] (-4,0) -- (4,0) node[above,blue]{$x$}; 
   \draw[->] (0,-4) -- (0.,4) node[left,blue]{$y$};
  \draw[orange] (O)  circle (3cm);
  \foreach  \p/\r in {A/-45,O/-45,A'/-135,B'/-45,B/45}
  \fill (\p) circle  (1pt) node[shift={(\r:3mm)},blue]{$\p$};
 \end{tikzpicture}
   }
 \loigiai{Ta  có
  \begin{itemchoice}
   \itemch $M$ thuộc cung $BA'$ thì  $\cot\alpha<0$.
   \itemch $M$ thuộc  cung $AB$ thì $\tan\alpha>0$.
   \itemch Ta có  $\sin^2\alpha =1-\cos^2\alpha =\dfrac{16}{25}$.\\
   Vì $M$ thuộc cung $A'B'$ nên $\sin\alpha<0 $, suy  ra $ \sin\alpha=-\dfrac{4}{5}$.
   \itemch Ta  có $\cos^2\alpha =\dfrac{1}{1+\tan^2\alpha}=\dfrac{1}{5}$.\\
    Vì $M$ thuộc cung $B'A$ nên $\cos\alpha>0$, suy ra $\cos \alpha=\dfrac{\sqrt{5}}{5}$.
  \end{itemchoice}
  }
\end{ex}

% nguồn: Ok_dungsai_lop11_Lan1
\begin{ex}%[1D1H2-2]%[1D1H2-3]%[ma:a617bd2a]
 Cho $\tan \alpha=2$ và $a, b$ là 2 số thực dương. Xét tính đúng - sai của các phát  biểu sau.
 \choiceTF 
   {\True $\cot \alpha=\dfrac{1}{2}$}
  {Khi  $0 \leq \alpha \leq \dfrac{\pi}{2}$ thì $\cos \alpha=-\dfrac{\sqrt{5}}{5}$}
  {\True Giá trị  của biểu thức $B=\dfrac{\sin \alpha-\cos \alpha}{\sin ^3 \alpha+3 \cos ^3 \alpha+2 \sin \alpha}=\dfrac{5}{31}$}
  {Cho $\dfrac{\sin ^4 \alpha}{a}+\dfrac{\cos ^4 \alpha}{b}=\dfrac{1}{a+b}$. Giá trị của biểu thức $A=\dfrac{\sin ^8 \alpha}{a^3}+\dfrac{\cos ^8  \alpha}{b^3}=\dfrac{1}{a^3+b^3}$}
 \loigiai{ 
  \begin{itemchoice}
   \itemch Đúng. Ta có $\tan \alpha \cdot \cot \alpha=1 \Rightarrow \cot \alpha=\dfrac{1}{\tan  \alpha}=\dfrac{1}{2}$.
   \itemch  Sai. Ta có $1+\tan ^2 \alpha=\dfrac{1}{\cos ^2 \alpha} \Rightarrow \cos ^2 \alpha=\dfrac{1}{1+\tan ^2 \alpha}=\dfrac{1}{5} \Rightarrow \cos \alpha= \pm \dfrac{\sqrt{5}}{5}$.\\
   Mà $0 \leq \alpha  \leq \dfrac{\pi}{2}$ nên $\cos \alpha=\dfrac{\sqrt{5}}{5}$.
   \itemch Đúng. $B=\dfrac{\dfrac{\sin \alpha}{\cos ^3 \alpha}-\dfrac{\cos \alpha}{\cos ^3 \alpha}}{\dfrac{\sin ^3  \alpha}{\cos ^3 \alpha}+\dfrac{3 \cos ^3 \alpha}{\cos ^3 \alpha}+\dfrac{2 \sin \alpha}{\cos ^3 \alpha}}=\dfrac{\tan \alpha\left(\tan ^2 \alpha+1\right)-\left(\tan ^2 \alpha+1\right)}{\tan ^3 \alpha+3+2 \tan \alpha\left(\tan ^2 \alpha+1\right)}$.\\
   Suy ra  $B=\dfrac{2(4+1)-(4+1)}{8+3+2\cdot 2(4+1)}=\dfrac{5}{31}$.
    \itemch Sai. Đặt $\cos ^2 \alpha=t$. Khi đó\allowdisplaybreaks
   $ 
   \begin{aligned}
    &\, \dfrac{(1-t)^2}{a}+\dfrac{t^2}{b}=\dfrac{1}{a+b}  \\
    \Leftrightarrow&\, b(1-t)^2+a  t^2=\dfrac{a b}{a+b} \\
     \Leftrightarrow&\, a t^2+b t^2-2 b t+b=\dfrac{a b}{a+b} \\
    \Leftrightarrow&\, (a+b) t^2-2 b t+b=\dfrac{a  b}{a+b} \\
     \Leftrightarrow&\,(a+b)^2 t^2-2 b(a+b) t+b^2=0\\
    \Leftrightarrow&\,  t=\dfrac{b}{a+b}.
   \end{aligned}
   $\\ 
   Suy ra $\heva{&\cos ^2 \alpha=\dfrac{b}{a+b} \\  &\sin ^2 \alpha=\dfrac{a}{a+b}.}$\\
   Vậy $\dfrac{\sin ^8 \alpha}{a^3}+\dfrac{\cos  ^8 \alpha}{b^3}=\dfrac{a}{(a+b)^4}+\dfrac{b}{(a+b)^4}=\dfrac{1}{(a+b)^3}$.
  \end{itemchoice}
 } 
\end{ex}

% nguồn: Ok_dungsai_lop11_Lan1
\begin{ex}%[1D1H2-2]%[ma:3c7e6e19]
 Cho $\tan a=\dfrac{1}{2}$, $0<a<\dfrac{\pi}{2}$, $ \tan b=\dfrac{1}{3}$, $ 0<b<\dfrac{\pi}{2}$. Xét tính đúng  - sai của các phát biểu~sau.
 \choiceTF[2] 
  {  \True$\cot a=2$}
   { \True$\tan (a+b)=1$}
  {  $\tan (a-b)=7$}
  {  \True $\sin (a-2 b)=-\dfrac{2 \sqrt{5}}{25}$}
  \loigiai{
  \begin{itemchoice}
   \itemch Đúng. Ta có $\tan  a \cdot \cot a=1 \Rightarrow \cot a=\dfrac{1}{\tan a}=\dfrac{1}{\dfrac{1}{2}}=2$.
   \itemch  Đúng. Ta có $\tan (a+b)=\dfrac{\tan a+\tan b}{1-\tan a \tan b}=\dfrac{\dfrac{1}{2}+\dfrac{1}{3}}{1-\dfrac{1}{2} \cdot \dfrac{1}{3}}=1$.
   \itemch Sai. Ta có $\tan  (a-b)=\dfrac{\tan a-\tan b}{1+\tan a \tan b}=\dfrac{\dfrac{1}{2}-\dfrac{1}{3}}{1+\dfrac{1}{2} \cdot \dfrac{1}{3}}=\dfrac{1}{7}$.
   \itemch Đúng. Ta có 
   \begin{itemize}
     \item $
    \begin{aligned}
     &\,\cos ^2 a=\dfrac{1}{1+\tan ^2  a}=\dfrac{1}{1+\left(\dfrac{1}{2}\right)^2}=\dfrac{4}{5}\\
     \Rightarrow&\,\hoac{& 
     \cos  a=\dfrac{2 \sqrt{5}}{5}\text{ (nhận)} \\
     &\cos a=-\dfrac{2 \sqrt{5}}{5}\text{  (loại)}
     } \Rightarrow \sin a=\tan  a \cdot \cos a=\dfrac{\sqrt{5}}{5}.
    \end{aligned}
     $
    \item  $
    \begin{aligned}
     &\,\cos ^2 b=\dfrac{1}{1+\tan ^2 b}=\dfrac{1}{1+\left(\dfrac{1}{3}\right)^2}=\dfrac{9}{10} \\ 
     \Rightarrow&\,\hoac{& 
     \cos  b=\dfrac{3 \sqrt{10}}{10}\text{ (nhận)} \\
     & \cos  b=-\dfrac{3 \sqrt{10}}{10}\text{ (loại)}
     } \Rightarrow \sin b=\tan  b \cdot \cos b=\dfrac{\sqrt{10}}{10}.
    \end{aligned}
    $ 
   \end{itemize}
   Suy ra $$\sin (a-2 b)=\sin a \cdot\cos 2 b-\cos a\cdot \sin 2 b=\sin a\cdot \left(1-2 \sin ^2 b\right)-\cos a \cdot (2 \sin b \cos b)=-\dfrac{2  \sqrt{5}}{25}.$$
  \end{itemchoice}
  }
\end{ex}

% nguồn: Ok_dungsai_lop11_Lan1
\begin{ex}%[1D1H2-2]%[ma:c34d37ac]
 Cho $0^{\circ}<\alpha<90^{\circ}$. Các khẳng định sau đây  đúng hay sai?
  \choiceTF
   {\True $ A=\sin\left(\alpha+90^{^\circ}\right)>0$}
   {\True $ B=\cos\left(\alpha-45^{^\circ}\right)>0$}
   {$ C=\tan\left(270^{^\circ}-\alpha\right)<0$}
  {$ D=\cos\left(2\alpha+90^{^\circ}\right)>0$} 
 \loigiai{ 
  \begin{itemchoice}
   \itemch Ta có $0^{\circ}<\alpha<90^{\circ}\Rightarrow  90^{\circ}<\alpha+90^{\circ}<180^{\circ}$
    $\Rightarrow\sin\left(\alpha+90^{\circ}\right)>0$.
   \itemch Ta có  $0^{\circ}<\alpha<90^{\circ}\Rightarrow-45^{\circ}<\alpha-45^{\circ}<45^{\circ}$
    $\Rightarrow\cos\left(\alpha-45^{\circ}\right)>0$.
   \itemch Ta  có $0^{\circ}<\alpha<90^{\circ}\Rightarrow-90^{\circ}<-\alpha<0^{\circ}$\\
   $\begin{aligned}
    &\Rightarrow 270^{\circ}+\left(-90^{\circ}\right)<270^{\circ}+(-\alpha)<270^{\circ}+0^{\circ}\\
   &\Rightarrow 180^{\circ}<270^{\circ}-\alpha<270^{\circ}\\ 
   &\Rightarrow\tan\left(270^{\circ}-\alpha\right)>0. 
  \end{aligned}$
  \itemch Ta  có $0^{\circ}<\alpha<90^{\circ}\Rightarrow 90^{\circ}<2\alpha+90^{\circ}<270^{\circ}$
  $\Rightarrow\cos\left(2\alpha+270^{\circ}\right)<0$. 
 \end{itemchoice}
 }
\end{ex}

% nguồn: Ok_dungsai_lop11_Lan1
\begin{ex}%[1D1H2-2]%[ma:62e22ac5]
 Biết  $\tan\alpha=\dfrac{2\sqrt{10}}{9},\pi<\alpha<\dfrac{3\pi}{2}$. Các khẳng định sau đây đúng hay sai?
 \choiceTF 
  {\True $\cot\alpha=\dfrac{9}{2\sqrt{10}}$} 
  {\True $\cos\alpha=-\dfrac{9}{11}$} 
  {\True $\left\{\begin{array}{*{35}{l}}
  \cos\alpha <0\\ 
  \sin\alpha <0\\ 
\end{array}\right.$}
{$\sin\alpha=-\dfrac{2\sqrt{10}}{11}$} 
 \loigiai{
 $\pi<\alpha<\dfrac{3\pi}{2}\Rightarrow\heva{\cos\alpha<0\\\sin\alpha<0}$;  $\cot\alpha=\dfrac{9}{2\sqrt{10}}$; \\
 $\cos\alpha=-\sqrt{\dfrac{1}{\tan ^2\alpha+1}}=-\sqrt{\dfrac{1}{\dfrac{40}{81}+1}}=-\dfrac{9}{11}$;\\ 
  $\sin\alpha=-\sqrt{1-\cos ^2\alpha}=-\sqrt{1-\dfrac{81}{121}}=\dfrac{2\sqrt{10}}{11}$.
} 
\end{ex}

% nguồn: Ok_dungsai_lop11_Lan1
\begin{ex}%[1D1H2-2]%[1D1V2-3]%[ma:64cd628c]
 Cho $\sin \alpha=\dfrac{2}{3}$, $\dfrac{\pi}{2}<\alpha<\pi$, xét tính đúng  - sai của các mệnh đề sau.
 \choiceTF 
  {\True $\cos  \alpha<0$}
  { \True $\cos  \alpha=\dfrac{-\sqrt{5}}{3}$}
   {$\sin \left(\dfrac{5 \pi}{2}-\alpha\right)=\dfrac{\sqrt{5}}{3}$}
  { Giá trị của biểu thức $P=\dfrac{\sqrt{5} \tan \alpha}{1+\tan  ^2 \alpha}$ là $\dfrac{10}{9}$}
  \loigiai{
  \begin{itemchoice}
   \itemch Đúng. Vì $\alpha$ nằm ở góc phần tư thứ  hai của đường tròn lượng giác.
   \itemch Đúng. Vì $\sin ^2 \alpha+\cos ^2 \alpha=1$ nên $\cos \alpha= \pm \sqrt{1-\sin ^2  \alpha}=\sqrt{1-\dfrac{4}{9}}=\dfrac{ \pm \sqrt{5}}{3}$ mà $\dfrac{\pi}{2}<\alpha<\pi$ nên $\cos \alpha<0$, suy ra $\cos \alpha=\dfrac{-\sqrt{5}}{3}$.
   \itemch Sai. Vì $\sin \left(\dfrac{5 \pi}{2}-\alpha\right)=\sin \left(\dfrac{\pi}{2}-\alpha+2 \pi\right)=\sin \left(\dfrac{\pi}{2}-\alpha\right)=\cos \alpha$  mà $\dfrac{\pi}{2}<\alpha<\pi$ nên $\cos \alpha<0$ mà $\dfrac{\sqrt{5}}{3}>0$.
   \itemch Sai. Ta có $\sin ^2 \alpha+\cos  ^2 \alpha=1$ nên $\cos \alpha= \pm \sqrt{1-\sin ^2 \alpha}=\sqrt{1-\dfrac{4}{9}}=\dfrac{ \pm \sqrt{5}}{3}$ mà $\dfrac{\pi}{2}<\alpha<\pi$ nên $\cos \alpha<0$.\\
    Suy ra $\cos \alpha=\dfrac{-\sqrt{5}}{3}$.\\
   Do đó $\tan \alpha=\dfrac{\sin \alpha}{\cos \alpha}=\dfrac{2}{3} \cdot\left(\dfrac{3}{-\sqrt{5}}\right)=\dfrac{-2}{\sqrt{5}}$.\\ 
   Thay vào biểu thức ta  có $P=\dfrac{\sqrt{5} \cdot\left(\dfrac{-2}{\sqrt{5}}\right)}{1+\dfrac{4}{5}}=\dfrac{-10}{9}$.
  \end{itemchoice}
 } 
\end{ex}

% nguồn: Ok_dungsai_lop11_Lan1
\begin{ex}%[1D1H2-2]%[ma:22fb8949]
 Cho $\sin \alpha =  \dfrac{12}{13}$ và $\dfrac{\pi}{2}<\alpha<\pi$. Khi đó
 \choiceTF 
  {\True  $\cos\alpha = -\dfrac{5}{13}$}
  {$\cos\alpha = \pm  \dfrac{5}{13}$}
   {\True $\tan\alpha = -\dfrac{12}{5}$}
  {$\cot\alpha  = \dfrac{5}{13}$}
  \loigiai{
  Với $\sin  \alpha = \dfrac{12}{13}$ và $\dfrac{\pi}{2}<\alpha<\pi$.\\
  Ta  có \\
  $\sin^2\alpha+\cos^2\alpha=1 \Leftrightarrow \cos\alpha = \pm\dfrac{5}{13} \Rightarrow \cos\alpha = -\dfrac{5}{13}\  \left(\dfrac{\pi}{2}<\alpha<\pi\right)$.\\
   $\Rightarrow \tan\alpha = \dfrac{\sin\alpha}{\cos\alpha} = -\dfrac{12}{5}$.\\
   $\Rightarrow \cot\alpha=\dfrac{1}{\tan\alpha}=-\dfrac{12}{5}$.
  }
\end{ex}

% nguồn: Ok_dungsai_lop11_Lan1
\begin{ex}%[1D1H2-2]%[ma:393b7390]
 Cho $\alpha=-\dfrac{\pi}{4}+(2k+1)\pi  $ (biết $k\in\mathbb{Z}$). Các khẳng định sau đây đúng hay sai?
  \choiceTF
   {\True $\sin\alpha=\dfrac{\sqrt{2}}{2}$}
  {\True $\cos\alpha=\dfrac{\sqrt{2}}{2}$} 
   {$\tan\alpha=-1$}
   {$\cot\alpha=-1$}
  \loigiai{
  Ta có  \\
  $\begin{aligned}
  \sin\left[-\dfrac{\pi}{4}+(2 k+1)\pi\right]& =\sin\left(-\dfrac{\pi}{4}+2 k\pi+\pi\right)=\sin\left(-\dfrac{\pi}{4}+\pi\right)  \\
   &=-\sin\left(-\dfrac{\pi}{4}\right)=\sin\dfrac{\pi}{4}=\dfrac{\sqrt{2}}{2}.
 \end{aligned}$\\
 $\begin{aligned}
  \cos\left[-\dfrac{\pi}{4}+(2k+1)\pi\right]& =\cos\left(-\dfrac{\pi}{4}+2k\pi+\pi\right)=\cos\left(-\dfrac{\pi}{4}+\pi\right)\\
  &=-\cos\left(-\dfrac{\pi}{4}\right)=-\cos\dfrac{\pi}{4}=-\dfrac{\sqrt{2}}{2}.
 \end{aligned}$\\
 $\begin{aligned}
 &\tan\left[-\dfrac{\pi}{4}+(2k+1)\pi\right]=\tan\left(-\dfrac{\pi}{4}\right)=-\tan\dfrac{\pi}{4}=-1.\\ 
  &\cot\left[-\dfrac{\pi}{4}+(2k+1)\pi\right]=\cot\left(-\dfrac{\pi}{4}\right)=-\cot\dfrac{\pi}{4}=-1.
 \end{aligned}$
} 
\end{ex}

% nguồn: Ok_dungsai_lop11_Lan1
\begin{ex}%[1D1H2-2]%[ma:82c95558]
  Cho góc $\alpha$ thỏa mãn $\cos \alpha=\dfrac{3}{5}$ và $-\dfrac{\pi}{2}<\alpha<0$. Xét tính đúng - sai của các phát biểu~sau.
 \choiceTF 
   {$\sin \alpha>0$}
  {\True $\tan  \alpha<0$}
  {\True $\sin  \alpha=-\dfrac{4}{5}$}
  {$\sin  \left(\dfrac{\pi}{2}-\alpha\right)-\sin (-\alpha)=\dfrac{7}{5}$}
 \loigiai{ 
  \begin{itemchoice}
   \itemch Sai. Vì $-\dfrac{\pi}{2}<\alpha<0$ nên $\sin  \alpha<0$. Do đó $\sin \alpha>0$.
   \itemch Đúng. Vì $-\dfrac{\pi}{2}<\alpha<0$ nên $\tan \alpha<0$. Do đó $\tan \alpha<0$. 
   \itemch Đúng. Vì $-\dfrac{\pi}{2}<\alpha<0$ nên  $\sin \alpha=-\sqrt{1-\cos ^2 \alpha}=-\dfrac{4}{5}$. Do đó $\sin \alpha=-\dfrac{4}{5}$.
   \itemch Sai. Vì $\sin \left(\dfrac{\pi}{2}-\alpha\right)-\sin (-\alpha)=\cos  \alpha+\sin \alpha=-\dfrac{1}{5}$.\\
  \end{itemchoice}
 } 
\end{ex}

% nguồn: Ok_dungsai_lop11_Lan1
\begin{ex}%[1D1H2-2]%[ma:27068552]
 Cho $\alpha=\dfrac{\pi}{3}+k2\pi $ (biết  $k\in\mathbb{Z}$). Các khẳng định sau đây đúng hay sai?
 \choiceTF 
  {$\sin\alpha=-\dfrac{\sqrt{3}}{2}$} 
  {$\cos\alpha=-\dfrac{1}{2}$} 
  {\True  $\tan\alpha=\sqrt{3}$}
   {$\cot\alpha=-\dfrac{\sqrt{3}}{3}$}
  \loigiai{
  \begin{itemchoice}
    \itemch $ \sin\left(\dfrac{\pi}{3}+k 2\pi\right)=\sin\dfrac{\pi}{3}=\dfrac{\sqrt{3}}{2}$.
   \itemch $  \cos\left(\dfrac{\pi}{3}+k 2\pi\right)=\cos\dfrac{\pi}{3}=\dfrac{1}{2}$.
   \itemch  $ \tan\left(\dfrac{\pi}{3}+k 2\pi\right)=\tan\dfrac{\pi}{3}=\sqrt{3}$.
   \itemch  $ \cot\left(\dfrac{\pi}{3}+k 2\pi\right)=\cot\dfrac{\pi}{3}=\dfrac{\sqrt{3}}{3}$
  \end{itemchoice}
  }
\end{ex}

% nguồn: Ok_dungsai_lop11_Lan1
\begin{ex}%[1D1H2-2]%[1D1H2-3]%[ma:b80dc480]
 Cho góc  $\alpha$, $\left(0^\circ<\alpha<180^\circ\right)$ thỏa mãn $\tan \alpha=3$.
 \choiceTF 
  {$\cot \alpha=\dfrac{1}{\sqrt{3}}$} 
   { \True $\cos \alpha>0$}
  {\True $\sin  \alpha=\dfrac{3 \sqrt{10}}{10}$}
  {$\mathrm{P}=\dfrac{2 \sin \alpha-3 \cos \alpha}{3 \sin \alpha+2  \cos \alpha}=\dfrac{-3}{11}$}
 \loigiai{ 
  \begin{itemchoice}
   \itemch Sai. $\tan \alpha \cdot \cot  \alpha=1 \Rightarrow \cot \alpha=\dfrac{1}{3}$.
   \itemch Đúng. $0^\circ<\alpha<180^\circ \Rightarrow  \sin \alpha>0$, mà $\tan \alpha>0 \Rightarrow \cos \alpha>0$.
   \itemch Đúng. Ta có  $0^\circ<\alpha<180^\circ \Rightarrow \sin \alpha>0$.
   $$ 
   1+\cot ^2 \alpha=\dfrac{1}{\sin ^2 \alpha} \Leftrightarrow \dfrac{1}{\sin ^2 \alpha}=1+\dfrac{1}{9}=\dfrac{10}{9} \Leftrightarrow \sin ^2 \alpha=\dfrac{9}{10} \Rightarrow \sin  \alpha=\dfrac{3 \sqrt{10}}{10}.
   $$ 
   \itemch Sai.  Vì $\tan \alpha=3$ nên $\cos \alpha \neq 0$, chia cả tử và mẫu cho $\cos \alpha$ ta được
   $$P=\dfrac{\dfrac{2 \sin \alpha-3 \cos \alpha}{\cos \alpha}}{\dfrac{3 \sin \alpha+2 \cos \alpha}{\cos \alpha}}=\dfrac{2\cdot \dfrac{\sin \alpha}{\cos \alpha}-3}{3\cdot \dfrac{\sin \alpha}{\cos \alpha}+2}=\dfrac{2 \tan \alpha-3}{3 \tan  \alpha+2}=\dfrac{2\cdot 3-3}{3\cdot 3+2}=\dfrac{3}{11}.$$
  \end{itemchoice}
 } 
\end{ex}

% nguồn: Ok_dungsai_lop11_Lan1
\begin{ex}%[1D1H2-2]%[ma:474b9732]
 Cho $\cos a=\dfrac{3}{4}$, $ \sin a>0 $, $\sin b=\dfrac{3}{5}$, $  \cos b<0$. Xét tính đúng - sai của các phát biểu sau.
 \choiceTF 
  {\True Giá  trị của $\tan a=\dfrac{\sqrt{7}}{3}$}
  {Giá trị của $\cot  b=-\dfrac{2}{3}$}
  {Giá trị của $\cos 2  a+\cos 2 b$ thuộc khoảng $\left(\dfrac{1}{2}; 1\right)$}
  {Giá trị  của $\cos (a+b)$ thuộc khoảng $\left(-\dfrac{1}{2};-\dfrac{1}{3}\right)$}
 \loigiai{ 
  \begin{itemchoice}
   \itemch Đúng. Ta có $\tan ^2 a+1=\dfrac{1}{\cos ^2 a} \Rightarrow \tan a=\sqrt{\dfrac{1}{\cos  ^2 a}-1}=\sqrt{\dfrac{16}{9}-1}=\dfrac{\sqrt{7}}{3}$.
   \itemch Sai. Ta có $\cot ^2 b+1=\dfrac{1}{\sin ^2 b} \Rightarrow \cot b=-\sqrt{\dfrac{1}{\sin ^2 a}-1}=-\sqrt{\dfrac{25}{9}-1}=-\dfrac{4}{3}$. 
   \itemch  Sai. Ta có $
   \begin{aligned}
    \cos 2  a+\cos 2 b
    &\,=2 \cos ^2 a-1+1-2 \sin ^2 b\\ 
    &\,=2 \cdot \dfrac{9}{16}-2 \cdot  \dfrac{9}{25}=\dfrac{81}{200}\\
    &\,=0{,}405  \notin\left(\dfrac{1}{2}; 1\right).
   \end{aligned}
    $
   \itemch Sai. Ta  có $\heva{&\cos a=\dfrac{3}{4} \\ &\sin a>0} \Rightarrow \sin a=\sqrt{1-\cos ^2 a}=\dfrac{\sqrt{7}}{4}$.\\
   $\heva{&\sin b=\dfrac{3}{5}  \\ &\cos b<0} \Rightarrow \cos b=-\sqrt{1-\sin ^2 b}=-\dfrac{4}{5}$.\\
    Suy ra $
   \begin{aligned}
    \cos  (a+b)&\, =\cos a \cos b-\sin a \sin b\\
    &\, =\dfrac{3}{4} \cdot\left(-\dfrac{4}{5}\right)-\dfrac{\sqrt{7}}{4} \cdot \dfrac{3}{5}\\ 
     &\, =-\dfrac{3}{5}\left(1+\dfrac{\sqrt{7}}{4}\right) \notin\left(-\dfrac{1}{2};-\dfrac{1}{3}\right).
   \end{aligned}
    $
  \end{itemchoice}
  }
\end{ex}

% nguồn: Ok_dungsai_lop11_Lan1
\begin{ex}%[1D1H2-2]%[1D1H2-3]%[ma:cfc055d5]
 Cho góc $\alpha$ thỏa mãn $\dfrac{-\pi}{2}<\alpha<0$ và $\cot \alpha=-3$. Xét tính đúng - sai của các phát biểu~sau. 
 \choiceTF 
   {$\sin \alpha>0$}
  {\True $\sin \alpha=-\dfrac{\sqrt{10}}{10}$} 
  {\True $\dfrac{\cos \alpha-\sin  \alpha}{\cos ^3 \alpha+3 \sin ^3 \alpha+2 \cos \alpha}=\dfrac{10}{21}$}
  {\True $\left[\tan \dfrac{17 \pi}{4}+\tan \left(\dfrac{7  \pi}{2}-\alpha\right)\right]^2+\left[\cot \dfrac{13 \pi}{4}+\cot (7 \pi-\alpha)\right]^2=20$}
  \loigiai{
  \begin{itemchoice}
   \itemch Sai. Do $\dfrac{-\pi}{2}<\alpha<0$ nên  $\sin \alpha<0$.
   \itemch Đúng. $\sin \alpha=-\dfrac{\sqrt{10}}{10}$ 
   Do $\dfrac{-\pi}{2}<\alpha<0$ nên  $\sin \alpha<0$.\\
   Từ hệ thức $1+\cot ^2 \alpha=\dfrac{1}{\sin ^2 \alpha}$, suy ra $\sin \alpha=-\sqrt{\dfrac{1}{1+\cot ^2 \alpha}}=-\dfrac{\sqrt{10}}{10}$.\\ 
   Vì  vậy mệnh đề $\sin \alpha=-\dfrac{\sqrt{10}}{10}$.
    \itemch Đúng.
   Xét biểu thức  $P=\dfrac{\cos \alpha-\sin \alpha}{\cos ^3 \alpha+3 \sin ^3 \alpha+2 \cos \alpha}$.\\
   Chia cả  tử và mẫu cho $\sin ^3 \alpha$ ta được
   $$P=\dfrac{\dfrac{\cos \alpha}{\sin ^3 \alpha}-\dfrac{1}{\sin ^2 \alpha}}{\dfrac{\cos  ^3 \alpha}{\sin ^3 \alpha}+3+2\cdot \dfrac{\cos \alpha}{\sin ^3 \alpha}}=\dfrac{\cot \alpha \cdot\left(1+\cot ^2 \alpha\right)-\left(1+\cot ^2 \alpha\right)}{\cot ^3 \alpha+3+2 \cdot \cot \alpha \cdot\left(1+\cot ^2 \alpha\right)}.$$
   Thay $\cot \alpha=-3$ vào $P$ ta được $P=\dfrac{10}{21}$. 
   \itemch Đúng. 
   Ta có $\tan \dfrac{17 \pi}{4}=\tan  \left(\dfrac{\pi}{4}+4 \pi\right)=\tan \dfrac{\pi}{4}=1$ và $\tan \left(\dfrac{7 \pi}{2}-\alpha\right)=\cot \alpha$.\\
   Và $\heva{&\cot \dfrac{13  \pi}{4}=\cot \left(\dfrac{\pi}{4}+3 \pi\right)=\cot \dfrac{\pi}{4}=1\\ & \cot (7 \pi-\alpha)=-\cot \alpha.}$\\
   Suy ra $ 
   \begin{aligned}
    &\, \left[\tan \dfrac{17 \pi}{4}+\tan \left(\dfrac{7 \pi}{2}-\alpha\right)\right]^2+\left[\cot \dfrac{13 \pi}{4}+\cot  (7 \pi-\alpha)\right]^2\\
    =&\,(1+\cot x)^2+(1-\cot x)^2=2+2 \cot ^2 x=2+2 \cdot(-3)^2=20  .
   \end{aligned}
    $
  \end{itemchoice}
 } 
\end{ex}

% nguồn: Ok_dungsai_lop11_Lan1
\begin{ex}%[1D1H2-2]%[ma:aa01e905]
 Cho  $\pi<\alpha<\dfrac{3 \pi}{2}$.
 \choiceTF 
  {\True $\sin \alpha  <0$}
  {\True $\cos  (\alpha+\pi)>0$}
  {$\sin \left(\dfrac{\pi}{2}-\alpha\right)>0$} 
  {\True  $\tan \left(\alpha+\dfrac{3 \pi}{2}\right)<0$}
  \loigiai{
  Vì $\pi<\alpha<\dfrac{3 \pi}{2}$ nên $\alpha$ thuộc góc phần tư thứ III  nên $\sin\alpha<0$, $\cos\alpha<0$, $\tan\alpha>0$, $\cot\alpha>0$.
  \begin{itemchoice}
   \itemch  Đúng
   \itemch Đúng  vì $\cos (\alpha+\pi)=-\cos \alpha>0$.
   \itemch Sai vì $\sin  \left(\dfrac{\pi}{2}-\alpha\right)=\cos \alpha<0$.
   \itemch Đúng vì $\tan \left(\alpha+\dfrac{3 \pi}{2}\right)=-\cot \alpha<0$. 
  \end{itemchoice}
 } 
\end{ex}

% nguồn: Ok_dungsai_lop11_Lan1
\begin{ex}%[1D1H2-2]%[ma:0cc8ff83]
 Biết $\sin x=-\dfrac{3}{5}$ với $\pi<x<\dfrac{3\pi}{2}$. Các khẳng định  sau đây đúng hay sai?
  \choiceTF
   {$\cos x>0$}
  {\True  $\cos x=-\dfrac{4}{5}$}
   {\True $\tan x=\dfrac{3}{4}$}
  {\True  $\cot x=\dfrac{4}{3}$}
 \loigiai{ 
  Do  $\pi<x<\dfrac{3\pi}{2}$ nên $\cos x<0$.\\
  Ta có $\sin  ^2 x+\cos ^2 x=1\Rightarrow\cos ^2 x=1-\sin ^2 x=1-\dfrac{9}{25}=\dfrac{16}{25}$.\\
 $\Rightarrow\cos x=-\dfrac{4}{5};\tan x=\dfrac{\sin x}{\cos  x}=\dfrac{3}{4};\cot x=\dfrac{\cos x}{\sin x}=\dfrac{4}{3}$.}
\end{ex}

% nguồn: Ok_dungsai_lop11_Lan1
\begin{ex}%[1D1H2-2]%[ma:1214f765]
  Cho góc $ \alpha \in \left(\dfrac{3\pi}{2};2\pi\right) $. Các mệnh đề sau đúng hay sai?
 \choiceTF 
   {\True $ \cot \alpha<0$}
  {$ \tan \left(\pi-\alpha\right)=\tan  \alpha <0 $}
  {\True Nếu $ \sin \alpha =\dfrac{-3}{5}$  thì $ \cos \alpha = \dfrac{4}{5}$}
   {\True Nếu $ \sin 2\alpha=\dfrac{-\sqrt 3}{2} $ thì $ \left(\sin \alpha+\cos \alpha \right)^2 = \dfrac{2-\sqrt3}{2} $}
 \loigiai{ 
  \begin{itemchoice}
   \itemch Đúng. Vì giả sử điểm $ M $ biểu diễn góc $ \alpha $ thì $ M $ thuộc cung phần tư thứ tư. Suy ra $ \cos\alpha=x_M>0$, $\sin\alpha=y_M<0$, $ \tan \alpha <0  $, $\cot \alpha <0 $.
   \itemch Sai. Vì $ \tan\left(\pi-\alpha\right)=-\tan  \alpha >0$.
   \itemch Đúng. Vì $ \sin^2 \alpha+\cos^2 \alpha=1  \Rightarrow \cos^2 \alpha=\dfrac{16}{25} \Rightarrow \cos \alpha =\dfrac{\pm 4}{5} $.\\
   Vì  $ \alpha \in \left(\dfrac{3\pi}{2};2\pi\right) $ nên $ \cos \alpha >0 $.\\
   Vậy $ \cos  \alpha =\dfrac{4}{5} $.
   \itemch Đúng. Vì $ \left(\sin \alpha+\cos \alpha\right)^2=\sin^2\alpha+2\sin\alpha\cos\alpha+\cos^2\alpha=1+\sin2\alpha=1-\dfrac{\sqrt{3}}{2}=\dfrac{2-\sqrt{3}}{2}  $.
  \end{itemchoice}
  }
\end{ex}

% nguồn: Ok_dungsai_lop11_Lan2
\begin{ex}%[1D1H2-2]%[ma:198f4f30]
	Tính được các giá trị lượng giác của góc $\alpha$, biết  $\tan \alpha=2$, $0< \alpha < \dfrac{\pi}{2}$. Khi đó
	\choiceTF 
		{\True $\cot \alpha=\dfrac{1}{2}$} 
		{\True $\cos^2\alpha=\dfrac{1}{5}$} 
		{$\cos  \alpha=-\dfrac{\sqrt{5}}{5}$}
		 {\True $\sin \alpha=\dfrac{2\sqrt{5}}{5}$}
	 \loigiai{
		\begin{itemchoice}
			 \itemch \textbf{Đúng}.\\
			Ta  có $\tan \alpha=2$, $0< \alpha < \dfrac{\pi}{2}$.\\
			Khi  đó $\cot \alpha=\dfrac{1}{\tan \alpha}=\dfrac{1}{2}$.
			\itemch  \textbf{Đúng}.\\
			$1+\tan ^2\alpha=\dfrac{1}{\cos  ^2\alpha} \Rightarrow \cos ^2\alpha=\dfrac{1}{1+\tan^2 \alpha}=\dfrac{1}{5}$.
			\itemch  \textbf{Sai}.\\
			Vì $0< \alpha < \dfrac{\pi}{2}$  nên $\cos \alpha=\dfrac{\sqrt{5}}{5}$.
			 \itemch \textbf{Đúng}.\\
			$\tan \alpha=\dfrac{\sin \alpha}{\cos \alpha} \Rightarrow \sin  \alpha=\dfrac{2\sqrt{5}}{5}$.
		\end{itemchoice}
	} 
\end{ex}

% nguồn: Ok_dungsai_lop11_Lan2
\begin{ex}%[1D1H2-2]%[ma:408e2345]
	Cho $0^\circ<\alpha<90^\circ$. Xét được dấu của  các biểu thức sau. Khi đó
	\choiceTF 
		{\True $A=\sin \left( \alpha  +{{90}^\circ} \right)>0$}
		{\True $B=\cos \left( \alpha -{{45}^\circ}  \right)>0$}
		{$C=\tan \left(  {{270}^\circ}-\alpha \right)<0$}
		{$D=\cos \left( 2\alpha  +{{90}^\circ} \right)>0$}
	 \loigiai{
		\begin{itemchoice}
			 \itemch \textbf{Đúng.}\\
			Ta có $0^\circ<\alpha<90^\circ  \Rightarrow 90^\circ<\alpha+90^\circ<180^\circ\Rightarrow \sin \left(\alpha+90^\circ\right)>0$.
			 \itemch \textbf{Đúng.}\\
			Ta có $0^\circ<\alpha<90^\circ \Rightarrow-45^\circ<\alpha-45^\circ<45^\circ\Rightarrow \cos  \left(\alpha-45^\circ\right)>0$.
			 \itemch \textbf{Sai.}\\
			Ta có  $0^\circ<\alpha<90^\circ \Rightarrow-90^\circ<-\alpha<0^\circ$\\
			${\begin{aligned}
			& \Rightarrow 270^\circ+\left(-90^\circ\right)<270^\circ+(-\alpha)<270^\circ+0^\circ \\ 
			 & \Rightarrow 180^\circ<270^\circ-\alpha<270^\circ \\
			& \Rightarrow \tan \left(270^\circ-\alpha\right)>0. 
		\end{aligned}}$
		\itemch  \textbf{Sai.}\\
		Ta có $0^\circ<\alpha<90^\circ \Rightarrow 90^\circ<2 \alpha+90^\circ<270^\circ\Rightarrow \cos \left(2 \alpha+270^\circ\right)<0$. 
	\end{itemchoice}
} 
\end{ex}

% nguồn: Ok_dungsai_lop11_Lan2
\begin{ex}%[1D1H2-2]%[ma:240d35cc]
	Cho  biết $\tan \alpha =\dfrac{1}{2}$.
	 \choiceTF
		{$\cot \alpha  =\dfrac{1}{2}$}
		{\True  $\dfrac{1}{\cos^2\alpha }=\dfrac{5}{4}$}
		{\True  $\cot \alpha =2$}
		 {\True $\tan^2\alpha =\dfrac{1}{4}$}
	 \loigiai{
		\begin{itemchoice}
			\itemch  \textbf{Sai.}\\
			Vì $\tan  \alpha\cdot\cot \alpha =1\Leftrightarrow \cot \alpha =\dfrac{1}{\tan \alpha }=2$.
			 \itemch \textbf{Đúng.}\\
			Vì $\dfrac{1}{\cos^2\alpha }=1+\tan^2\alpha=1+\dfrac{1}{4}=\dfrac{5}{4}$. 
			 \itemch \textbf{Đúng.}\\
			Ta  có $\tan \alpha\cdot\cot \alpha =1\Leftrightarrow \cot \alpha =\dfrac{1}{\tan \alpha }=2$.
			 \itemch \textbf{Đúng.}\\
			Vì  $\tan^2\alpha ={{(\tan \alpha)}^{2}}=\dfrac{1}{4}$.
		\end{itemchoice}
	} 
\end{ex}

% nguồn: Ok_dungsai_lop11_Lan2
\begin{ex}%[1D1H2-2]%[ma:57d6015e]
	Tính được các giá trị lượng giác của  góc $\alpha$, biết $\sin \alpha=-\dfrac{\sqrt{7}}{4}$, $-\dfrac{\pi}{2} < \alpha < 0$. Khi đó
	 \choiceTF
		{\True  $\cos^2\alpha=\dfrac{9}{16}$}
		 {$\cos \alpha=-\dfrac{3}{4}$}
		{\True  $\cot \alpha=-\dfrac{3}{\sqrt{7}}$}
		 {$\tan \alpha+\cot \alpha=-\dfrac{16\sqrt{7}}{23}$}
	 \loigiai{
		\begin{itemchoice}
			\itemch \textbf{Đúng}.\\ 
			$\sin \alpha= - \dfrac{\sqrt{7}}{4}$, $-\dfrac{\pi}{2} < \alpha <  0$.\\
			 $\sin^2\alpha+ \cos^2 \alpha=1\Rightarrow \cos^2\alpha= 1 - \sin^2\alpha = 1 - \left(-\dfrac{\sqrt{7}}{4}\right)^2=\dfrac{9}{16}$.
			 \itemch \textbf{Sai}.\\
			Vì $-\dfrac{\pi}{2} < \alpha < 0$ nên $\cos \alpha >  0\Rightarrow \cos \alpha=\dfrac{3}{4}$.
			\itemch  \textbf{Đúng}.\\
			$\tan \alpha=\dfrac{\sin \alpha}{\cos \alpha}=-\dfrac{\sqrt{7}}{3}$. 
			$\cot  \alpha=-\dfrac{3}{\sqrt{7}}$.
			\itemch  \textbf{Sai}.\\
			 $\tan \alpha+\cot \alpha=-\dfrac{16\sqrt{7}}{21}$.
		\end{itemchoice}
	} 
\end{ex}

% nguồn: Ok_dungsai_lop11_Lan2
\begin{ex}%[1D1H2-2]%[ma:8aed975e]
	Cho $\cos x=\dfrac{1}{4} $ với  $0<x<\dfrac{\pi}{2} $. Khi đó
	 \choiceTF
		{$\sin  x<0$}
		{$\sin  x=-\dfrac{\sqrt{15}}{4} $}
		{\True $\tan  x=\sqrt{15} $}
		{$\cot x=-\dfrac{1}{\sqrt{15}} $} 
	\loigiai{ 
		\begin{itemchoice}
			\itemch Do  $0<x<\dfrac{\pi}{2} $ nên $\sin x>0$.
			\itemch Ta  có $\sin ^{2} x+\cos ^{2} x=1\Rightarrow \sin ^{2} x=1-\cos ^{2} x=1-\dfrac{1}{16} =\dfrac{15}{16} \Rightarrow \sin x=\dfrac{\sqrt{15}}{4} $.
			\itemch $\tan x=\dfrac{\sin x}{\cos x}  =\sqrt{15} $.
			\itemch $\cot x=\dfrac{\cos x}{\sin x} =\dfrac{1}{\sqrt{15}} $. 
		\end{itemchoice}
	} 
\end{ex}

% nguồn: Ok_dungsai_lop11_Lan2
\begin{ex}%[1D1H2-2]%[ma:5b291ed6]
	Cho $\sin\alpha=-\dfrac{12}{13},\,\,\dfrac{3\pi}{2}<\alpha <  2\pi $ . Xét tính đúng sai của các mệnh đề sau
	\choiceTF 
		{\True $\cos\left(\dfrac{\pi}{3}-\alpha\right)=\dfrac{5-12\sqrt  3}{26}$}
		{\True $\tan\left(\alpha+\dfrac{\pi}{3}\right)=\dfrac{5\sqrt  3-12}{12\sqrt 3+1}$}
		{$\sin\left(\dfrac{\pi}{3}+\alpha\right)=\dfrac{12\sqrt  3+12}{26}$}
		 {$\cot\left(\alpha+\dfrac{\pi}{3}\right)=\dfrac{48-25\sqrt 3}{11}$}
	 \loigiai{
		Ta có  $\dfrac{3\pi}{2}<\alpha < 2\pi \Rightarrow\cos\alpha > 0; \tan\alpha < 0; \cot\alpha < 0$.\\
		 Vì $\sin\alpha=-\dfrac{12}{13}\Rightarrow\cos\alpha=\sqrt{1-\sin^2\alpha}=\sqrt{1-\left(\dfrac{12}{13}\right)^2}=\dfrac{5}{13}$.\\
		Suy ra  $\tan\alpha=\dfrac{\sin\alpha}{\cos\alpha}=-\dfrac{12}{5}; \cot\alpha=-\dfrac{5}{12}.$\\
		\begin{itemchoice}
			 \itemch \textbf{Đúng.}\\
			 Vì $\cos\left(\dfrac{\pi}{3}-\alpha\right)=\cos\dfrac{\pi}{3}\cdot\cos\alpha+\sin\dfrac{\pi}{3}\cdot\sin\alpha=\dfrac{1}{2}.\dfrac{5}{13}+\dfrac{\sqrt 3}{2}.\left(-\dfrac{12}{13}\right)=\dfrac{5-12\sqrt 3}{26}$.\\
			\itemch \textbf{Đúng.}\\ 
			 $\tan\left(\alpha+\dfrac{\pi}{3}\right)=\dfrac{\tan\alpha+\tan\dfrac{\pi}{3}}{1-\tan\alpha .\tan\dfrac{\pi}{3}}=\dfrac{-\dfrac{12}{5}+\sqrt 3}{1+\dfrac{12}{5}.\sqrt 3}=\dfrac{5\sqrt 3-12}{12\sqrt 3+1}$.
			\itemch \textbf{Sai.}\\ 
			 $\sin\left(\dfrac{\pi}{3}+\alpha\right)=\sin\dfrac{\pi}{3}\cdot\cos\alpha+\cos\dfrac{\pi}{3}\cdot\sin\alpha=\dfrac{\sqrt 3}{2}.\dfrac{5}{13}+\dfrac{1}{2}.\left(-\dfrac{12}{13}\right)=\dfrac{5\sqrt 3-12}{26}$.
			 \itemch \textbf{Sai.}\\
			$\cot\left(\alpha+\dfrac{\pi}{3}\right)=\dfrac{1}{\tan\left(\alpha+\dfrac{\pi}{3}\right)}=\dfrac{12\sqrt  3+1}{5\sqrt 3-12}$.
		\end{itemchoice}
	} 
\end{ex}

% nguồn: Ok_dungsai_lop11_Lan2
\begin{ex}%[1D1H2-2]%[ma:22fd5aa1]
	Tính các giá trị lượng giác còn lại của góc $x$, biết: $\sin x=-\dfrac{3}{5}$ với $\pi<x<\dfrac{3 \pi}{2}$.  Khi đó
	\choiceTF 
		 {$\cos x>0$}
		{\True $\cos x=-\dfrac{4}{5}$; $\tan x=\dfrac{\sin x}{\cos x}=\dfrac{3}{4}$; $\cot x=\dfrac{\cos x}{\sin x}=\dfrac{4}{3}$} 
		{\True $\tan  x=\dfrac{3}{4}$}
		{\True $\cot x=\dfrac{4}{3}$} 
	\loigiai{ 
		Do $\pi<x<\dfrac{3 \pi}{2}$ nên $\cos x<0$.\\ 
		Ta có $\sin ^2 x+\cos ^2  x=1 \Rightarrow \cos ^2 x=1-\sin ^2 x=1-\dfrac{9}{25}=\dfrac{16}{25}$\\
		$\Rightarrow  \cos x=-\dfrac{4}{5}$; $\tan x=\dfrac{\sin x}{\cos x}=\dfrac{3}{4}$; $\cot x=\dfrac{\cos x}{\sin x}=\dfrac{4}{3}$.
		\begin{itemchoice}
			\itemch  \textbf{Sai.}
			\itemch  \textbf{Đúng.}
			\itemch \textbf{Đúng.} 
			\itemch \textbf{Đúng.} 
		\end{itemchoice}
	 }
\end{ex}

% nguồn: Ok_dungsai_lop11_Lan2
\begin{ex}%[1D1H2-2]%[ma:b81d8572]
	Tính được các giá trị lượng giác còn lại của góc $x$, biết $\cos x=-\dfrac{5}{13}$ với  $180^\circ < x < 270^\circ$, khi đó
	\choiceTF 
		 {\True $\sin x < 0$}
		{\True  $\tan x=\dfrac{12}{5}$}
		 {\True $\cot x=\dfrac{5}{12}$}
		{$\sin  x-\cos x=-\dfrac{12}{13}$}
	\loigiai{ 
		\begin{itemchoice}
			\itemch  \textbf{Đúng}.\\
			Do $180^\circ < x < 270^\circ \Rightarrow  \sin x < 0$.
			\itemch \textbf{Đúng}.\\ 
			Ta có $\sin  ^2x+\cos ^2x=1\Rightarrow \sin ^2x=1-\cos ^2x=1-\dfrac{25}{169}=\dfrac{144}{169}$
			$\Rightarrow \sin  x=-\dfrac{12}{13}$.\\
			$\tan x=\dfrac{\sin  x}{\cos x}=\dfrac{12}{5}$.
			\itemch  \textbf{Đúng}.\\
			$\cot  x=\dfrac{\cos x}{\sin x}=\dfrac{5}{12}$.\\
			 \itemch \textbf{Sai}.\\
			$\sin x-\cos x=  -\dfrac{12}{13} + \dfrac{5}{13} = -\dfrac{7}{13}$.
		\end{itemchoice}
	 }
\end{ex}

% nguồn: Ok_dungsai_lop11_Lan2
\begin{ex}%[1D1H2-2]%[ma:26a9b9b4]
	Tính được các giá trị lượng giác còn lại của góc $x$, biết $\tan x= - \dfrac{1}{3}$ với $\dfrac{\pi}{2} < x  < \pi$, khi đó
	\choiceTF{\True $\cos x < 0$} 
		{$\cos x=-\dfrac{\sqrt{10}}{10}$} 
		{$\sin  x=-\dfrac{\sqrt{10}}{10}$}
		 {$\sin x+\cos x=-\dfrac{\sqrt{10}}{5}$}
	\loigiai{ 
		\begin{itemchoice}
			\itemch  \textbf{Đúng}.\\
			Vì $\dfrac{\pi}{2}  < x < \pi$ nên $\cos x < 0$.
			 \itemch \textbf{Sai}.\\
			 Ta có $\dfrac{1}{\cos ^2x}=1+\tan ^2x=1+\dfrac{1}{9}=\dfrac{10}{9} \Rightarrow \cos ^2x=\dfrac{9}{10} \Rightarrow \cos x=-\dfrac{3\sqrt{10}}{10}$.
			\itemch \textbf{Sai}.\\ 
			$\sin x=\cos x \tan x=-\dfrac{3\sqrt{10}}{10}  \cdot \left(-\dfrac{1}{3}\right)=\dfrac{\sqrt{10}}{10}$.
			 \itemch \textbf{Sai}.\\
			$\sin x+\cos x=  \dfrac{\sqrt{10}}{10} -\dfrac{3\sqrt{10}}{10} = -\dfrac{\sqrt{10}}{5}$.
		\end{itemchoice}
	 }
\end{ex}

% nguồn: Ok_dungsai_lop11_Lan2
\begin{ex}%[1D1H2-2]%[ma:8daa66c0]
	Tính được các giá trị lượng giác của góc $\alpha$, biết $\sin \alpha=\dfrac{2}{3}$,  $\dfrac{\pi}{2} < \alpha < \pi$. Khi đó
	 \choiceTF
		 {\True $\dfrac{\pi}{2} < \alpha < \pi \Rightarrow \cos \alpha < 0$}
		{$\cos \alpha=\dfrac{\sqrt{5}}{3}$} 
		{\True  $\tan \alpha=-\dfrac{2}{\sqrt{5}}$}
		{\True $\cot  \alpha=-\dfrac{\sqrt{5}}{2}$}
	\loigiai{ 
		\begin{itemchoice}
			\itemch  \textbf{Đúng}.\\
			Vì  $\dfrac{\pi}{2} < \alpha < \pi$ nên $\cos \alpha < 0$.
			\itemch \textbf{Sai}.\\ 
			$\cos^2  \alpha=1-\sin ^2\alpha=1-\dfrac{4}{9}=\dfrac{5}{9}$.\\
			Vì $\cos \alpha < 0 $ nên $\cos \alpha  = -\dfrac{\sqrt{5}}{3}$
			\itemch  \textbf{Đúng}.\\
			$\tan \alpha=\dfrac{\sin \alpha}{\cos  \alpha}=\dfrac{\dfrac{2}{3}}{-\dfrac{\sqrt{5}}{3}}=-\dfrac{2}{\sqrt{5}}$.
			\itemch  \textbf{Đúng}.\\
			$\cot \alpha=\dfrac{1}{\tan \alpha}=\dfrac{1}{-\dfrac{2}{\sqrt{5}}}=-\dfrac{\sqrt{5}}{2}$. 
		\end{itemchoice}
	 }
\end{ex}

% nguồn: Ok_dungsai_lop11_Lan2
\begin{ex}%[1D1H2-2]%[ma:909c275d]
	Tính được các giá trị lượng giác của góc $\alpha$, biết $\cos \alpha=-\dfrac{7}{15}$ sao  cho
	$\pi < \alpha < \dfrac{3\pi}{2}$. Khi đó 
	\choiceTF 
		 {\True $\sin^2\alpha=\dfrac{176}{225}$}
		{$\sin \alpha=\dfrac{\sqrt{176}}{15}$} 
		{\True $\tan \alpha=\dfrac{\sqrt{176}}{7}$} 
		 {$\cot \alpha=-\dfrac{7}{\sqrt{176}}$}
	 \loigiai{
		\begin{itemchoice}
			\itemch  \textbf{Đúng}.\\
			Ta có  $\sin^2\alpha+\cos ^2\alpha=1$ nên $\sin ^2\alpha = 1- \cos ^2\alpha= 1 - \left(-\dfrac{7}{15}\right)^2=\dfrac{176}{225}$.
			\itemch \textbf{Sai}.\\ 
			Vì $\pi < \alpha < \dfrac{3\pi}{2}$ nên $\sin \alpha < 0$ suy ra $\sin \alpha=-\dfrac{\sqrt{176}}{15}$. 
			\itemch \textbf{Đúng}.\\ 
			$\tan  \alpha=\dfrac{\sin \alpha}{\cos \alpha}=\dfrac{\sqrt{176}}{7}$.
			\itemch \textbf{Sai}.\\ 
			$\cot \alpha=\dfrac{\cos \alpha}{\sin \alpha}=\dfrac{7}{\sqrt{176}}$. 
		\end{itemchoice}
	 }
\end{ex}

% nguồn: Ok_dungsai_lop11_Lan2
\begin{ex}%[1D1H2-2]%[ma:eafd4a65]
	Tính được giá trị các biểu thức. Khi  đó
	\choiceTF 
		 {\True $A=\dfrac{1}{\sqrt3}\cos30^\circ-3\sqrt2\sin45^\circ+\cot45^\circ=-\dfrac{3}{2}$}
		{\True  $B=\sin^260^\circ+\tan^230^\circ-2=-\dfrac{11}{12}$}
		 {\True $C=2\sin\dfrac{\pi}{6}+\sqrt{3}\cos\dfrac{\pi}{6}=\dfrac{5}{2}$}
		{$D=\dfrac{\tan\dfrac{\pi}{4}+1}{\dfrac{1}{\sqrt{2}}\cot\dfrac{\pi}{2}-2}=\dfrac{6}{5}$} 
	\loigiai{ 
		\begin{itemchoice}
			\itemch  \textbf{Đúng.}\\
			 $A=\dfrac{1}{\sqrt{3}} \cdot \dfrac{\sqrt{3}}{2}-3 \sqrt{2} \cdot \dfrac{\sqrt{2}}{2}+1=\dfrac{1}{2}-3+1=-\dfrac{3}{2}$.
			 \itemch \textbf{Đúng.}\\
			$B=\left(\dfrac{\sqrt{3}}{2}\right)^2+\left(\dfrac{1}{\sqrt{3}}\right)^2-2=\dfrac{3}{4}+\dfrac{1}{3}-2=-\dfrac{11}{12}$. 
			\itemch \textbf{Đúng.}\\ 
			$C=2 \cdot \dfrac{1}{2}+\sqrt{3}  \cdot \dfrac{\sqrt{3}}{2}=\dfrac{5}{2}$.
			 \itemch \textbf{Sai.}\\
			$D=\dfrac{1+1}{\dfrac{1}{\sqrt{3}} \cdot \dfrac{1}{\sqrt{3}}-2}=-\dfrac{6}{5}$. 
		\end{itemchoice}
	 }
\end{ex}

% nguồn: Ok_dungsai_lop11_Lan2
\begin{ex}%[1D1H2-2]%[ma:e2868687]
	Tính được các giá trị  lượng giác của góc $\alpha$, biết $\cot \alpha=\sqrt{2}+1$, $0< \alpha < \dfrac{\pi}{2}$. Khi đó
	 \choiceTF
		{\True 
		$\heva{ 
		&\cos  \alpha > 0\\
		 &\sin \alpha > 0.}$}
		 {$\tan \alpha=\sqrt{2}+1$}
		 {\True
		 $\sin \alpha=\dfrac{\sqrt{2-\sqrt{2}}}{2}$}
		{\True $\cos  \alpha=\dfrac{\sqrt{2+\sqrt{2}}}{2}$}
	\loigiai{ 
		\begin{itemchoice}
			 \itemch \textbf{Đúng}.\\
			Vì $0< \alpha < \dfrac{\pi}{2}$ nên $0< \alpha < \dfrac{\pi}{2}  \Rightarrow \heva{&\cos \alpha > 0\\ &\sin \alpha > 0.}$
			\itemch \textbf{Sai}.\\ 
			$\tan \alpha  = \dfrac{1}{\cot \alpha}=\dfrac{1}{\sqrt{2}+1}=\sqrt{2}-1$.
			\itemch  \textbf{Đúng}.\\
			$\sin  ^2 \alpha=\dfrac{1}{\cot^2\alpha+1}=\dfrac{1}{(\sqrt{2}+1)^2+1}$.\\
			Vì $\sin \alpha > 0$ nên $\sin \alpha = \dfrac{\sqrt{2-\sqrt{2}}}{2}$ 
			 \itemch \textbf{Đúng}.\\
			$\cos^2 \alpha = 1 - \sin^2 \alpha = 1 - \dfrac{2  - \sqrt{2}}{4}$.\\
			Vì  $\cos \alpha > 0$ nên $\cos \alpha = \dfrac{\sqrt{2+\sqrt{2}}}{2}$.
		\end{itemchoice}
	} 
\end{ex}

% nguồn: Ok_dungsai_lop11_Lan2
\begin{ex}%[1D1H2-2]%[ma:aabfb3ef]
	Cho $\cos x=-\dfrac{5}{13} $ với $180^\circ  <x<270^\circ $. Khi đó
	 \choiceTF
		{\True $\sin x<0$} 
		{\True $\tan x=\dfrac{12}{5} $} 
		{\True  $\cot x=\dfrac{5}{12} $}
		 {$\sin x-\cos x=-\dfrac{12}{13} $}
	 \loigiai{
		\begin{itemchoice}
			\itemch Do  $180^\circ <x<270^\circ \Rightarrow \sin x<0$.
			\itemch Ta có  $\sin ^{2} x+\cos ^{2} x=1\Rightarrow \sin ^{2} x=1-\cos ^{2} x=1-\dfrac{25}{169} =\dfrac{144}{169}\Rightarrow \sin x=-\dfrac{12}{13} $; $\tan x=\dfrac{\sin x}{\cos x} =\dfrac{12}{5}$.
			\itemch $\cot x=\dfrac{\cos  x}{\sin x} =\dfrac{5}{12} $.
			 \itemch $\sin x-\cos x=-\dfrac{7}{13} $.
	\end{itemchoice}}
\end{ex}

% nguồn: Ok_dungsai_lop11_Lan2
\begin{ex}%[1D1H2-2]%[ma:92db9825]
	Gọi $M$  là điểm cuối khi biểu diễn cung lượng giác $\alpha $ trên đường tròn lượng giác.
	 \choiceTF
		{Nếu $M$ nằm  bên phải trục tung thì $\cos \alpha <0$}
		{\True Nếu $M$ thuộc góc phần tư thứ tư thì $\sin \alpha <0$  và $\cos \alpha >0$}
		{Nếu $M$ thuộc góc phần tư thứ hai thì $\sin \alpha >0$ và $\cos  \alpha >0$}
		{\True Nếu $M$ nằm phía trên  trục hoành thì $\sin \alpha >0$}
	 \loigiai{
		Theo  lí thuyết ta có
		\begin{itemchoice}
			\itemch \textbf{Sai.} 
			\itemch \textbf{Đúng.} 
			\itemch  \textbf{Sai.}
			\itemch  \textbf{Đúng.}
		\end{itemchoice}
	 }
\end{ex}

% nguồn: Ok_dungsai_lop11_Lan2
\begin{ex}%[1D1H2-2]%[ma:00616556]
	Tính được các giá trị lượng giác của góc  $\alpha =-\dfrac{\pi}{4}+(2k+1)\pi$ (biết ${k \in \mathbb{Z}}$). Khi đó
	\choiceTF 
		{\True $\sin  \alpha =\dfrac{\sqrt{2}}{2}$}
		{\True  $\cos \alpha =\dfrac{\sqrt{2}}{2}$}
		{$\tan \alpha =-1$} 
		 {$\cot \alpha =-1$}
	 \loigiai{
		\begin{itemchoice}
			 \itemch \textbf{Đúng.}\\
			Ta  có
			\begin{align*}
				\sin \left[-\dfrac{\pi}{4}+(2 k+1) \pi\right]&=\sin \left(-\dfrac{\pi}{4}+2 k  \pi+\pi\right)\\
				&=\sin \left( -\dfrac{\pi}{4}+\pi \right)=-\sin \left( -\dfrac{\pi}{4} \right)=\sin \dfrac{\pi}{4}=\dfrac{\sqrt{2}}{2}. 
			\end{align*}
			 \itemch \textbf{Đúng.}\\
			 Ta có
			\begin{align*}
				\cos \left[-\dfrac{\pi}{4}+(2k+1)\pi \right]&=\cos \left( -\dfrac{\pi}{4}+2k\pi +\pi \right)=\cos \left(  -\dfrac{\pi}{4}+\pi \right) \\
				&=-\cos \left( -\dfrac{\pi}{4}  \right)=-\cos \dfrac{\pi}{4}=-\dfrac{\sqrt{2}}{2}.
			\end{align*}
			 \itemch \textbf{Sai.}\\
			 Ta có
			\begin{align*}
				\tan \left[-\dfrac{\pi}{4}+(2k+1)\pi \right]&=\tan \left( -\dfrac{\pi}{4}  \right)=-\tan \dfrac{\pi}{4}=-1.
			\end{align*}
			\itemch \textbf{Sai.}\\ 
			Ta  có
			\begin{align*}
				\cot \left[-\dfrac{\pi}{4}+(2k+1)\pi \right]&=\cot  \left( -\dfrac{\pi}{4} \right)=-\cot \dfrac{\pi}{4}=-1.
			\end{align*}
		\end{itemchoice}
	} 
\end{ex}

% nguồn: Ok_dungsai_lop11_Lan2
\begin{ex}%[1D1H2-2]%[ma:60420686]
	Cho $\cot  x=2$. Tính được các biểu thức $B_1=\dfrac{2\sin x+3\cos x}{3\sin x-2\cos x}$, $B_2=\dfrac{2}{\cos ^2x-\sin x\cos x}$, khi đó
	 \choiceTF
		{\True Vì $\cot x=2$ nên  $\sin x\ne 0$}
		{\True  $B_1=-8$}
		 {$B_2=-5$}
		 {$B_1+B_2=-13$}
	\loigiai{ 
		\begin{itemchoice}
			\itemch \textbf{Đúng}.\\ 
			Vì $\cot x=2$ nên $\sin  x\ne 0$.
			\itemch \textbf{Đúng}.\\ 
			Chia cả tử và mẫu của  biểu thức $B_1$ cho $\sin x$, ta được
			\[B_1=\dfrac{2\dfrac{\sin  x}{\sin x}+3\dfrac{\cos x}{\sin x}}{3\dfrac{\sin x}{\sin x}-2\dfrac{\cos x}{\sin x}}=\dfrac{2+3\cot x}{3-2\cot x}=\dfrac{2+3\cdot 2}{3-2\cdot 2}=-8.\]
			 \itemch \textbf{Sai}.\\
			Chia cả tử và mẫu của  biểu thức $B_2$ cho $\sin ^2x$, ta được
			\[B_2=\dfrac{\dfrac{2}{\sin ^2x}}{\dfrac{\cos  ^2x}{\sin ^2x}-\dfrac{\sin x\cos x}{\sin ^2x}}=\dfrac{2\left(1+\cot ^2x\right)}{\cot ^2x-\cot x}=\dfrac{2\left(1+2^2\right)}{2^2-2}=5.\]
			 \itemch \textbf{Sai}.\\
			Ta có  $B_1+B_2=-8+5=-3$.
		\end{itemchoice}
	} 
\end{ex}

% nguồn: Ok_dungsai_lop11_Lan2
\begin{ex}%[1D1H2-2]%[ma:407a6a35]
	Cho  $\sin x=-\dfrac{3}{5} $ với $\pi <x<\dfrac{3\pi}{2} $. Khi đó
	 \choiceTF
		 {$\cos x>0$}
		{\True  $\cos x=-\dfrac{4}{5} $}
		{\True $\tan x=\dfrac{3}{4} $} 
		{\True $\cot  x=\dfrac{4}{3} $}
	\loigiai{ 
		\begin{itemchoice}
			\itemch Do $\pi <x<\dfrac{3\pi}{2} $  nên $\cos x<0$.
			\itemch Ta có  $\sin ^{2} x+\cos ^{2} x=1\Rightarrow \cos ^{2} x=1-\sin ^{2} x=1-\dfrac{9}{25} =\dfrac{16}{25} $$\Rightarrow \cos x=-\dfrac{4}{5} $.
			\itemch $\tan x=\dfrac{\sin x}{\cos x}  =\dfrac{3}{4} $.
			\itemch $\cot  x=\dfrac{\cos x}{\sin x} =\dfrac{4}{3} $.
	\end{itemchoice}}
\end{ex}

% nguồn: Ok_dungsai_lop11_Lan2
\begin{ex}%[1D1H2-2]%[ma:36b608d8]
	Tính được các giá trị lượng giác  của góc $\alpha$, biết $\cos \alpha=-\dfrac{3}{4}$, $-\dfrac{3\pi}{2} < \alpha <-\pi$. Khi đó
	 \choiceTF
		{$\sin \alpha < 0$} 
		{$\sin \alpha=-\dfrac{\sqrt{7}}{4}$} 
		 {\True $\tan \alpha=-\dfrac{\sqrt{7}}{3}$}
		{\True $\cot  \alpha=-\dfrac{3}{\sqrt{7}}$}
	\loigiai{ 
		\begin{itemchoice}
			\itemch  \textbf{Sai}.\\
			Vì $-\dfrac{3\pi}{2}  < \alpha <-\pi$ nên $\sin \alpha > 0$.
			\itemch  \textbf{Sai}.\\
			$\sin^2 \alpha = 1 - \cos^2 \alpha =  1 - \dfrac{9}{16} = \dfrac{7}{16}$.\\
			Vì $\sin \alpha > 0$ nên $\sin \alpha = \dfrac{\sqrt{7}}{4}$ 
			 \itemch \textbf{Đúng}.\\
			$\tan \alpha=\dfrac{\sin  \alpha}{\cos \alpha}=\dfrac{\dfrac{\sqrt{7}}{4}}{-\dfrac{3}{4}}=-\dfrac{\sqrt{7}}{3}$.
			\itemch \textbf{Đúng}.\\ 
			 $\cot \alpha=-\dfrac{3}{\sqrt{7}}$.
		\end{itemchoice}
	 }
\end{ex}

% nguồn: Ok_dungsai_lop11_Lan2
\begin{ex}%[1D1H2-2]%[ma:4ab5ad10]
	Tính được các giá trị lượng giác của góc $\alpha =-\dfrac{\pi }{4} +(2k+1)\pi $ ($k\in \mathbb{Z}$).  Khi đó
	 \choiceTF
		{\True $\sin \alpha =\dfrac{\sqrt{2} }{2}  $}
		{\True $\cos \alpha =\dfrac{\sqrt{2} }{2} $} 
		{$\tan \alpha =-1$} 
		{$\cot \alpha  =-1$}
	\loigiai{ 
		\begin{itemchoice}
			\itemch Ta có \begin{eqnarray*}
			\sin \left[-\dfrac{\pi }{4}  +(2k+1)\pi \right]&=&\sin \left(-\dfrac{\pi }{4} +2k\pi +\pi \right)\\
			 &=&\sin \left(-\dfrac{\pi }{4} +\pi \right)\\
			&=&-\sin  \left(-\dfrac{\pi }{4} \right)\\
			&=&\sin  \dfrac{\pi }{4} =\dfrac{\sqrt{2}}{2}.
		\end{eqnarray*}
		\itemch Ta  có
		\begin{eqnarray*}
			\cos \left[-\dfrac{\pi }{4} +(2k+1)\pi \right] &=&\cos \left(-\dfrac{\pi }{4} +2k\pi  +\pi \right)\\
			&=&\cos \left(-\dfrac{\pi }{4} +\pi  \right)\\
			&=&-\cos  \left(-\dfrac{\pi }{4} \right)\\
			&=&-\cos  \dfrac{\pi }{4} =-\dfrac{\sqrt{2}}{2}.
		\end{eqnarray*}
		\itemch  $\tan \left[-\dfrac{\pi }{4} +(2k+1)\pi \right]=\tan \left(-\dfrac{\pi }{4} \right)=-\tan \dfrac{\pi }{4} =-1$.
		\itemch $\cot \left[-\dfrac{\pi }{4} +(2k+1)\pi \right]=\cot \left(-\dfrac{\pi }{4} \right)=-\cot \dfrac{\pi }{4} =-1$. 
	\end{itemchoice}
 }
\end{ex}

% nguồn: Ok_dungsai_lop11_Lan2
\begin{ex}%[1D1H2-2]%[ma:95ad16fa]
	Cho $0<\alpha <\dfrac{\pi}{2} $.  Khi đó
	 \choiceTF
		{\True  $A=\cos (\alpha +\pi)<0$}
		{\True $B=\tan  (\alpha -\pi)>0$}
		{$C=\sin \left(\alpha +\dfrac{2\pi}{5}\right)<0$} 
		{$D=\cos \left(\alpha  -\dfrac{3\pi}{8}\right)<0$}
	 \loigiai{
		\begin{itemchoice}
			\itemch Vì $0<\alpha <\dfrac{\pi}{2} \Rightarrow \pi <\alpha  +\pi <\dfrac{3\pi}{2} \Rightarrow \cos (\alpha +\pi)<0$.
			\itemch Vì $0<\alpha <\dfrac{\pi}{2}  \Rightarrow -\pi <\alpha -\pi <-\dfrac{\pi}{2} \Rightarrow \tan (\alpha -\pi)>0$.
			 \itemch Vì $0<\alpha <\dfrac{\pi}{2} \Rightarrow \dfrac{2\pi}{5} <\alpha +\dfrac{2\pi}{5} <\dfrac{9\pi}{10} \Rightarrow \sin \left(\alpha +\dfrac{2\pi}{5}\right)>0$.
			\itemch Vì $0<\alpha <\dfrac{\pi}{2} \Rightarrow -\dfrac{3\pi}{8} <\alpha -\dfrac{3\pi}{8} <\dfrac{\pi}{8} \Rightarrow  \cos \left(\alpha -\dfrac{3\pi}{8}\right)>0$.
	\end{itemchoice}}
\end{ex}

% nguồn: Ok_dungsai_lop11_Lan2
\begin{ex}%[1D1H2-2]%[ma:eb45f2fc]
	Cho góc lượng giác $\alpha$ thỏa mãn $\cos  \alpha=\dfrac{3}{5}$, $0< \alpha < \dfrac{\pi}{2}$.
	 \choiceTF
		{Giá trị $\sin$ của góc lượng giác $\alpha$ âm} 
		{\True Giá trị $\tan$ của góc lượng  giác $\alpha$ dương}
		{\True Giá trị của biểu thức $P=\sin  ^2\alpha-2\cos ^2\alpha=\dfrac{2}{25}$}
		{\True Giá trị của biểu thức $Q=\tan \alpha+\cot \alpha > 2$} 
	\loigiai{ 
		\begin{itemchoice}
			 \itemch \textbf{Sai}.\\
			Vì $0< \alpha  < \dfrac{\pi}{2}$ nên $\sin \alpha > 0$
			\itemch \textbf{Đúng}.\\ Vì $0< \alpha < \dfrac{\pi}{2}$ nên $\sin \alpha  > 0, \cos \alpha > 0\Rightarrow \tan \alpha > 0$
			 \itemch \textbf{Đúng}.\\
			Vì $\cos \alpha=\dfrac{3}{5},  0< \alpha < \dfrac{\pi}{2}$.\\
			$\Rightarrow \sin \alpha=\sqrt{1-\cos  ^2\alpha}=\dfrac{4}{5}$.\\
			\[P=\sin ^2\alpha-2\cos ^2\alpha=\dfrac{16}{25}-\dfrac{18}{25}=-\dfrac{2}{25}.\] 
			 \itemch \textbf{Đúng}.\\
			 $Q=\tan \alpha+\cot \alpha=\dfrac{4}{3}+\dfrac{3}{4}=\dfrac{25}{12} > 2$.
		\end{itemchoice}
	} 
\end{ex}

% nguồn: Ok_dungsai_lop11_Lan2
\begin{ex}%[1D1H2-2]%[ma:1ff230c9]
	 Cho $0^\circ <\alpha <90^\circ $. Khi đó
	 \choiceTF
		 {\True $A=\sin (\alpha +90^\circ)>0$}
		{\True $B=\cos (\alpha -45^\circ)>0$} 
		{$C=\tan (270^\circ -\alpha)<0$} 
		{$D=\cos  (2\alpha +90^\circ)>0$}
	 \loigiai{
		\begin{itemchoice}
			\itemch Ta có $0^\circ <\alpha <90^\circ \Rightarrow 90^\circ <\alpha +90^\circ <180^\circ $$\Rightarrow  \sin (\alpha +90^\circ)>0$.
			\itemch Ta có $0^\circ <\alpha <90^\circ \Rightarrow -45^\circ <\alpha -45^\circ  <45^\circ $$\Rightarrow \cos (\alpha -45^\circ)>0$.
			\itemch  Ta có
			\begin{eqnarray*}
				0^\circ <\alpha <90^\circ &\Rightarrow& -90^\circ <-\alpha <0^\circ\\ &\Rightarrow& 270^\circ +(-90^\circ)<270^\circ  +(-\alpha)<270^\circ +0^\circ\\
				&\Rightarrow& 180^\circ <270^\circ -\alpha <270^\circ. 
			\end{eqnarray*}
			Suy ra  $ \tan (270^\circ -\alpha)>0$.
			\itemch Ta có $0^\circ <\alpha <90^\circ \Rightarrow 90^\circ <2\alpha +90^\circ <270^\circ $$\Rightarrow  \cos (2\alpha +270^\circ)<0$
	\end{itemchoice}}
\end{ex}

% nguồn: Ok_dungsai_lop11_Lan2
\begin{ex}%[1D1H2-2]%[ma:3d4dbcd5]
	Cho $\tan x=-2$. Tính được các biểu thức $A_1=\dfrac{5\cot x+4\tan x}{5\cot x-4\tan x}$, $A_2=\dfrac{2\sin x+\cos x}{\cos  x-3\sin x}$, khi đó
	\choiceTF 
		{\True $\cot x=-\dfrac{1}{2}$} 
		{Vì $\tan  x=-2$ nên $\cos x=0$}
		 {\True $A_1=-\dfrac{21}{11}$}
		 {$A_2=\dfrac{3}{7}$}
	 \loigiai{
		\begin{itemchoice}
			\itemch  \textbf{Đúng}.\\
			Ta có $\tan x=-2\Rightarrow \cot x=-\dfrac{1}{2}$. 
			 \itemch \textbf{Sai}.\\
			Ta có $\tan x=-2$ nên $\cos x\ne  0$.
			\itemch \textbf{Đúng}.\\ $A_1=\dfrac{-\dfrac{5}{2}+4\cdot \left(-2\right)}{-\dfrac{5}{2}-4\cdot  \left(-2\right)}=-\dfrac{21}{11}$.
			\itemch  \textbf{Sai}.\\
			Chia tử và mẫu của biểu thức  $A_2$ cho $\cos x$, ta được
			\[A_2=\dfrac{\dfrac{2\sin x}{\cos x}+\dfrac{\cos x}{\cos x}}{\dfrac{\cos x}{\cos x}-\dfrac{3\sin x}{\cos x}}=\dfrac{2\tan  x+1}{1-3\tan x}=\dfrac{2\cdot \left(-2\right)+1}{1-3\cdot \left(-2\right)}=-\dfrac{3}{7}.\]
		\end{itemchoice}
	} 
\end{ex}

% nguồn: Ok_dungsai_lop11_Lan2
\begin{ex}%[1D1H2-2]%[ma:91721033]
	Trên đường tròn lượng giác, điểm $M$ thỏa mãn  $\left( Ox,OM \right)=500^\circ $. Xét tính đúng sai của các mệnh đề sau:
	\choiceTF 
		{\True $M$ nằm ở góc phần tư thứ  II}
		{\True  $\sin \left( Ox,OM \right)>0$; $\cos\left( Ox,OM \right)<0$}
		 {$M$ có hoành độ âm và tung độ âm}
		{$M$ có hoành độ dương và tung độ  âm}
	\loigiai{ 
		\begin{itemchoice}
			\itemch \textbf{Đúng.}\\ 
			Điểm $M$ thỏa mãn $\left( Ox,OM \right)=500^\circ $ thì nằm ở góc phần tư thứ II vì $500^\circ -360^\circ =140^\circ \in \left(  90^\circ;180^\circ \right)$.
			\itemch \textbf{Đúng.}\\ 
			 Ta có $\sin(Ox,OM)=\sin(140^\circ+360^\circ)=\sin140^\circ$. Mà $140^\circ\in(90^\circ;180^\circ)$ nên $\sin(Ox,OM)>0$.\\
			 Tương tự, $\cos(Ox,OM)<0$.
			\itemch  \textbf{Sai.}\\
			 Do $M$ thuộc góc phần tư thứ II.
			\itemch  \textbf{Sai.}\\
			Do $M$ thuộc góc  phần tư thứ II.
		\end{itemchoice}
	 }
\end{ex}

% nguồn: Ok_dungsai_lop11_Lan2
\begin{ex}%[1D1H2-2]%[ma:c964dabd]
	Tính được các giá trị lượng giác của góc $\alpha =\dfrac{\pi }{3} +k2\pi $ $(k\in  \mathbb{Z})$. Khi đó
	\choiceTF 
		{$\sin \alpha =-\dfrac{\sqrt{3}}{2} $} 
		 {$\cos \alpha =-\dfrac{1}{2} $}
		{\True $\tan \alpha =\sqrt{3}  $}
		{$\cot  \alpha =-\dfrac{\sqrt{3}}{3} $}
	\loigiai{ 
		\begin{itemchoice}
			\itemch $\sin \left(\dfrac{\pi }{3} +k2\pi \right)=\sin \dfrac{\pi  }{3} =\dfrac{\sqrt{3}}{2} $.
			\itemch $\cos \left(\dfrac{\pi }{3} +k2\pi \right)=\cos \dfrac{\pi }{3} =\dfrac{1}{2} $. 
			\itemch $\tan \left(\dfrac{\pi }{3} +k2\pi \right)=\tan  \dfrac{\pi }{3} =\sqrt{3} $.
			\itemch $\cot \left(\dfrac{\pi }{3} +k2\pi \right)=\cot \dfrac{\pi }{3} =\dfrac{\sqrt{3}}{3}  $.
		\end{itemchoice}
	 }
\end{ex}

% nguồn: Ok_dungsai_lop11_Lan2
\begin{ex}%[1D1H2-2]%[ma:b26a1ed1]
	Tính được các giá trị  lượng giác của góc $\alpha$, biết $\tan \alpha=\dfrac{2\sqrt{10}}{9}$, $\pi < \alpha < \dfrac{3\pi}{2}$. Khi đó
	 \choiceTF
		 {\True $\heva{
		 &\cos \alpha < 0\\
		&\sin  \alpha < 0
		 }$}
		{\True  $\cos \alpha=-\dfrac{9}{11}$}
		{\True $\cot \alpha=\dfrac{9}{2\sqrt{10}}$} 
		 {\True $\sin \alpha=-\dfrac{2\sqrt{10}}{11}$}
	 \loigiai{
		\begin{itemchoice}
			\itemch  \textbf{Đúng}.\\
			Vì  $\pi < \alpha < \dfrac{3\pi}{2}$ nên $\heva{
			&\cos \alpha < 0\\ 
			 &\sin \alpha < 0.}$
			 \itemch \textbf{Đúng}.\\
			Ta có $1 + \tan^2 \alpha = \dfrac{1}{\cos^2 \alpha}$ suy ra $\cos^2  \alpha = \dfrac{1}{1+\tan^2 \alpha} = \dfrac{1}{\dfrac{40}{81}+1} = \dfrac{81}{121}$.\\
			Vì $\cos \alpha < 0$  nên $\cos \alpha = \dfrac{9}{11}$\\
			\itemch  \textbf{Đúng}.\\
			$\cot \alpha=\dfrac{9}{2\sqrt{10}}$. 
			 \itemch \textbf{Đúng}.\\
			$\sin  \alpha = 1 - \cos^2 \alpha = 1 - \dfrac{81}{121} = \dfrac{40}{121}$\\
			Vì $\sin \alpha < 0 $  nên $\sin \alpha = \dfrac{2\sqrt{10}}{11}$.
		\end{itemchoice}
	 }
\end{ex}

% nguồn: Ok_dungsai_lop11_Lan2
\begin{ex}%[1D1H2-2]%[ma:6e3822f7]
	Cho góc $\alpha $ thỏa mãn $2\pi <\alpha <\dfrac{5\pi}{2}$. 
	 \choiceTF
		 {$\tan \alpha <0$}
		{\True $\cot \alpha  >0$}
		{\True $\sin \alpha >0$} 
		{\True $\cos  \alpha >0$}
	\loigiai{ 
		 Với $2\pi <\alpha <\dfrac{5\pi}{2}$ góc phần tư thứ nhất nên ta có $\sin \alpha >0$, $\cos \alpha >0$, $\tan \alpha >0$, $\cot \alpha >0$.
		Do  đó
		\begin{itemchoice}
			 \itemch \textbf{Sai.}
			 \itemch \textbf{Đúng.}
			\itemch  \textbf{Đúng.}
			\itemch \textbf{Đúng.} 
		\end{itemchoice}
	 }
\end{ex}

% nguồn: Ok_dungsai_lop11_Lan2
\begin{ex}%[1D1H2-2]%[ma:c784bb40]
	Tính các giá trị lượng giác còn lại của góc $x$, biết: $\cos x=\dfrac{1}{4}$ với $0<x<\dfrac{\pi}{2}$. Khi  đó
	\choiceTF 
		{$\sin x<0$} 
		{$\sin  x=-\dfrac{\sqrt{15}}{4}$}
		{\True $\tan  x=\sqrt{15}$}
		 {$\cot x=-\dfrac{1}{\sqrt{15}}$}
	\loigiai{ 
		\begin{itemchoice}
			\itemch  \textbf{Sai.}
			Do $0<x<\dfrac{\pi}{2}$ nên ${\sin  x>0}$.
			 \itemch \textbf{Sai.}
			 Do $0<x<\dfrac{\pi}{2}$ nên ${\sin x>0}$.\\
			Ta có $\sin ^2 x+\cos ^2 x=1 \Rightarrow \sin ^2 x=1-\cos  ^2 x=1-\dfrac{1}{16}=\dfrac{15}{16}$\\
			$\Rightarrow \sin x=\dfrac{\sqrt{15}}{4}$. 
			 \itemch \textbf{Đúng.}
			Do $0<x<\dfrac{\pi}{2}$  nên ${\sin x>0}$.\\
			Ta có $\sin ^2 x+\cos  ^2 x=1 \Rightarrow \sin ^2 x=1-\cos ^2 x=1-\dfrac{1}{16}=\dfrac{15}{16}$\\
			$\Rightarrow \sin x=\dfrac{\sqrt{15}}{4}$; $\tan  x=\dfrac{\sin x}{\cos x}=\sqrt{15}$.
			\itemch  \textbf{Sai.}
			 Do $0<x<\dfrac{\pi}{2}$ nên ${\sin x>0}$.\\
			Ta có $\sin ^2 x+\cos ^2 x=1 \Rightarrow \sin  ^2 x=1-\cos ^2 x=1-\dfrac{1}{16}=\dfrac{15}{16}$\\
			$\Rightarrow  \sin x=\dfrac{\sqrt{15}}{4}$; $\tan x=\dfrac{\sin x}{\cos x}=\sqrt{15}$; $\cot x=\dfrac{\cos x}{\sin x}=\dfrac{1}{\sqrt{15}}$.
		\end{itemchoice}
	 }
\end{ex}

% nguồn: Ok_dungsai_lop11_Lan2
\begin{ex}%[1D1H2-2]%[ma:77f6a3d0]
	Cho  $\cot x=2$. Tính được các biểu thức $B_1=\dfrac{2\sin x+3\cos x}{3\sin x-2\cos x}$, $B_2=\dfrac{2}{\cos ^2x-\sin x \cos x}$, khi đó
	 \choiceTF
		{\True $\sin x \neq  0$}
		{\True  $B_1=-8$}
		{$B_2=-5$} 
		{$B_1+B_2=-13$} 
	\loigiai{ 
		\begin{itemchoice}
			\itemch  \textbf{Đúng}.\\
			Vì $\cot x=2$ nên $\sin x  \neq 0$.\\
			\itemch  \textbf{Đúng}.\\
			Chia cả tử và  mẫu của biểu thức $B_1$ cho $\sin x$, ta được
			\[B_1=\dfrac{2\cdot \dfrac{\sin x}{\sin x}+3\cdot \dfrac{\cos x}{\sin x}}{3\cdot \dfrac{\sin x}{\sin x}-2\cdot \dfrac{\cos x}{\sin x}}=\dfrac{2+3\cot x}{3-2\cot  x}=\dfrac{2+3\cdot 2}{3-2\cdot 2}=-8.
			 \]
			 \itemch \textbf{Sai}.\\
			Chia cả tử và mẫu của biểu thức $B_2$ cho ${\sin ^2x}$,  ta được
			\[B_2=\dfrac{\dfrac{2}{\sin ^2x}}{\dfrac{\cos ^2x}{\sin ^2x}-\dfrac{\sin  x \cos x}{\sin ^2x}}=\dfrac{2\left(1+\cot ^2x\right)}{\cot ^2x-\cot x}=\dfrac{2\left(1+2^2\right)}{2^2-2}=5.
			 \]
			\itemch \textbf{Sai}.\\ 
			$B_1 + B_2  = -3$
		\end{itemchoice}
	} 
\end{ex}

% nguồn: Ok_dungsai_lop11_Lan1
\begin{ex}%[1D1V2-2]%[ma:e6e4e2ed]
 Cho góc  lượng giác $\alpha$ thỏa mãn $\cos\alpha=\dfrac{3}{5},0<\alpha<\dfrac{\pi}{2}$. Xét tính đúng-sai của các phát biểu sau:
 \choiceTF 
  {Giá trị sin của góc lượng giác $\alpha$  âm}
  {\True  Giá trị tan của góc lượng giác $\alpha$ dương}
  {Giá  trị của biểu thức $P=\sin^2\alpha-2\cos^2\alpha=\dfrac{2}{25}$}
  {\True  Giá trị của biểu thức $Q=\tan\alpha+\cot\alpha>2$}
  \loigiai{
  \begin{enumerate}
   \item \textbf{Sai  .}Vì $0<\alpha<\dfrac{\pi}{2}$ nên $\sin\alpha>0$.
   \item \textbf{Đúng.} Vì $0<\alpha<\dfrac{\pi}{2}$ nên $\sin\alpha>0,\cos\alpha>0\Rightarrow\tan\alpha>0$. 
   \item \textbf{Sai.}  Vì
   \[\cos\alpha=\dfrac{3}{5}, 0<\alpha<\dfrac{\pi}{2}  \Rightarrow\sin\alpha=\sqrt{1-\cos^2 \alpha}=\dfrac{4}{5}\]
   \[\Rightarrow P=\sin^2 \alpha-2  \cos^2 \alpha=\dfrac{16}{25}-\dfrac{18}{25}=-\dfrac{2}{25}.\]
    \item \textbf{Đúng.} $Q=\tan\alpha+\cot\alpha=\dfrac{4}{3}+\dfrac{3}{4}=\dfrac{25}{12}>2$.
  \end{enumerate}
 } 
\end{ex}

% nguồn: Ok_dungsai_lop11_Lan2
\begin{ex}%[1D1V2-2]%[ma:7365f64c]
	Cho góc $\alpha$ thỏa mãn $\dfrac{-\pi}{2} < \alpha  < 0$ và $\cot \alpha=-3$.
	 \choiceTF
		{$\sin \alpha > 0$} 
		{\True $\sin \alpha=-\dfrac{\sqrt{10}}{10}$} 
		{\True $\dfrac{\cos\alpha-\sin \alpha}{\cos^3\alpha+3\sin  ^3\alpha+2\cos\alpha}=\dfrac{10}{21}$}
		{\True $\left[\tan \dfrac{17\pi}{4}+\tan \left(\dfrac{7\pi}{2}-\alpha \right)\right]^2+\left[\cot \dfrac{13\pi}{4}+\cot  \left(7\pi-\alpha \right)\right]^2=20$}
	\loigiai{ 
		\begin{itemchoice}
			\itemch  \textbf{Sai}.\\
			Do $\dfrac{-\pi}{2} <  \alpha < 0$ nên $\sin \alpha < 0$.
			\itemch  \textbf{Đúng}.\\
			Vì $\dfrac{-\pi}{2} < \alpha < 0$ nên $\sin \alpha < 0$.\\ 
			Từ hệ thức $1+\cot  ^2\alpha=\dfrac{1}{\sin ^2\alpha}$.\\
			Suy ra $\sin \alpha=-\sqrt{\dfrac{1}{1+\cot  ^2\alpha}}=-\dfrac{\sqrt{10}}{10}$.
			\itemch \textbf{Đúng}.\\ 
			Xét biểu thức  $P=\dfrac{\cos\alpha-\sin \alpha}{\cos^3\alpha+3\sin ^3\alpha+2\cos\alpha}$.
			Chia cả tử và mẫu cho $\sin ^3\alpha$  ta được\\
			\[P=\dfrac{\dfrac{\cos\alpha}{\sin ^3\alpha}-\dfrac{1}{\sin ^2\alpha}}{\dfrac{\cos^3\alpha}{\sin ^3\alpha}+3+2\dfrac{\cos\alpha}{\sin  ^3\alpha}}=\dfrac{\cot \alpha.\left(1+\cot ^2\alpha \right)-\left(1+\cot ^2\alpha \right)}{\cot ^3\alpha+3+2\cdot \cot \alpha\cdot \left(1+\cot ^2\alpha \right)}.\]
			Thay $\cot \alpha=-3$ vào  $P$ ta được $P=\dfrac{10}{21}$.
			\itemch \textbf{Đúng}.\\ 
			 Ta có
			\begin{itemize}
				 \item $\tan \dfrac{17\pi}{4}=\tan \left(\dfrac{\pi}{4}+4\pi \right)=\tan \dfrac{\pi}{4}=1$.
				\item $\tan  \left(\dfrac{7\pi}{2}-\alpha \right)=\cot \alpha$.
				\item $\cot  \dfrac{13\pi}{4}=\cot \left(\dfrac{\pi}{4}+3\pi \right)=\cot \dfrac{\pi}{4}=1$.
				\item $\cot \left(7\pi-\alpha \right)=-\cot \alpha$. 
			\end{itemize}
			Suy  ra
			\begin{align*}
				 &\left[\tan \dfrac{17\pi}{4}+\tan \left(\dfrac{7\pi}{2}-\alpha \right)\right]^2+\left[\cot \dfrac{13\pi}{4}+\cot \left(7\pi-\alpha \right)\right]^2\\
				&=\left(1+\cot  x\right)^2+\left(1-\cot x\right)^2=2+2\cot ^2x=2+2\cdot \left(-3\right)^2=20.
			\end{align*}
		\end{itemchoice}
	} 
\end{ex}

% nguồn: Ok_dungsai_lop11_Lan2
\begin{ex}%[1D1V2-2]%[ma:4b1f7b37]
	Cho $\tan  \alpha=2$.
	\choiceTF 
		{\True $\cot  \alpha=\dfrac{1}{2}$}
		 {Khi $0\le \alpha \le \dfrac{\pi}{2}$ thì $\cos \alpha=-\dfrac{\sqrt{5}}{5}$}
		{\True Giá trị của biểu thức $B=\dfrac{\sin \alpha-\cos \alpha}{\sin ^3\alpha+3\cos  ^3\alpha+2\sin \alpha}=\dfrac{5}{31}$}
		{Cho $\dfrac{\sin ^4\alpha}{a}+\dfrac{\cos ^4\alpha}{b}=\dfrac{1}{a+b}$. Giá trị của biểu thức $A=\dfrac{\sin  ^8\alpha}{a^3}+\dfrac{\cos ^8\alpha}{b^3}=\dfrac{1}{a^3+b^3}$}
	\loigiai{ 
		\begin{itemchoice}
			\itemch  \textbf{Đúng}.\\
			Ta có $\tan  \alpha\cdot \cot \alpha=1\Rightarrow \cot \alpha=\dfrac{1}{\tan \alpha}=\dfrac{1}{2}$.
			\itemch  \textbf{Sai}.\\
			Ta  có $1+\tan ^2\alpha=\dfrac{1}{\cos ^2\alpha} \Rightarrow \cos ^2\alpha=\dfrac{1}{1+\tan ^2\alpha}=\dfrac{1}{5} \Rightarrow \cos \alpha=\pm \dfrac{\sqrt{5}}{5}$.\\
			Mà $0\le \alpha \le \dfrac{\pi}{2}$ nên $\cos \alpha=\dfrac{\sqrt{5}}{5}$. 
			\itemch  \textbf{Đúng}.\\
			 Chia tử và mẫu cho $\cos^3 x$, ta được
			\[B=\dfrac{\dfrac{\sin \alpha}{\cos ^3\alpha}-\dfrac{\cos \alpha}{\cos ^3\alpha}}{\dfrac{\sin ^3\alpha}{\cos  ^3\alpha}+\dfrac{3\cos ^3\alpha}{\cos ^3\alpha}+\dfrac{2\sin \alpha}{\cos ^3\alpha}}=\dfrac{\tan \alpha \left(\tan ^2\alpha+1\right)-\left(\tan ^2\alpha+1\right)}{\tan ^3\alpha+3+2\tan \alpha \left(\tan ^2\alpha+1\right)}.\]
			Suy ra $B=\dfrac{2\left(4+1\right)-\left(4+1\right)}{8+3+2.2\left(4+1\right)}=\dfrac{5}{31}$. 
			 \itemch \textbf{Sai}.\\
			Đặt $\cos  ^2\alpha=t$.\\
			Ta  có
			\begin{align*}
				 &\dfrac{\left(1-t\right)^2}{a}+\dfrac{t^2}{b}=\dfrac{1}{a+b}\\
				 \Leftrightarrow& b\left(1-t\right)^2+at^2=\dfrac{ab}{a+b}\\
				 \Leftrightarrow& at^2+bt^2-2bt+b=\dfrac{ab}{a+b}\\
				 \Leftrightarrow& \left(a+b\right)t^2-2bt+b=\dfrac{ab}{a+b}\\
				 \Leftrightarrow& \left(a+b\right)^2t^2-2b\left(a+b\right)t+b^2=0\\
				 \Leftrightarrow& t=\dfrac{b}{a+b}.
			\end{align*}
			Suy  ra $\cos ^2\alpha=\dfrac{b}{a+b}$ và $\sin ^2\alpha=\dfrac{a}{a+b}$.\\
			Vậy $\dfrac{\sin ^8\alpha}{a^3}+\dfrac{\cos ^8\alpha}{b^3}=\dfrac{a}{\left(a+b\right)^4}+\dfrac{b}{\left(a+b\right)^4}=\dfrac{1}{\left(a+b\right)^3}$. 
		\end{itemchoice}
	 }
\end{ex}


%%%%%%%%%%%%%%%% TRẢ LỜI NGẮN %%%%%%%%%%%%%%%%

% nguồn: Ok_TraLoiNgan_Lop11_Lan2
\begin{ex}%[1D1H2-2]%[ma:e86ee221]
	Biểu thức $A=\tan  \left(\dfrac{17\pi}{2}-x\right)+2\cot (5\pi+x)=k \cot x$, khi đó $k$ bằng
	 \shortans{3}
	\loigiai{ 
		Ta có $A=\tan \left(8\pi+\dfrac{\pi}{2}-x\right)+2\cot x=\tan \left(\dfrac{\pi}{2}-x\right)+2\cot x=\cot  x+2\cot x=3\cot x$.\\
		Vậy  $k=1$.
	 }
\end{ex}

% nguồn: Ok_TraLoiNgan_Lop11_Lan2
\begin{ex}%[1D1H2-2]%[ma:4500efd1]
	Cho góc $\alpha $ thỏa mãn $\sin \alpha =\dfrac{12}{13} $ và $\dfrac{\pi }{2} <\alpha <\pi $. Tính $\cos \alpha$  (\textit{kết quả làm tròn đến hàng phần chục}).
	\shortans{-0,4} 
	\loigiai{ 
		Ta có $\heva{&\cos \alpha =\pm \sqrt{1-\sin ^{2} \alpha } =\pm \dfrac{5}{13}\\ &\dfrac{\pi }{2}  <\alpha <\pi \Rightarrow \cos\alpha <0}$.\\
		Suy ra $\cos \alpha =-\dfrac{5}{13} \approx -0{,}4$. 
	 }
\end{ex}

% nguồn: Ok_TraLoiNgan_Lop11_Lan2
\begin{ex}%[1D1H2-2]%[ma:806f106b]
	Cho $\cos \alpha=\dfrac{3}{4}$. Tính giá trị của biểu thức $B=\dfrac{8\tan  \alpha+24\cot \alpha}{\tan \alpha+\cot \alpha}$.
	\shortans{17} 
	\loigiai{ 
		 Ta có $\cos \alpha=\dfrac{3}{4} \Rightarrow 1+\tan ^2\alpha=\dfrac{1}{\cos ^2\alpha}=\dfrac{16}{9} \Rightarrow \tan ^2\alpha=\dfrac{7}{9}$.\\
		Khi  đó $B=\dfrac{8\tan \alpha+\dfrac{24}{\tan \alpha}}{\tan \alpha+\dfrac{1}{\tan \alpha}}=\dfrac{8\tan ^2\alpha+24}{\tan ^2\alpha+1}=\dfrac{8\cdot \dfrac{7}{9}+24}{\dfrac{7}{9}+1}=17$.
	 }
\end{ex}

% nguồn: Ok_TraLoiNgan_Lop11_Lan2
\begin{ex}%[1D1H2-2]%[ma:7ad8be4b]
	Cho góc $\alpha $ thỏa mãn $\cos \alpha =-\dfrac{\sqrt{5} }{3} $ và $\pi  <\alpha <\dfrac{3\pi }{2} $. Tính $\tan \alpha$ (\textit{kết quả làm tròn đến hàng phần trăm}).
	 \shortans{0,89}
	\loigiai{ 
		Ta có $\heva{&\sin \alpha =\pm \sqrt{1-\cos ^{2} \alpha  } =\pm \dfrac{2}{3}\\ &\pi <\alpha <\dfrac{3\pi }{2}\Rightarrow \sin\alpha <0 }$.\\
		Suy ra $\sin\alpha =-\dfrac{2}{3}  \Rightarrow\tan \alpha =\dfrac{\sin\alpha }{\cos \alpha } =\dfrac{2}{\sqrt{5}} \approx 0{,}89$.
	 }
\end{ex}

% nguồn: Ok_TraLoiNgan_Lop11_Lan2
\begin{ex}%[1D1H2-2]%[ma:3d66b3ac]
	Cho tam giác $ABC$, khi đó biểu thức $T=\dfrac{\sin ^3\dfrac{B}{2}}{\cos \left(\dfrac{A+C}{2}\right)}+\dfrac{\cos  ^3\dfrac{B}{2}}{\sin \left(\dfrac{A+C}{2}\right)}-\dfrac{\cos (A+C)}{\sin B} \tan B$ bằng
	 \shortans{2}
	 \loigiai{
		 Vì $A+B+C=180^{\circ}$ nên $A+C=180^{\circ}-B$.\\
		 Khi đó\allowdisplaybreaks
		\begin{eqnarray*}
			T&=&\dfrac{\sin ^3\dfrac{B}{2}}{\cos \left(\dfrac{180^{\circ}-B}{2}\right)}+\dfrac{\cos ^3\dfrac{B}{2}}{\sin \left(\dfrac{180^{\circ}-B}{2}\right)}-\dfrac{\cos \left(180^{\circ}-B\right)}{\sin B} \cdot  \tan B\\
			&=&\dfrac{\sin ^3\dfrac{B}{2}}{\sin \dfrac{B}{2}}+\dfrac{\cos ^3\dfrac{B}{2}}{\cos \dfrac{B}{2}}-\dfrac{-\cos B}{\sin B} \cdot  \tan B\\
			&=&\sin  ^2\dfrac{B}{2}+\cos ^2\dfrac{B}{2}+\cot B\cdot \tan B=1+1=2.
		\end{eqnarray*}
	} 
\end{ex}

% nguồn: Ok_TraLoiNgan_Lop11_Lan2
\begin{ex}%[1D1H2-2]%[ma:d6645e22]
	 Cho góc $\alpha$ thỏa mãn $\tan \alpha=2$. Giá trị biểu thức $P=\dfrac{2\sin ^2\alpha+3\sin \alpha\cdot \cos \alpha+4\cos ^2\alpha}{5\sin ^2\alpha+6\cos ^2\alpha}$.(Kết quả làm tròn đến hàng phần trăm)
	\shortans{0,69} 
	 \loigiai{
		Chia cả tử và mẫu của $P$ cho $\cos ^2\alpha$ ta  được
		\[P=\dfrac{2\tan ^2\alpha+3\tan \alpha+4}{5\tan ^2\alpha+6}=\dfrac{2\cdot 2^2+3\cdot  2+4}{5\cdot 2^2+6}=\dfrac{9}{13} \approx 0{,}69.\]
	 }
\end{ex}

% nguồn: Ok_TraLoiNgan_Lop11_Lan2
\begin{ex}%[1D1H2-2]%[ma:6d64b61e]
	Biểu thức sau: $S=\dfrac{\sin \left(\dfrac{15\pi}{2}-x\right)-2\cos (x-\pi)}{\cos \left(\dfrac{5\pi}{2}-x\right)}=k \cot x$.  Khi đó $k$ bằng
	 \shortans{1}
	 \loigiai{
		Ta  có
		\allowdisplaybreaks 
		\begin{eqnarray*}
			S  &=&\dfrac{\sin \left(7\pi+\dfrac{\pi}{2}-x\right)-2\cos (\pi-x)}{\cos \left(2\pi+\dfrac{\pi}{2}-x\right)}
			=\dfrac{\sin \left(\pi+\dfrac{\pi}{2}-x\right)+2\cos x}{\cos \left(\dfrac{\pi}{2}-x\right)}\\ 
			&=&\dfrac{-\sin \left(\dfrac{\pi}{2}-x\right)+2\cos  x}{\sin x}
			=\dfrac{-\cos x+2\cos x}{\sin x}\\ 
			&=&\dfrac{\cos x}{\sin  x}=\cot x.
		\end{eqnarray*}
	} 
\end{ex}

% nguồn: Ok_TraLoiNgan_Lop11_Lan2
\begin{ex}%[1D1H2-2]%[ma:d87ec0ab]
	Cho góc $\alpha $ thỏa mãn $\cos \alpha =-\dfrac{12}{13} $  và $\dfrac{\pi }{2} <\alpha <\pi .$ Tính $\tan \alpha$ (\textit{kết quả làm tròn đến hàng phần chục}).
	\shortans{-0,4} 
	 \loigiai{
		Ta có $\heva{&\sin \alpha =\pm \sqrt{1-\cos ^{2}  \alpha } =\pm \dfrac{5}{13}\\ &\dfrac{\pi }{2} <\alpha <\pi \Rightarrow\sin \alpha >0.}$\\
		Suy ra $\sin \alpha =\dfrac{5}{13}  \Rightarrow\tan \alpha =\dfrac{\sin \alpha }{\cos \alpha } =-\dfrac{5}{12} \approx -0{,}4$.
	 }
\end{ex}

% nguồn: Ok_TraLoiNgan_Lop11_Lan2
\begin{ex}%[1D1H2-2]%[ma:f47c384b]
	Cho góc $\alpha$ thỏa mãn $\dfrac{\pi}{2} < \alpha < \pi$ và $\tan \alpha-\cot \alpha=1$. Tính  $P=\tan \alpha+\cot \alpha$. (Kết quả làm tròn đến hàng phần mười)
	 \shortans{-2,2}
	 \loigiai{
		Ta có 
		\begin{align*}
			&\tan \alpha-\cot  \alpha=1\\
			\Leftrightarrow&  \tan \alpha-\dfrac{1}{\tan \alpha}=1\\
			 \Leftrightarrow& \tan ^2\alpha-\tan \alpha-1=0\\
			\Leftrightarrow&  \tan \alpha=\dfrac{1\pm \sqrt{5}}{2}.
		\end{align*}
		 Do $\dfrac{\pi}{2} < \alpha < \pi$ suy ra $\tan \alpha < 0$ nên $\tan \alpha=\dfrac{1-\sqrt{5}}{2} \Rightarrow \cot \alpha=\dfrac{1}{\tan \alpha}=\dfrac{2}{1-\sqrt{5}}$.\\
		Thay $\tan \alpha=\dfrac{1-\sqrt{5}}{2}$ và $\cot \alpha=\dfrac{2}{1-\sqrt{5}}$ vào $P$ ta được \[P=\dfrac{1-\sqrt{5}}{2}+\dfrac{2}{1-\sqrt{5}}=-\sqrt{5}  \approx-2{,}2.\]
	 }
\end{ex}

% nguồn: Ok_TraLoiNgan_Lop11_Lan2
\begin{ex}%[1D1H2-2]%[ma:d663b30f]
	Cho góc $\alpha $ thỏa mãn $\tan \alpha =-\dfrac{4}{3} $ và $\dfrac{2017\pi }{2} <\alpha <\dfrac{2019\pi }{2} $. Tính $\sin  \alpha$.
	\shortans{0,8} 
	\loigiai{ 
		Ta có 
		\begin{itemize}
			\item $1+\tan ^{2} \alpha  =\dfrac{1}{\cos ^{2} \alpha } \Rightarrow1+\left(-\dfrac{4}{3} \right)^{2} =\dfrac{1}{\cos ^{2} \alpha} \Rightarrow \cos \alpha =\pm\dfrac{3}{5}$.
			\item $\dfrac{2017\pi }{2} <\alpha <\dfrac{2019\pi  }{2} \Rightarrow \dfrac{\pi }{2} +504\cdot2\pi <\alpha <\dfrac{3\pi }{2} +504\cdot 2\pi\Rightarrow \cos\alpha<0 $.
		\end{itemize}
		Suy ra $\cos \alpha =-\dfrac{3}{5} $. Mà $\tan \alpha =\dfrac{\sin  \alpha }{\cos \alpha } \Leftrightarrow -\dfrac{4}{3} =\dfrac{\sin \alpha }{-\dfrac{3}{5}} \Rightarrow\sin \alpha =\dfrac{4}{5} $.
	 }
\end{ex}

% nguồn: Ok_TraLoiNgan_Lop11_Lan2
\begin{ex}%[1D1H2-2]%[ma:433858b7]
	Biểu thức sau: $T=2\sin \left(\dfrac{9\pi}{2}-x\right)+3\cos (19\pi-x)=k \cos x$. Khi đó giá trị $k$ bằng 
	 \shortans{-1}
	\loigiai{ 
		Ta có 
		\allowdisplaybreaks 
		\begin{eqnarray*}
			T &=& 2\sin \left(4\pi+\dfrac{\pi}{2}-x\right)+3\cos  (18\pi+\pi-x)\\
			&=& 2\sin  \left(\dfrac{\pi}{2}-x\right)+3\cos (\pi-x)\\
			&=&  2\cos x-3\cos x\\
			 &=&-\cos x.
		\end{eqnarray*}
		Vậy  $k=-1$.
	} 
\end{ex}

% nguồn: Ok_TraLoiNgan_Lop11_Lan2
\begin{ex}%[1D1H2-2]%[ma:a20a72d3]
	Cho $\sin\alpha=-\dfrac{1}{3}$, với $\left(180^\circ<\alpha <270^\circ\right)$.  Tính $\cos\alpha$ (Kết quả làm tròn đến hàng phần trăm).
	 \shortans{0,94}
	\loigiai{ 
		Vì $270^\circ<\alpha <360^\circ$ nên $\cos\alpha>0\Rightarrow\cos  \alpha=\sqrt{1-\sin^2\alpha}=\sqrt{\dfrac{8}{9}}$.\\
		 Do đó $\cos\alpha=\dfrac{2\sqrt 2}{3}\approx 0{,}94$.
	} 
\end{ex}

% nguồn: Ok_TraLoiNgan_Lop11_Lan2
\begin{ex}%[1D1H2-2]%[ma:ca47a7fb]
	 Cho $\cos \alpha =\dfrac{1}{3} $. Khi đó $\sin \left(\alpha -\dfrac{3\pi }{2} \right)$ bằng bao nhiêu? (\textit{kết quả làm tròn đến hàng phần trăm}).
	\shortans{0,33} 
	\loigiai{ 
		Ta có $\sin \left(\alpha -\dfrac{3\pi }{2}  \right)=\sin \left(\alpha +\dfrac{\pi }{2} -2\pi \right)=\sin \left(\alpha +\dfrac{\pi }{2} \right)=\cos \alpha =\dfrac{1}{3}$.
	 }
\end{ex}

% nguồn: Ok_TraLoiNgan_Lop11_Lan2
\begin{ex}%[1D1H2-2]%[ma:a6163f5a]
	Tính giá trị của $\cos \left[\dfrac{\pi  }{3} +(2k+1)\pi \right]$.
	\shortans{-0,5} 
	 \loigiai{
		Ta có $\cos \left[\dfrac{\pi }{3} +(2k+1)\pi \right]=\cos \left(\dfrac{\pi }{3} +\pi +k2\pi \right)=\cos \left(\dfrac{\pi }{3} +\pi \right)=-\cos  \dfrac{\pi }{3} =-\dfrac{1}{2}$.
	} 
\end{ex}

% nguồn: Ok_TraLoiNgan_Lop11_Lan2
\begin{ex}%[1D1H2-2]%[ma:a6a98a78]
	Biểu thức $B=\dfrac{2\sin (x-4\pi)+\cos \left(x-\dfrac{5\pi}{2}\right)}{\sin  \left(\dfrac{3\pi}{2}-x\right)}=k \tan x$, khi đó $k$ bằng
	\shortans{-3} 
	\loigiai{ 
		Ta  có
		 \allowdisplaybreaks
		\begin{eqnarray*}
			B&=&\dfrac{2 \sin x+\cos \left(x-\dfrac{\pi}{2}-2\pi\right)}{\sin \left(\pi+\dfrac{\pi}{2}-x\right)}=\dfrac{2 \sin x+\cos \left(x-\dfrac{\pi}{2}\right)}{-\sin \left(\dfrac{\pi}{2}-x\right)}=\dfrac{2 \sin x+\cos \left(\dfrac{\pi}{2}-x\right)}{-\cos x} \\ 
			&=&\dfrac{2 \sin x+\sin x}{-\cos x}=-3 \tan  x.
		\end{eqnarray*}
		Vậy $k=-3$. 
	} 
\end{ex}

% nguồn: Ok_TraLoiNgan_Lop11_Lan2
\begin{ex}%[1D1H2-2]%[ma:ea5a03f0]
	Cho góc $\alpha$ thỏa mãn $\pi < \alpha < \dfrac{3\pi}{2}$ và $\sin \alpha-2\cos \alpha=1$.  Tính $P=2\tan \alpha-\cot \alpha$.
	(Kết quả làm tròn  đến hàng phần trăm)
	\shortans{0,17} 
	\loigiai{ 
		Với $\pi < \alpha < \dfrac{3\pi}{2}$ suy ra  $\heva{&\sin \alpha < 0\\&\cos \alpha < 0.}$\\
		Ta có $\heva{&\sin  \alpha-2\cos \alpha=1\\&\sin ^2\alpha+\cos ^2\alpha=1.}$\\
		\begin{align*}
			&\left(1+2\cos \alpha \right)^2+\cos  ^2\alpha=1\\
			\Leftrightarrow& 5\cos ^2\alpha+4\cos \alpha=0\\ 
			\Leftrightarrow&  \hoac{&\cos \alpha=0\,\text{(loại)}\\&\cos \alpha=-\dfrac{4}{5}.}
		\end{align*}
		Từ hệ thức $\sin  ^2\alpha+\cos ^2\alpha=1$.
		\begin{align*}
			 &\left(-\dfrac{4}{5}\right)^2+\sin^2x=1\\
			\Leftrightarrow& \hoac{&\sin \alpha=-\dfrac{3}{5}\\&\sin \alpha=\dfrac{3}{5}.} 
		\end{align*}
		Do $\sin \alpha  < 0$ nên $\sin \alpha=-\dfrac{3}{5}$.\\
		Ta có $\tan  \alpha=\dfrac{\sin \alpha}{\cos \alpha}=\dfrac{3}{4}$
		và $\cot \alpha=\dfrac{\cos  \alpha}{\sin \alpha}=\dfrac{4}{3}$.\\
		Thay $\tan \alpha=\dfrac{3}{4}$ và $\cot \alpha=\dfrac{4}{3}$ vào $P$, ta được $P=\dfrac{1}{6}  \approx 0{,}17$.
	} 
\end{ex}

% nguồn: Ok_TraLoiNgan_Lop11_Lan2
\begin{ex}%[1D1H2-2]%[ma:9664c379]
	Cho $\sin\alpha=\dfrac{3}{5}$ và  $\left(90^\circ<\alpha <180^\circ\right)$. Tính $\cos\alpha$.
	 \shortans{-0,8}
	\loigiai{ 
		$90^\circ<\alpha <180^\circ$ nên  $\cos\alpha < 0$.\\
		Ta  có $\sin ^2\alpha+\cos ^2\alpha=1\Rightarrow{\cos ^2}\alpha=1-\sin ^2\alpha=1-\left(\dfrac{3}{5}\right)^2=\dfrac{16}{25}\Rightarrow\cos\alpha=-\dfrac{4}{5}$.\\
		 Vậy $\cos\alpha=-\dfrac{4}{5}=-0{,}8$.
	 }
\end{ex}

% nguồn: Ok_TraLoiNgan_Lop11_Lan2
\begin{ex}%[1D1V2-2]%[ma:61916b4f]
	Cho góc $\alpha$ thỏa mãn $\sin \alpha+\cos \alpha=\dfrac{\sqrt{2}}{2}$. Tính  $P=\tan ^2\alpha+\cot ^2\alpha$.
	\shortans{14} 
	 \loigiai{
		 Ta có
		\begin{align*}
			&\sin \alpha+\cos \alpha=\dfrac{\sqrt{2}}{2}\\ 
			\Rightarrow& \left(\sin  \alpha+\cos \alpha \right)^2=\dfrac{1}{2}\\
			\Leftrightarrow& \sin \alpha  \cdot \cos \alpha=-\dfrac{1}{4}.
		\end{align*}
		Khi  đó
		\begin{align*}
			P&=\dfrac{\sin ^2\alpha}{\cos ^2\alpha}+\dfrac{\cos  ^2\alpha}{\sin ^2\alpha}\\
			&=\dfrac{\sin ^4\alpha+\cos  ^4\alpha}{\sin ^2\alpha\cdot \cos ^2\alpha}\\ &=\dfrac{\left(\sin ^2\alpha+\cos ^2\alpha \right)^2-2\sin ^2\alpha\cdot \cos ^2\alpha}{\sin ^2\alpha\cdot \cos ^2\alpha}\\
			&=\dfrac{1-2\left(\sin \alpha  \cdot \cos \alpha \right)^2}{\left(\sin \alpha \cdot \cos \alpha \right)^2}=14.
		\end{align*}
	 }
\end{ex}

% nguồn: Ok_TraLoiNgan_Lop11_Lan2
\begin{ex}%[1D1V2-2]%[ma:b9963a37]
	Cho góc $\alpha$ thỏa mãn $\sin \alpha \cdot \cos \alpha=\dfrac{12}{25}$ và $\sin \alpha+\cos \alpha > 0$. Tính $P=\sin ^3\alpha+\cos ^3\alpha$. (Kết quả làm  tròn đến hàng phần trăm)
	 \shortans{0,73}
	\loigiai{ 
		Ta  có
		\[P=\sin ^3\alpha+\cos ^3\alpha=\left(\sin \alpha+\cos \alpha  \right)^3-3\sin \alpha \cdot \cos \alpha \left(\sin \alpha+\cos \alpha \right).\]
		Với 
		\begin{align*}
			&\left(\sin \alpha+\cos \alpha \right)^2=\sin  ^2\alpha+2\sin \alpha \cos \alpha+\cos ^2\alpha=1+\dfrac{24}{25}=\dfrac{49}{25}\\
			\Leftrightarrow& \hoac{&\sin  \alpha+\cos \alpha=\dfrac{7}{5}\\&\sin \alpha+\cos \alpha=-\dfrac{7}{5}.}
		\end{align*}
		 Vì $\sin \alpha+\cos \alpha > 0$ nên ta chọn $\sin \alpha+\cos \alpha=\dfrac{7}{5}$.\\
		Thay $\heva{&\sin \alpha+\cos \alpha=\dfrac{7}{5}\\&\sin \alpha \cos \alpha=\dfrac{12}{25}}$ vào $P$, ta được 
		\[P=\left(\dfrac{7}{5} \right)^3-3\cdot \dfrac{12}{25}\cdot \dfrac{7}{5}=\dfrac{91}{125}  \approx 0{,}73.\]
	 }
\end{ex}


%%%%%%%%%%%%%%%% TỰ LUẬN %%%%%%%%%%%%%%%%

% nguồn: Ok_TuanLuan_Toan11_Lan2
\begin{ex}%[1D1H2-2]%[ma:2a609cdf]
	Giá trị lớn  nhất của $Q=\sin ^6x+\cos ^6x$ bằng
	 \loigiai{
		Ta có $Q=\sin ^6x+\cos  ^6x=1-\dfrac{3}{4} \sin ^22x$.\\
		Vì $0\le \sin ^22x\le 1$, $\forall x\in\mathbb{R}$ nên $-\dfrac{3}{4} \le-\dfrac{3}{4}  \sin ^22x\le 0\Leftrightarrow \dfrac{1}{4} \le 1-\dfrac{3}{4} \sin ^22x\le 1$.\\
		 Vậy giá trị lớn nhất của $Q$ bằng $1$. Dấu "$=$" xảy ra khi $\sin 2x=0\Leftrightarrow x=k\dfrac{\pi}{2}$ ($k\in\mathbb{Z}$).
	} 
\end{ex}

% nguồn: Ok_TuanLuan_Toan11_Lan2
\begin{ex}%[1D1H2-2]%[ma:3c29c62c]
	Tính giá trị lượng  giác còn lại của góc $x$ biết:
	\begin{enumerate}
		\item  $\sin x=-\dfrac{3}{5}$ với $\pi < x < \dfrac{3\pi}{2}$.
		\item $\cos x=\dfrac{1}{4}$ với $0< x < \dfrac{\pi}{2}$. 
		 \item $\cos x=\dfrac{3}{5}$ với $0< x < 90^\circ$.
		\item $\cos x=-\dfrac{5}{13}$ với $180^\circ< x  < 270^\circ$.
	\end{enumerate}
	\loigiai{ 
		\begin{enumerate}
			\item Do  $\pi < x < \dfrac{3\pi}{2} \Rightarrow \heva{&\sin x < 0\\
			&\cos x  < 0 \\
			 &\tan x > 0 \\
			&\cot x  > 0.}$
			Từ đó với $\sin x=-\dfrac{3}{5} \Rightarrow \cos x=-\sqrt{1-\sin ^2x}=-\dfrac{4}{5} \Rightarrow \heva{&\tan  x=\dfrac{\sin x}{\cos x}=\dfrac{3}{4} \\
			&\cot x=\dfrac{\cos  x}{\sin x}=\dfrac{4}{3}.}$
			\item Do $0< x < \dfrac{\pi}{2} \Rightarrow \heva{&\sin x  > 0 \\
			&\cos x > 0  \\
			 &\tan x > 0 \\
			&\cot  x > 0.
			}$ 
			Từ đó với  $\cos x=\dfrac{1}{4} \Rightarrow \sin x=\sqrt{1-\cos ^2x}=\dfrac{\sqrt{15}}{4} \Rightarrow \heva{&\tan x=\dfrac{\sin x}{\cos x}=\sqrt{15} \\
			&\cot x=\dfrac{\cos  x}{\sin x}=\dfrac{1}{\sqrt{15}}.
			 }$
			 \item Do $0< x < 90^\circ\Rightarrow \heva{&\sin x > 0 \\
			&\cos  x > 0 \\
			&\tan x > 0 \\ 
			&\cot x > 0.}$ 
			Từ đó với $\cos  x=\dfrac{3}{5} \Rightarrow \sin x=\sqrt{1-\cos ^2x}=\dfrac{4}{5} \Rightarrow \heva{&\tan x=\dfrac{\sin x}{\cos x}=\dfrac{4}{3} \\
			&\cot x=\dfrac{\cos  x}{\sin x}=\dfrac{3}{4}.
			 }$
			\item Do $180^\circ< x < 270^\circ\Rightarrow \heva{&\sin x  < 0 \\
			&\cos x  < 0\\
			&\tan x >  0\\
			&\cot x >  0.}$
			Từ  đó với $\cos x=-\dfrac{5}{13} \Rightarrow \sin x=-\sqrt{1-\cos ^2x}=-\dfrac{12}{13} \Rightarrow \heva{&\tan x=\dfrac{\sin x}{\cos x}=\dfrac{12}{5} \\
			&\cot x=\dfrac{\cos x}{\sin x}=\dfrac{5}{12}. 
			 }$
		\end{enumerate}
	 }
\end{ex}

% nguồn: Ok_TuanLuan_Toan11_Lan2
\begin{ex}%[1D1H2-2]%[ma:e2b9707f]
	Tính  các giá trị lượng giác của mỗi góc sau:
	\begin{enumerate}
		\item $\dfrac{\pi}{3}  +k2\pi \,\, (k\in \mathbb{Z})$;
		\item $\dfrac{\pi}{2} +k\pi \, \,  (k\in \mathbb{Z})$.
	\end{enumerate}
	\loigiai{ 
		\begin{enumerate}
			\item Ta có 
			\begin{itemize}
				\item $\sin \left(\dfrac{\pi}{3} +k2\pi \right)=\sin \dfrac{\pi}{3}  =\dfrac{\sqrt{3}}{2}$.
				\item  $\cos \left(\dfrac{\pi}{3} +k2\pi \right)=\cos \dfrac{\pi}{3} =\dfrac{1}{2}$.
				 \item $\tan \left(\dfrac{\pi}{3} +k2\pi \right)=\tan \dfrac{\pi}{3} =\sqrt{3}$.
				\item $\cot \left(\dfrac{\pi}{3} +k2\pi \right)=\cot  \dfrac{\pi}{3} =\dfrac{1}{\sqrt{3}} $.
			\end{itemize}
			 \item $\tan \left(\dfrac{\pi}{2} +k\pi \right)$ không xác định; $\cot \left(\dfrac{\pi}{2} +k\pi \right)=\cot \dfrac{\pi}{2} =0$.
			\begin{itemize}
				\item Với  $k=2l\, (l\in \mathbb{Z})$:\\
				$\sin \left(\dfrac{\pi}{2} +k\pi \right)=\sin \left(\dfrac{\pi}{2} +l2\pi \right)=\sin \dfrac{\pi}{2} =1$;  \\
				$\cos \left(\dfrac{\pi}{2} +k\pi \right)=\cos \left(\dfrac{\pi}{2} +l2\pi \right)=\cos  \dfrac{\pi}{2} =0$.
				\item Với $k=2l+1\, \,  (l\in \mathbb{Z})$:\\
				$\sin \left(\dfrac{\pi}{2} +k\pi \right)=\sin  \left(\dfrac{\pi}{2} +\pi +l2\pi \right)=\sin \dfrac{3\pi}{2} =-1$;\\
				$\cos \left(\dfrac{\pi}{2} +k\pi \right)=\cos  \left(\dfrac{\pi}{2} +\pi +l2\pi \right)=\cos \dfrac{3\pi}{2} =0$.
			\end{itemize}
	\end{enumerate}}
\end{ex}

% nguồn: Ok_TuanLuan_Toan11_Lan2
\begin{ex}%[1D1H2-2]%[ma:303277ab]
	Tính giá trị lượng  giác còn lại của góc $x$ biết
	\begin{enumerate}
		\item $\tan x=3$ với $\pi < x < \dfrac{3\pi}{2}$. 
		\item $\tan x=-2$ với $\dfrac{\pi}{2} <  x < \pi$.
		\item $\tan x=-\dfrac{1}{2}$ với $\dfrac{\pi}{2} < x  < \pi$.
		\item $\cot x=3$ với $\pi <  x < \dfrac{3\pi}{2}$.
	\end{enumerate}
	 \loigiai{
		\begin{enumerate}
			 \item $\tan x=3\Rightarrow \cot x=\dfrac{1}{3}$.
			$\tan x=3\Rightarrow \dfrac{\sin x}{\cos x}=3\Rightarrow  \sin ^2x=9\cos ^2x\Rightarrow \sin ^2x-9\left(1-\sin ^2x\right)=0\Rightarrow \sin ^2x=\dfrac{9}{10}.$
			Vì $\pi < x < \dfrac{3\pi}{2}$ nên $\heva{&\sin x <  0 \\ &\cos x < 0.}$
			Do đó  $\sin x=-\dfrac{3\sqrt{10}}{10}$; $\cos x=-\dfrac{\sqrt{10}}{10}$.
			\item $\tan x=-2$ $\Rightarrow \cot  x=-\dfrac{1}{2}$
			$\tan  x=-2\Rightarrow \dfrac{\sin x}{\cos x}=-2\Rightarrow \sin ^2x=4\cos ^2x\Rightarrow \sin ^2x-4\left(1-\sin ^2x\right)=0\Rightarrow \sin ^2x=\dfrac{4}{5}.
			$ 
			Vì $\dfrac{\pi}{2} < x <  \pi$ nên $\heva{&\sin x > 0\\ &\cos x < 0.}$
			Do đó $\sin x=\dfrac{2\sqrt{5}}{5}$; $\cos x=-\dfrac{\sqrt{5}}{5}$. 
			\item $\tan x=-\dfrac{1}{2}$ $\Rightarrow \cot  x=-2$
			$\tan  x=-\dfrac{1}{2} \Rightarrow \dfrac{\sin x}{\cos x}=-\dfrac{1}{2} \Rightarrow 4\sin ^2x=\cos ^2x\Rightarrow 4\sin ^2x-1\left(1-\sin ^2x\right)=0\Rightarrow \sin ^2x=\dfrac{1}{5}.
			$ 
			Vì $\dfrac{\pi}{2} < x < \pi$ nên $\heva{&\sin x > 0 \\&\cos x < 0.}$ 
			Do đó $\sin x=\dfrac{\sqrt{5}}{5}$; $\cos x=-\dfrac{2\sqrt{5}}{5}$. 
			\item $\cot x=3$ $\Rightarrow  \tan x=\dfrac{1}{3}$
			$\tan x=\dfrac{1}{3}  \Rightarrow \dfrac{\sin x}{\cos x}=\dfrac{1}{3} \Rightarrow 9\sin ^2x=\cos ^2x\Rightarrow 9\sin ^2x-\left(1-\sin ^2x\right)=0\Rightarrow \sin ^2x=\dfrac{1}{10}.
			$ 
			Vì $\pi < x < \dfrac{3\pi}{2}$ nên $\heva{&\sin x < 0 \\ &\cos x <  0.
			 }$
			Do đó $\sin x=-\dfrac{\sqrt{10}}{10}$; $\cos x=-\dfrac{3\sqrt{10}}{10}$. 
		\end{enumerate}
	} 
\end{ex}

% nguồn: Ok_TuanLuan_Toan11_Lan2
\begin{ex}%[1D1H2-2]%[ma:c2f8d0ef]
	Giá trị lớn nhất của biểu thức $M=7\cos ^2x-2\sin ^2x$  là
	\loigiai{ 
		Ta có $M=7\left(1-\sin ^2x\right)-2\sin  ^2x=7-9\sin ^2x.$\\
		Lại có $0\le \sin ^2x\le 1$, $\forall x\in \mathbb{R}\Leftrightarrow 0\ge-9\sin ^2x\ge-9$, $\forall x\in \mathbb{R}\Leftrightarrow 7\ge 7-2\sin  ^2x\ge-2$, $\forall x\in \mathbb{R}$.\\
		Vậy giá trị lớn nhất của $M$ bằng $7$. Dấu "$=$" xảy ra khi và  chỉ khi $\sin x=0\Leftrightarrow x=k\pi$ ($k\in\mathbb{Z}$).
	 }
\end{ex}

% nguồn: Ok_TuanLuan_Toan11_Lan2
\begin{ex}%[1D1H2-2]%[ma:95991a52]
	Tính  các giá trị lượng giác của góc $\alpha $, biết
	\begin{enumerate}
		\item $\cos\alpha =\dfrac{1}{5} $ và $0<\alpha  <\dfrac{\pi}{2} $;
		\item $\sin\alpha  =\dfrac{2}{3} $ và $\dfrac{\pi}{2} <\alpha <\pi $;
		\item  $\tan\alpha =\sqrt{5} $ và $\pi <\alpha <\dfrac{3\pi}{2} $;
		\item $\cot\alpha =-\dfrac{1}{\sqrt{2}} $ và $\dfrac{3\pi}{2}  <\alpha <2\pi $.
	\end{enumerate}
	\loigiai{ 
		\begin{enumerate}
			\item Vì $0<a<\dfrac{\pi}{2} $  nên sin $a>0$.\\
			Mặt  khác, từ $\sin^2 a+\cos^2 a=1$ suy ra $\sin\alpha =\sqrt{1-\cos^2 \alpha} =\sqrt{1-\left(\dfrac{1}{5} \right)^{2}} =\dfrac{2\sqrt{6}}{5} $.\\
			Do đó $\tan\alpha =\dfrac{\sin\alpha}{\cos\alpha} =\dfrac{\dfrac{2\sqrt{6}}{5}}{\dfrac{1}{5}}  =2$ và $\cot\alpha =\dfrac{1}{\tan\alpha} =\dfrac{1}{2\sqrt{6}} =\dfrac{\sqrt{6}}{12} $.
			\item  Vì $\dfrac{\pi}{2} <\alpha <\pi $ nên $\cos\alpha <0$. \\
			Mặt khác từ $\sin^2 \alpha +\cos^2 \alpha =1$ suy ra $\cos\alpha  =-\sqrt{1-\sin^2 \alpha} =-\sqrt{1-\left(\dfrac{2}{3} \right)^{2}} =-\dfrac{\sqrt{5}}{3} $.\\
			Do đó $\tan\alpha =\dfrac{\sin\alpha}{\cos\alpha} =\dfrac{\dfrac{2}{3}}{-\dfrac{\sqrt{5}}{3}} =-\dfrac{2}{\sqrt{5}} =-\dfrac{2\sqrt{5}}{5} $ và $\cot\alpha  =\dfrac{1}{\tan\alpha} =\dfrac{1}{-\dfrac{2\sqrt{5}}{5}} =-\dfrac{\sqrt{5}}{2} $.
			\item Ta có $\cot\alpha =\dfrac{1}{\tan\alpha} =\dfrac{1}{\sqrt{5}} =\dfrac{\sqrt{5}}{5} $. Do  $\pi <\alpha <\dfrac{3\pi}{2} $ nên cos $\alpha <0$. \\
			Mặt  khác từ $1+\tan^2 \alpha =\dfrac{1}{\cos^2 \alpha} $ suy ra $\cos\alpha =-\sqrt{\dfrac{1}{1+\tan^2 \alpha}} =-\sqrt{\dfrac{1}{1+(\sqrt{5})^{2}}} =-\dfrac{\sqrt{6}}{6} $.\\
			 Mà ta có $\tan\alpha =\dfrac{\sin\alpha}{\cos\alpha} \Rightarrow \sin\alpha =\tan\alpha \cdot\cot\alpha =\sqrt{5} \cdot\left(-\dfrac{\sqrt{6}}{6} \right)=-\dfrac{\sqrt{30}}{6} $.
			\item Ta có $\tan\alpha =\dfrac{1}{\cot\alpha} =\dfrac{1}{-\dfrac{1}{\sqrt{2}}} =-\sqrt{2} $. Vì $\dfrac{3\pi}{2}  <\alpha <2\pi $ nên $\cos\alpha > 0$. \\
			Mặt khác từ $1+\tan^2 \alpha =\dfrac{1}{\cos^2 \alpha} $ suy ra $\cos\alpha =\sqrt{\dfrac{1}{1+\tan^2  \alpha}} =\sqrt{\dfrac{1}{1+(-\sqrt{2})^{2}}} =\dfrac{\sqrt{3}}{3} $.\\
			Mà $\tan\alpha =\dfrac{\sin\alpha}{\cos\alpha} \Rightarrow \sin\alpha  =\tan\alpha \cdot\cot\alpha =-\sqrt{2} \cdot\left(\dfrac{\sqrt{3}}{3} \right)=-\dfrac{\sqrt{6}}{3} $.
	\end{enumerate}}
\end{ex}

% nguồn: Ok_TuanLuan_Toan11_Lan2
\begin{ex}%[1D1H2-2]%[ma:b0695712]
	Cho $\cot x=\dfrac{3}{4}\left(\pi < x < \dfrac{3\pi}{2}\right)$.  Tính giá trị của các giá trị lượng giác còn lại.
	 \loigiai{
		Vì $\pi <  x < \dfrac{3\pi}{2}$ nên $\sin x < 0$.\\
		Ta có $\tan x \cdot \cot x=1\Rightarrow \tan x=\dfrac{1}{\cot x}=\dfrac{1}{\dfrac{3}{4}}=\dfrac{4}{3}$.\\ 
		Mà $\dfrac{1}{\sin  ^2x}=1+\cot ^2x=1+\left(\dfrac{3}{4}\right)^2=\dfrac{25}{16} \Rightarrow \sin ^2x=\dfrac{16}{25}$.\\
		Suy ra  $\sin x=-\dfrac{4}{5}$.\\
		Và $\cot x=\dfrac{\cos x}{\sin x} \Rightarrow \cos x=\cot x  \cdot \sin x=\dfrac{3}{4} \cdot\left(-\dfrac{4}{5}\right)=-\dfrac{3}{5}$.
	} 
\end{ex}

% nguồn: Ok_TuanLuan_Toan11_Lan2
\begin{ex}%[1D1H2-2]%[ma:df505a87]
	Tính  các giá trị lượng giác của góc $\alpha $, biết
	\begin{enumerate}
		\item $\sin \alpha =\dfrac{1}{3} $ và $-\dfrac{\pi}{2} <\alpha  <0$;
		\item $\cos \alpha =-0,7$ và $\pi <\alpha <\dfrac{3\pi}{2}  $;
		\item $\tan \alpha  =2$ và $0<\alpha <\dfrac{\pi}{2} $;
		\item $\cot \alpha =\dfrac{7}{3} $ và $\pi <\alpha <\dfrac{3\pi}{2}  $.
	\end{enumerate}
	 \loigiai{
		\begin{enumerate}
			\item Ta có $\cos ^{2} \alpha =1-\sin ^{2} \alpha =1-\left(\dfrac{1}{3} \right)^{2} =\dfrac{8}{9} $. Vì $-\dfrac{\pi}{2} <\alpha <0$  nên $\cos \alpha >0$.\\
			 Do đó $\cos \alpha =\dfrac{2\sqrt{2}}{3} $ suy ra $\tan \alpha =\dfrac{\sin \alpha}{\cos \alpha} =\dfrac{\dfrac{1}{3}}{\dfrac{2\sqrt{2}}{3}} =\dfrac{\sqrt{2}}{4} $ và $\cot \alpha =\dfrac{1}{\tan \alpha} =2\sqrt{2} $.
			\item Ta có $\sin ^{2} \alpha =1-\cos ^{2} \alpha =1-(-0,7)^{2}  =\dfrac{51}{100} $. Vì $\pi <\alpha <\dfrac{3\pi}{2} $ nên $\sin \alpha <0$.\\
			Do đó $\sin \alpha =-\dfrac{\sqrt{51}}{10} $ suy  ra $\tan \alpha =\dfrac{\sin \alpha}{\cos \alpha} =\dfrac{-\dfrac{\sqrt{51}}{10}}{-0,7} =\dfrac{\sqrt{51}}{7} $ và $\cot \alpha =\dfrac{1}{\tan \alpha} =\dfrac{7\sqrt{51}}{51} $.
			\item  Ta có
			\begin{itemize}
				\item $\cot \alpha =\dfrac{1}{\tan \alpha} =\dfrac{1}{2}  $;
				\item $\dfrac{1}{\cos ^{2} \alpha} =1+\tan ^{2}  \alpha =1+2^{2} =5$.
			\end{itemize}
			Suy ra $\cos ^{2} \alpha =\dfrac{1}{5} $. Vì $0<\alpha <\dfrac{\pi}{2} $ nên $\cos \alpha >0$. Do đó $\cos \alpha =\dfrac{\sqrt{5}}{5}  $.\\
			Suy ra $\sin \alpha =\tan \alpha\cdot \cos \alpha =\dfrac{2\sqrt{5}}{5}  $.
			\item Ta có 
			\begin{itemize}
				\item  $\tan \alpha =\dfrac{1}{\cot \alpha} =\dfrac{3}{7} $;
				\item $\dfrac{1}{\sin ^{2} \alpha} =1+\cot ^{2}  \alpha =1+\left(\dfrac{7}{3} \right)^{2} =\dfrac{58}{9} $.
			\end{itemize}
			Suy ra $\sin ^{2} \alpha =\dfrac{9}{58} $. Vì $\pi <\alpha <\dfrac{3\pi}{2} $ nên $\sin \alpha <0$. Do đó  $\sin \alpha =-\dfrac{3\sqrt{58}}{58} $.\\
			Suy ra  $\cos \alpha =\cot \alpha \sin \alpha =\dfrac{7}{3} \cdot \left(-\dfrac{3\sqrt{58}}{58} \right)=-\dfrac{7\sqrt{58}}{58} $.
	\end{enumerate}}
\end{ex}

% nguồn: Ok_TuanLuan_Toan11_Lan2
\begin{ex}%[1D1H2-2]%[ma:53138d70]
	Cho $\sin x=\dfrac{3}{5}\left(\dfrac{\pi}{2} < x < \pi\right)$. Tính giá trị của các giá  trị lượng giác còn lại.
	\loigiai{ 
		Vì $ \dfrac{\pi}{2} <  x < \pi$ nên $\cos x < 0$.\\
		Ta có $\sin ^2x+\cos ^2x=1\Rightarrow \cos ^2x=1-\sin  ^2x=1-\left(\dfrac{3}{5}\right)^2=\dfrac{16}{25}$.\\
		Suy ra $\cos x=-\dfrac{4}{5}$.\\ 
		Khi đó $\tan x=\dfrac{\sin x}{\cos x}=\dfrac{\dfrac{3}{5}}{-\dfrac{4}{5}}=-\dfrac{3}{4}$.\\ 
		Và $\cot x=\dfrac{\cos x}{\sin  x}=\dfrac{-\dfrac{4}{5}}{\dfrac{3}{5}}=-\dfrac{4}{3}$.
	} 
\end{ex}

% nguồn: Ok_TuanLuan_Toan11_Lan2
\begin{ex}%[1D1H2-2]%[ma:4fa312c7]
	Tính giá trị lượng giác còn  lại của góc $x$ biết
	\begin{enumerate}
		\item  $\cos x=\dfrac{2}{\sqrt{5}}$ với $-\dfrac{\pi}{2} < x < 0$.
		\item $\cos x=\dfrac{4}{5}$  với $270^\circ < x < 360^\circ$.
		\item $\sin x=\dfrac{5}{13}$  với $\dfrac{\pi}{2} < x < \pi$
		\item $\sin x=-\dfrac{1}{3}$ với $180^\circ <  x < 270^\circ$.
	\end{enumerate}
	\loigiai{ 
		\begin{enumerate}
			\item Do $-\dfrac{\pi}{2} < x < 0$ $\Rightarrow \heva{&\sin x < 0  \\
			&\cos x > 0 \\ 
			&\tan x  < 0 \\
			&\cot x < 0.}$ 
			Từ  đó với $\cos x=\dfrac{2}{\sqrt{5}}\Rightarrow \sin x=-\sqrt{1-\cos ^2x}=-\dfrac{\sqrt{5}}{5}\Rightarrow
			 \heva{&\tan x=\dfrac{\sin x}{\cos x}=-\dfrac{1}{2}\\
			&\cot x=\dfrac{1}{\tan x}=-2.}$ 
			\item Do $270^\circ < x  < 360^\circ\Rightarrow
			\heva{&\sin x < 0  \\
			&\cos x > 0  \\
			 &\tan x < 0 \\
			&\cot x < 0. 
			 }$
			Từ đó với $\cos x=\dfrac{4}{5}\Rightarrow \sin x=-\sqrt{1-\cos  ^2x}=-\dfrac{3}{5}\Rightarrow
			\heva{&\tan x=\dfrac{\sin x}{\cos  x}=-\dfrac{3}{4} \\
			&\cot  x=\dfrac{1}{\tan x}=-\dfrac{4}{3}.
			}$ 
			\item Do $\dfrac{\pi}{2} <  x < \pi\Rightarrow
			\heva{&\sin x > 0 \\ 
			&\cos x < 0 \\ 
			&\tan  x < 0 \\
			&\cot x <  0.
			 }$
			Từ đó với $\sin  x=\dfrac{5}{13}\Rightarrow \cos x=-\sqrt{1-\sin ^2x}=-\dfrac{12}{13}\Rightarrow
			 \heva{&\tan x=\dfrac{\sin x}{\cos x}=-\dfrac{5}{12} \\
			&\cot x=\dfrac{1}{\tan  x}=-\dfrac{12}{5}.
			}$ 
			\item Do $180^\circ < x < 270^\circ\Rightarrow 
			\heva{&\sin  x < 0 \\
			&\cos x < 0  \\
			&\tan x > 0 \\ 
			&\cot x  > 0.
			}$ 
			Từ đó  với $\sin x=-\dfrac{1}{3}\Rightarrow \cos x=-\sqrt{1-\sin ^2x}=-\dfrac{2\sqrt{2}}{3} \Rightarrow
			\heva{&\tan x=\dfrac{\sin x}{\cos  x}=\dfrac{\sqrt{2}}{4} \\
			&\cot x=\dfrac{1}{\tan x}=2\sqrt{2}. 
			 }$
		\end{enumerate}
	 }
\end{ex}

% nguồn: Ok_TuanLuan_Toan11_Lan2
\begin{ex}%[1D1H2-2]%[ma:cd2dea35]
	Cho $\tan x=\dfrac{3}{4}\left(-\pi < x <-\dfrac{\pi}{2}\right)$. Tính  giá trị của các giá trị lượng giác còn lại.
	\loigiai{ 
		Ta có $\tan  x \cdot \cot x=1 \Rightarrow \cot x=\dfrac{1}{\tan x}=\dfrac{1}{\dfrac{3}{4}}=\dfrac{4}{3}$.\\
		Vì $-\pi < x  <-\dfrac{\pi}{2}$ nên $\cos x < 0$.\\
		Mà $\dfrac{1}{\cos ^2x}=1+\tan ^2x=1+\left(\dfrac{3}{4}\right)^2=\dfrac{25}{16} \Rightarrow \cos  ^2x=\dfrac{16}{25}$.\\
		Suy ra  $\cos x=-\dfrac{4}{5}$.\\
		Và $\tan x=\dfrac{\sin x}{\cos x} \Rightarrow \sin x=\tan x \cdot \cos  x=\dfrac{3}{4} \cdot\left(-\dfrac{4}{5}\right)=-\dfrac{3}{5}$.
	} 
\end{ex}

% nguồn: Ok_TuanLuan_Toan11_Lan2
\begin{ex}%[1D1H2-2]%[ma:50bb5dc0]
	Cho $\cos x=\dfrac{2}{\sqrt{5}}\left(-\dfrac{\pi}{2} < x < 0\right)$. Tính giá trị của các giá trị lượng giác còn  lại.
	\loigiai{ 
		Vì $-\dfrac{\pi}{2} < x < 0$  nên $\sin x < 0$.\\
		Ta  có $\sin ^2x+\cos ^2x=1\Rightarrow \sin ^2x=1-\cos ^2x=1-\left(\dfrac{2}{\sqrt{5}}\right)^2=\dfrac{1}{5}$.\\
		Suy ra  $\sin x=-\dfrac{1}{\sqrt{5}}$.\\
		 Khi đó
		$\tan x=\dfrac{\sin x}{\cos x}=\dfrac{-\dfrac{1}{\sqrt{5}}}{\dfrac{2}{\sqrt{5}}}=-\dfrac{1}{2}.$\\ 
		Và $\cot  x=\dfrac{\cos x}{\sin x}=\dfrac{\dfrac{2}{\sqrt{5}}}{-\dfrac{1}{\sqrt{5}}}=-2.$
	 }
\end{ex}

% nguồn: Ok_TuanLuan_Toan11_Lan2
\begin{ex}%[1D1V2-2]%[ma:6a85ba8a]
	Tính giá  trị lượng giác của các biểu thức sau
	\begin{enumerate}
		\item Cho $\tan x=-2$. Tính $A_1=\dfrac{5\cot x+4\tan x}{5\cot x-4\tan x}$, $A_2=\dfrac{2\sin  x+\cos x}{\cos x-3\sin x}$.
		\item Cho $\cot x=\sqrt{2}$. Tính  $B_1=\dfrac{3\sin x-\cos x}{\sin x+\cos x}$, $B_2=\dfrac{\sin x-3\cos x}{\sin x+3\cos x}$.
		\item Cho $\cot x=2$. Tính $C_1=\dfrac{2\sin x+3\cos x}{3\sin x-2\cos x}$, $C_2=\dfrac{2}{\cos ^2x-\sin x\cos  x}$.
		\item Cho $\sin x=\dfrac{3}{5}$, $0< x < \dfrac{\pi}{2}$. Tính $E=\dfrac{\cot x+\tan x}{\cot  x-\tan x}$.
		\item Cho $\sin x=\dfrac{1}{5}$, $90^\circ< x < 180^\circ$. Tính  $F=\dfrac{8\tan ^2x+3\cot x-1}{\tan x+\cot x}$.
	\end{enumerate}
	\loigiai{ 
		\begin{enumerate}
			\item $\tan x=-2\Rightarrow \cot x=-\dfrac{1}{2} \Rightarrow  A_1=\dfrac{5\cot x+4\tan x}{5\cot x-4\tan x}=\dfrac{-\dfrac{5}{2}+4.\left(-2\right)}{-\dfrac{5}{2}-4\cdot\left(-2\right)}=-\dfrac{21}{11}$.
			$\tan x=-2\Rightarrow A_2=\dfrac{2\sin x+\cos x}{\cos x-3\sin x}=\dfrac{2\tan x+1}{1-3\tan x}=\dfrac{2\cdot (-2)+1}{1-3\cdot (-2)}=-\dfrac{3}{7} 
			 $.
			 \item $\cot x=\sqrt{2} \Rightarrow B_1=\dfrac{3\sin x-\cos x}{\sin x+\cos x}=\dfrac{3-\cot x}{1+\cot x}=\dfrac{3-\sqrt{2}}{1+\sqrt{2}}=-5+4\sqrt{2}$.
			$\cot x=\sqrt{2} \Rightarrow B_2=\dfrac{\sin x-3\cos x}{\sin  x+3\cos x}=\dfrac{1-3\cot x}{1+3\cot x}=\dfrac{1-3\sqrt{2}}{1+3\sqrt{2}}=\dfrac{-19+6\sqrt{2}}{17}
			 $.
			\item $\cot x=2\Rightarrow C_1=\dfrac{2\sin x+3\cos  x}{3\sin x-2\cos x}=\dfrac{2+3\cot x}{3-2\cot x}=\dfrac{2+3\cdot 2}{3-2\cdot 2}=-8$.
			$\cot x=2\Rightarrow C_2=\dfrac{2}{\cos ^2x-\sin x\cos x}=\dfrac{2\left(1+\cot ^2x\right)}{3\cot ^2x-\cot x}=\dfrac{2\left(1+2^2\right)}{3\cdot 2^2-2}=1 
			 $.
			\item $0< x < \dfrac{\pi}{2} \Rightarrow \heva{&{\sin x >  0} \\
			&{\cos x >  0}
			}$ $\Rightarrow \cos x=\sqrt{1-\left(\dfrac{3}{5} \right)^2}=\dfrac{4}{5}\Rightarrow \tan x=\dfrac{\sin x}{\cos x}=\dfrac{3}{4}$; $\cot x=\dfrac{4}{3}$. 
			$\to  E=\dfrac{\cot x+\tan x}{\cot x-\tan x}=\dfrac{\dfrac{4}{3}+\dfrac{3}{4}}{\dfrac{4}{3}-\dfrac{3}{4}}=\dfrac{25}{7}
			$. 
			\item Ta có $90^{\circ} < x <  180^{\circ} \Rightarrow \heva{&{\sin x > 0} \\
			&{\cos x <  0}
			}$ $\Rightarrow \cos x=-\sqrt{1-\left(\dfrac{1}{3} \right)^2}=-\dfrac{2\sqrt{2}}{3}$. 
			$\Rightarrow \tan x=\dfrac{\sin x}{\cos  x}=\dfrac{-1}{2\sqrt{2}}$; $\cot x=-2\sqrt{2}
			$. 
			Do đó $F=\dfrac{8\tan ^2x+3\cot  x-1}{\tan x+\cot x}=\dfrac{8\cdot\left(\dfrac{-1}{2\sqrt{2}} \right)^2-3\cdot 2\sqrt{2}-1}{-\dfrac{1}{2\sqrt{2}}-2\sqrt{2}}=\dfrac{8}{3}$.
		\end{enumerate}
	 }
\end{ex}

% nguồn: Ok_TuanLuan_Toan11_Lan2
\begin{ex}%[1D1V2-2]%[ma:5718b834]
	Cho  $\dfrac{\sin ^4\alpha}{a}+\dfrac{\cos ^4\alpha}{b}=\dfrac{1}{a+b}$. Tính giá trị của biểu thức: $A=\dfrac{\sin ^8\alpha}{a^3}+\dfrac{\cos ^8\alpha}{b^3}$.
	 \loigiai{
		Đặt $\cos ^2\alpha=t $, theo  đề, ta có
		\begin{eqnarray*}
			&&  \dfrac{\left(1-t\right)^2}{a}+\dfrac{t^2}{b}=\dfrac{1}{a+b}\\
			\Leftrightarrow  && b\left(1-t\right)^2+at^2=\dfrac{ab}{a+b} \\
			\Leftrightarrow && at^2+bt^2-2bt+b=\dfrac{ab}{a+b} \\ 
			\Leftrightarrow && \left(a+b\right)t^2-2bt+b=\dfrac{ab}{a+b} \\ 
			\Leftrightarrow  && \left(a+b\right)^2t^2-2b\left(a+b\right)t+b^2=0\\
			\Leftrightarrow && t=\dfrac{b}{a+b}. 
		\end{eqnarray*}
		Suy ra  $\cos ^2\alpha=\dfrac{b}{a+b}$; $\sin ^2\alpha=\dfrac{a}{a+b}$.\\
		Vậy $\dfrac{\sin ^8\alpha}{a^3}+\dfrac{\cos  ^8\alpha}{b^3}=\dfrac{a}{\left(a+b\right)^4}+\dfrac{b}{\left(a+b\right)^4}=\dfrac{1}{\left(a+b\right)^3}$.
	 }
\end{ex}

% nguồn: Ok_TuanLuan_Toan11_Lan2
\begin{ex}%[1D1V2-2]%[ma:a30c27cd]
	Tìm giá trị nhỏ nhất của biểu thức $P=\cot ^4a+\cot ^4b+2\tan  ^2a\cdot\tan ^2b+2$.
	 \loigiai{
		\begin{align*}
			P&=\left(\cot ^2a-\cot ^2b\right)^2+2\cot ^2a\cdot \cot ^2b+2\tan ^2a\cdot \tan  ^2b+2 \\
			&=\left(\cot ^2a-\cot ^2b\right)^2+2\left(\cot ^2a\cdot \cot ^2b+\tan ^2a\cdot  \tan ^2b-2\right)+6 \\
			&=\left(\cot ^2a-\cot ^2b\right)^2+2\left(\cot  ^2a\cdot \cot ^2b+\tan ^2a\cdot \tan ^2b-2\cot a\cdot \cot b\cdot \tan a\cdot \tan b\right)+6\\
			&=\left(\cot ^2a-\cot  ^2b\right)^2+2\left(\cot a\cdot \cot b-\tan a\cdot \tan b\right)^2+6\ge 6.
		\end{align*}
		Dấu  "$=$" xảy ra khi
		$\heva{&\cot  ^2a=\cot ^2b\\
		 &\cot a\cdot \cot b=\tan a\cdot \tan b}
		\Leftrightarrow  \heva{&\cot ^2a=1 \\
		&\cot ^2b=1} 
		\Leftrightarrow a=b=\dfrac{\pi}{4}+\dfrac{k\pi}{2}, (k\in {\bf \mathbb{Z}}). 
		$ 
	 }
\end{ex}
